% ---------- Page and fonts ----------
\documentclass[11pt,a4paper,openany,oneside]{book}
\usepackage[T1]{fontenc}
\usepackage[utf8]{inputenc}
\usepackage[sc]{mathpazo}
\linespread{1.08}
\usepackage[scaled=0.9]{helvet}
\usepackage{amsmath,amssymb,mathtools}
\usepackage{bm}
\usepackage{cancel}
\usepackage[a4paper,top=24mm,bottom=24mm,inner=24mm,outer=24mm,headheight=15pt,headsep=8mm,footskip=12mm]{geometry}
\usepackage{microtype}
\usepackage[dvipsnames,svgnames]{xcolor}
\usepackage{graphicx}
\usepackage{booktabs,array,tabularx}
\usepackage{multicol}
\usepackage{enumitem}
\setlist{nosep,leftmargin=1.4em,topsep=3pt}
\usepackage{tikz}
\usetikzlibrary{angles,quotes,arrows.meta,calc,decorations.pathreplacing,patterns,positioning,shapes.geometric,decorations.markings,3d,intersections}
\usepackage{pgfplots}
\pgfplotsset{compat=1.17}
\usepgfplotslibrary{fillbetween,groupplots}
\usepackage{tikz-3dplot}
\usepackage[most]{tcolorbox}
\usepackage{titlesec}
\usepackage{fancyhdr}
\usepackage{etoolbox}
\usepackage[hidelinks,unicode,psdextra,bookmarksnumbered=true,bookmarksopen=true,bookmarksopenlevel=1,pdfpagemode=UseOutlines]{hyperref}
\usepackage{bookmark}
\pdfstringdefDisableCommands{%
  \def\vect#1{#1}%
  \def\uv#1{#1}%
  \def\mathbf#1{#1}%
  \def\boldsymbol#1{#1}%
  \def\why#1{}%
  \def\bk#1{ [Book #1]}%
}
\hypersetup{
  pdftitle={Linear Algebra: The Complete Equation and Derivation Book (ASD-STE100 Edition)},
  pdfauthor={Based on Gilbert Strang, Introduction to Linear Algebra, 4th Edition},
  pdfsubject={Linear Algebra Derivations from Plus-Two Base, In-Text Examples, and Unique Worked Problems}
}

% ---------- Overflow Prevention ----------
\emergencystretch=3em
\allowdisplaybreaks

% ---------- Colours ----------
\definecolor{ink}{HTML}{1F2A33}
\definecolor{teal}{HTML}{0F6E6E}
\definecolor{tealsoft}{HTML}{E8F3F2}
\definecolor{amber}{HTML}{B4690E}
\definecolor{ambersoft}{HTML}{FFF6E8}
\definecolor{plum}{HTML}{6A3D7A}
\definecolor{plumsoft}{HTML}{F5EFF7}
\definecolor{rose}{HTML}{B03A48}
\definecolor{rosesoft}{HTML}{FCEEF0}
\definecolor{slate}{HTML}{5B6770}
\definecolor{slatesoft}{HTML}{F3F5F6}
\definecolor{bluec}{HTML}{2459A8}
\color{ink}

% ---------- Headings ----------
\newcommand{\sffamilybold}{\fontfamily{phv}\selectfont\bfseries}
\titleformat{\part}[display]{\centering\sffamilybold\color{teal}}{\Large\MakeUppercase{\thepart}}{8pt}{\Huge}
\renewcommand{\thepart}{Module \arabic{part}}
\usepackage{titletoc}
\titlecontents{part}[0pt]{\addvspace{12pt}\fontfamily{phv}\selectfont\bfseries\color{teal}}{\thecontentslabel\quad}{}{\hfill\contentspage}[\addvspace{3pt}]
\titleformat{\chapter}[display]
  {\sffamilybold\color{ink}}
  {\color{teal}\large CHAPTER \thechapter}{4pt}{\huge}[\vspace{2pt}{\color{teal}\titlerule[1.2pt]}]
\titlespacing*{\chapter}{0pt}{0pt}{18pt}
\titleformat{\section}{\sffamilybold\Large\color{teal}}{\thesection}{0.7em}{}
\titleformat{\subsection}{\sffamilybold\large\color{ink}}{\thesubsection}{0.7em}{}
\titleformat{\subsubsection}{\sffamilybold\normalsize\color{slate}}{}{0em}{}
\titlespacing*{\section}{0pt}{16pt}{7pt}
\titlespacing*{\subsection}{0pt}{13pt}{5pt}
\titlespacing*{\subsubsection}{0pt}{9pt}{3pt}
\setcounter{secnumdepth}{2}
\setcounter{tocdepth}{2}

% ---------- Header / footer ----------
\pagestyle{fancy}
\fancyhf{}
\renewcommand{\headrulewidth}{0.3pt}
\renewcommand{\chaptermark}[1]{\markboth{#1}{}}
\renewcommand{\sectionmark}[1]{\markright{\thesection\ \ #1}}
\fancyhead[L]{\small\fontfamily{phv}\selectfont\color{slate}\leftmark}
\fancyhead[R]{\small\fontfamily{phv}\selectfont\color{slate}\rightmark}
\fancyfoot[C]{\small\fontfamily{phv}\selectfont\color{slate}\thepage}
\fancypagestyle{plain}{\fancyhf{}\renewcommand{\headrulewidth}{0pt}\fancyfoot[C]{\small\fontfamily{phv}\selectfont\color{slate}\thepage}}

\setlength{\parindent}{0pt}
\setlength{\parskip}{5pt}

% ---------- Boxes (Fixed top padding & breakable to eliminate overflow) ----------
\tcbset{common/.style={enhanced,breakable,boxrule=0pt,arc=2mm,left=3mm,right=3mm,top=4.5mm,bottom=2.5mm,
  fonttitle=\fontfamily{phv}\selectfont\bfseries\small,coltitle=white,before skip=8pt,after skip=8pt}}
% Goal of a derivation
\newtcolorbox{goal}{common,colback=tealsoft,colframe=teal,borderline west={3pt}{0pt}{teal},
  title={GOAL},attach boxed title to top left={xshift=3mm,yshift=-2.5mm},boxed title style={colback=teal,arc=1mm}}
% Plus-two base facts
\newtcolorbox{base}[1][]{common,colback=ambersoft,colframe=amber,borderline west={3pt}{0pt}{amber},
  title={FROM YOUR +2 BASE#1},attach boxed title to top left={xshift=3mm,yshift=-2.5mm},boxed title style={colback=amber,arc=1mm}}
% Final result
\newtcolorbox{result}[1][]{common,breakable=true,colback=white,colframe=teal,boxrule=1.2pt,
  title={RESULT#1},attach boxed title to top left={xshift=3mm,yshift=-2.5mm},boxed title style={colback=teal,arc=1mm}}
% Meaning
\newtcolorbox{means}{common,colback=plumsoft,colframe=plum,borderline west={3pt}{0pt}{plum},
  title={WHAT IT MEANS},attach boxed title to top left={xshift=3mm,yshift=-2.5mm},boxed title style={colback=plum,arc=1mm}}
% Exam trap
\newtcolorbox{trap}{common,colback=rosesoft,colframe=rose,borderline west={3pt}{0pt}{rose},
  title={WATCH OUT},attach boxed title to top left={xshift=3mm,yshift=-2.5mm},boxed title style={colback=rose,arc=1mm}}
% Worked example (Question box)
\newtcolorbox{example}[2]{common,breakable=true,colback=slatesoft,colframe=slate,borderline west={3pt}{0pt}{slate},
  before upper={\belowpdfbookmark{Book Example #1}{ex#1}},
  title={BOOK EXAMPLE\ \ #1\quad #2},attach boxed title to top left={xshift=3mm,yshift=-2.5mm},boxed title style={colback=slate,arc=1mm}}
% End-of-page Solution box
\newtcolorbox{ pageend_sol }[1]{common,breakable=true,colback=tealsoft!45,colframe=teal!80!black,borderline west={3pt}{0pt}{teal!80!black},
  title={END-OF-PAGE SOLUTION FOR EXAMPLE\ \ #1},attach boxed title to top left={xshift=3mm,yshift=-2.5mm},boxed title style={colback=teal!80!black,arc=1mm}}
% Idea / picture first
\newtcolorbox{idea}{common,colback=white,colframe=teal!50,boxrule=0.6pt,borderline west={3pt}{0pt}{teal!50},
  title={THE IDEA},coltitle=teal,attach boxed title to top left={xshift=3mm,yshift=-2.5mm},boxed title style={colback=white,colframe=teal!50,boxrule=0.6pt,arc=1mm}}
% Syllabus strip
\newtcolorbox{syll}{enhanced,breakable,colback=slatesoft,colframe=slatesoft,arc=1mm,left=3mm,right=3mm,top=2mm,bottom=2mm,fontupper=\small\fontfamily{phv}\selectfont}

% ---------- Derivation helpers ----------
\newcommand{\why}[1]{\text{\footnotesize\color{slate}\itshape (#1)}}
\newcommand{\step}[1]{\par\smallskip{\fontfamily{phv}\selectfont\bfseries\small\color{teal}#1}\enspace}
\newcommand{\term}[1]{\textbf{\color{teal}#1}}
\newcommand{\eqbox}[1]{\boxed{\,#1\,}}
\newcommand{\vect}[1]{\mathbf{#1}}
\newcommand{\uv}[1]{\hat{\mathbf{#1}}}
\newcommand{\ihat}{\boldsymbol{i}}
\newcommand{\jhat}{\boldsymbol{j}}
\newcommand{\khat}{\boldsymbol{k}}
\newcommand{\dd}{\,\mathrm{d}}
\newcommand{\diff}{\mathrm{d}}
\newcommand{\T}{\mathrm{T}}
\newcommand{\Herm}{\mathrm{H}}
\newcommand{\bookref}[1]{{\small\color{slate}[Book: #1]}}
\newcounter{figno}[chapter]
\newcommand{\figcap}[1]{\par\nopagebreak[4]\smallskip\refstepcounter{figno}{\centering\small\color{slate}\fontfamily{phv}\selectfont \textbf{Fig.\ \thechapter.\arabic{figno}}\quad #1\par}\medskip}

% figure environment wrapper, no float (keeps math next to picture)
\newenvironment{pic}{\par\medskip\noindent\begin{minipage}{\linewidth}\centering}{\end{minipage}\par\medskip}
\usepackage{needspace}
\pretocmd{\section}{\needspace{7\baselineskip}}{}{}
\pretocmd{\subsection}{\needspace{5\baselineskip}}{}{}
\pretocmd{\subsubsection}{\needspace{4\baselineskip}}{}{}

% equation numbering per chapter
\numberwithin{equation}{chapter}

% TikZ styles (SVG-equivalent vector graphics)
\tikzset{
  atom/.style={circle,fill=teal,inner sep=0pt,minimum size=5pt},
  atomB/.style={circle,fill=amber,inner sep=0pt,minimum size=5pt},
  latt/.style={circle,fill=ink,inner sep=0pt,minimum size=3pt},
  vec/.style={-{Stealth[length=2.6mm]},thick},
  ax/.style={-{Stealth[length=2mm]},thin,slate},
  dim/.style={{Stealth[length=1.6mm]}-{Stealth[length=1.6mm]},thin,slate},
  lbl/.style={font=\small},
}

% ---------- Toolkit items (T1, T2, ...) ----------
\newcounter{tool}
\newtcolorbox{toolbox}[2]{common,breakable=true,colback=white,colframe=amber!70,boxrule=0.6pt,
  borderline west={3pt}{0pt}{amber},
  title={\refstepcounter{tool}\label{#2}T\thetool\quad #1},coltitle=amber!80!black,
  before upper={\belowpdfbookmark{T\thetool\ \ #1}{tool\thetool}},
  attach boxed title to top left={xshift=3mm,yshift=-2.5mm},
  boxed title style={colback=white,colframe=amber!70,boxrule=0.6pt,arc=1mm}}
\newcommand{\tref}[1]{\hyperref[#1]{\textbf{\color{amber!80!black}T\ref*{#1}}}}

% ---------- 3D oblique projection ----------
\tikzset{
  cub/.style={x={(1cm,0cm)},y={(0.42cm,0.30cm)},z={(0cm,1cm)}},
}
\raggedbottom

% ---------- Formula sheets ----------
\newcommand{\formulasheet}[1]{%
  \cleardoublepage\phantomsection
  \chapter*{#1}\addcontentsline{toc}{chapter}{#1}\markboth{#1}{}%
}
\newtcolorbox{fs}[1]{enhanced,breakable=true,colback=white,colframe=teal!60,boxrule=0.5pt,arc=1.5mm,
  left=2.5mm,right=2.5mm,top=2mm,bottom=2mm,before skip=5pt,after skip=5pt,
  title={#1},fonttitle=\fontfamily{phv}\selectfont\bfseries\footnotesize,colbacktitle=tealsoft,coltitle=teal}

\newcommand{\bk}[1]{\texorpdfstring{\enspace{\normalfont\small\color{slate}[Book #1]}}{ [Book #1]}}

\begin{document}

% ============================================================================
% TITLE PAGE
% ============================================================================
\begin{titlepage}
\centering
\vspace*{15mm}
{\fontfamily{phv}\selectfont\bfseries\large\color{teal} ASD-STE100 TECHNICAL MATHEMATICS REFERENCE\par}
\vspace{6mm}
{\fontfamily{phv}\selectfont\bfseries\Huge\color{ink} LINEAR ALGEBRA:\par}
\vspace{3mm}
{\fontfamily{phv}\selectfont\bfseries\LARGE\color{teal} Equations, First-Principle Derivations,\par}
{\fontfamily{phv}\selectfont\bfseries\LARGE\color{teal} and Complete Solved Problems\par}
\vspace{8mm}
{\color{teal}\rule{\linewidth}{1.5pt}\par}
\vspace{6mm}
{\large\color{slate} Written in Simplified Technical English (ASD-STE100).\\
All syllabus equations derived step-by-step from Plus-Two (+2) base knowledge.\\
Includes vector graphics diagrams, verbatim in-text examples with end-of-page solutions,\\
and all unique Worked Examples at the end of each module.\par}
\vspace{10mm}

\begin{tcolorbox}[enhanced,colback=slatesoft,colframe=teal,boxrule=1pt,arc=2mm,width=0.92\linewidth,left=5mm,right=5mm,top=4mm,bottom=4mm]
\fontfamily{phv}\selectfont\small
\textbf{\color{teal}SYLLABUS COVERAGE (GILBERT STRANG, 4TH EDITION):}\\[4pt]
\textbf{Module 1:} Vectors and Linear Combinations (1.1) $\cdot$ Spaces of Vectors (3.1) $\cdot$ Nullspace of $A$ (3.2) $\cdot$ Rank and Row Reduced Form $R$ (3.3) $\cdot$ Complete Solution to $A\vect{x}=\vect{b}$ (3.4) $\cdot$ Independence, Basis, Dimension (3.5) $\cdot$ Dimensions of the Four Subspaces (3.6).\\[4pt]
\textbf{Module 2:} Vectors and Linear Equations (2.1) $\cdot$ Idea of Elimination (2.2) $\cdot$ Elimination Using Matrices (2.3) $\cdot$ Rules for Matrix Operations (2.4) $\cdot$ Inverse Matrices (2.5) $\cdot$ Factorization $A=LU$ (2.6) $\cdot$ Transposes and Permutations (2.7) $\cdot$ Properties of Determinants (5.1) $\cdot$ Permutations and Cofactors (5.2) $\cdot$ Cramer's Rule, Inverses, Volumes (5.3).\\[4pt]
\textbf{Module 3:} Lengths and Dot Products (1.2) $\cdot$ Orthogonality of the Four Subspaces (4.1) $\cdot$ Projections (4.2) $\cdot$ Least Squares Approximations (4.3) $\cdot$ Orthogonal Bases and Gram-Schmidt (4.4) $\cdot$ Positive Definite Matrices (6.5) $\cdot$ Linear Algebra for Statistics and Probability (8.6).\\[4pt]
\textbf{Module 4:} Introduction to Eigenvalues (6.1) $\cdot$ Diagonalizing a Matrix (6.2) $\cdot$ Symmetric Matrices (6.4) $\cdot$ Similar Matrices and Jordan Form (6.6) $\cdot$ Singular Value Decomposition (6.7) $\cdot$ Markov Matrices and Economics (8.3) $\cdot$ Complex Numbers (10.1) $\cdot$ Hermitian and Unitary Matrices (10.2--10.3).
\end{tcolorbox}
\vfill
{\fontfamily{phv}\selectfont\small\color{slate} Compiled for Direct PDF Rendering $\cdot$ Overflow-Protected Layout\par}
\end{titlepage}

\tableofcontents
\cleardoublepage

% ============================================================================
% CHAPTER 0: PLUS-TWO (+2) BASE TOOLKIT
% ============================================================================
\chapter[Plus-Two (+2) Mathematical Toolkit]{Plus-Two (+2) Mathematical Toolkit}

\begin{syll}
Use these twelve base facts from Plus-Two (+2) mathematics to derive all equations in Modules 1 to 4.
\end{syll}

\begin{toolbox}{Ordered Pairs, Triples, and Pythagoras' Theorem}{tool:pyth}
In the 2D plane, a point has coordinates $(v_1, v_2)$. A right triangle with horizontal leg $v_1$ and vertical leg $v_2$ has hypotenuse $c$:
\begin{equation}
c^2 = v_1^2 + v_2^2 \implies c = \sqrt{v_1^2 + v_2^2}.
\end{equation}
In 3D space, apply Pythagoras' theorem twice to a box with edges $v_1, v_2, v_3$:
\begin{equation}
c^2 = (v_1^2 + v_2^2) + v_3^2 = v_1^2 + v_2^2 + v_3^2.
\end{equation}
\end{toolbox}

\begin{toolbox}{Trigonometric Identities and Cosine Addition Formula}{tool:trig}
For any angle $\theta$ on the unit circle (radius $r=1$), the coordinates are $(\cos\theta, \sin\theta)$:
\begin{equation}
\cos^2\theta + \sin^2\theta = 1, \qquad -1 \le \cos\theta \le 1.
\end{equation}
For two angles $\alpha$ and $\beta$, the cosine of the difference $\theta = \beta - \alpha$ is:
\begin{equation}
\cos(\beta - \alpha) = \cos\beta\cos\alpha + \sin\beta\sin\alpha.
\end{equation}
\end{toolbox}

\begin{toolbox}{Algebraic Squares and Quadratic Formula}{tool:quad}
For any real numbers $a$ and $b$:
\begin{equation}
(a \pm b)^2 = a^2 \pm 2ab + b^2 \ge 0 \implies 2ab \le a^2 + b^2.
\end{equation}
The roots of the quadratic equation $a\lambda^2 + b\lambda + c = 0$ ($a\neq 0$) are:
\begin{equation}
\lambda = \frac{-b \pm \sqrt{b^2 - 4ac}}{2a}.
\end{equation}
The sum of the roots is $\lambda_1 + \lambda_2 = -b/a$. The product of the roots is $\lambda_1\lambda_2 = c/a$.
\end{toolbox}

\begin{toolbox}{Simultaneous Linear Equations in Two Unknowns}{tool:sim2}
Consider two linear equations: $ax + by = f$ and $cx + dy = g$.
Multiply equation 1 by $c/a$ (when $a\neq 0$) and subtract it from equation 2:
\begin{equation}
\left(d - \frac{c}{a}b\right)y = g - \frac{c}{a}f \implies (ad - bc)y = ag - cf.
\end{equation}
A unique solution exists if and only if $ad - bc \neq 0$.
\end{toolbox}

\begin{toolbox}{Sigma Summation and Geometric Series}{tool:sigma}
The sum of $n$ products $a_1b_1 + a_2b_2 + \dots + a_nb_n$ is written as $\sum_{i=1}^n a_i b_i$.
The sum of the first $n$ integers and first $n$ squares are:
\begin{equation}
\sum_{k=1}^n k = \frac{n(n+1)}{2}, \qquad \sum_{k=1}^n k^2 = \frac{n(n+\frac{1}{2})(n+1)}{3} \approx \frac{1}{3}n^3.
\end{equation}
For $|x| < 1$, the infinite geometric series converges:
\begin{equation}
1 + x + x^2 + x^3 + \dots = \frac{1}{1-x}.
\end{equation}
\end{toolbox}

\begin{toolbox}{Calculus: Derivatives for Minimum and Taylor Series}{tool:calc}
A differentiable function $E(x)$ reaches its minimum where its derivative is zero:
\begin{equation}
\frac{\diff E}{\diff x} = 0, \qquad \frac{\diff^2 E}{\diff x^2} > 0.
\end{equation}
For a function of two variables $E(C,D)$, both partial derivatives must be zero at a minimum: $\partial E/\partial C = 0$ and $\partial E/\partial D = 0$.
The power series for $e^x$, $\cos\theta$, and $\sin\theta$ are:
\begin{equation}
e^x = 1 + x + \frac{x^2}{2!} + \frac{x^3}{3!} + \dots, \quad \cos\theta = 1 - \frac{\theta^2}{2!} + \frac{\theta^4}{4!} - \dots, \quad \sin\theta = \theta - \frac{\theta^3}{3!} + \frac{\theta^5}{5!} - \dots
\end{equation}
\end{toolbox}

% ============================================================================
% MODULE 1: VECTORS, VECTOR SPACES, AND SUBSPACES (PDF 2)
% ============================================================================
\part{Vectors, Vector Spaces, and Subspaces}

\chapter[Vectors and Linear Combinations (1.1)]{Vectors and Linear Combinations\bk{1.1}}

\begin{syll}
Syllabus topics: Column vectors $\cdot$ Vector addition $\cdot$ Scalar multiplication $\cdot$ Linear combinations $c\vect{v}+d\vect{w}$ $\cdot$ Vectors in $\mathbf{R}^2$, $\mathbf{R}^3$, and $\mathbf{R}^n$ $\cdot$ Lines, planes, and spaces from linear combinations.
\end{syll}

\section{Vector Addition and Scalar Multiplication\bk{1.1}}

\begin{goal}
Define column vectors, vector addition $\vect{v}+\vect{w}$, and scalar multiplication $c\vect{v}$ from Plus-Two coordinate pairs.
\end{goal}

\begin{base}
In Plus-Two coordinate geometry, a point in the plane has an $x$-coordinate $v_1$ and a $y$-coordinate $v_2$. You add horizontal steps to horizontal steps, and vertical steps to vertical steps.
\end{base}

\step{Step 1: Write a pair of numbers as a column vector.}
Put the two separate numbers $v_1$ and $v_2$ into a single column vector $\vect{v}$:
\begin{equation}
\vect{v} = \begin{bmatrix} v_1 \\ v_2 \end{bmatrix}, \qquad v_1 = \text{first component}, \quad v_2 = \text{second component}.
\end{equation}
To save space in text, we also write $\vect{v} = (v_1, v_2)$ with parentheses and commas. This is still a column vector.

\step{Step 2: Add vectors component by component.}
Add the first components together and add the second components together:
\begin{align}
\vect{v} + \vect{w} &= \begin{bmatrix} v_1 \\ v_2 \end{bmatrix} + \begin{bmatrix} w_1 \\ w_2 \end{bmatrix} = \begin{bmatrix} v_1 + w_1 \\ v_2 + w_2 \end{bmatrix} &&\why{vector addition} \\
\vect{v} - \vect{w} &= \begin{bmatrix} v_1 \\ v_2 \end{bmatrix} - \begin{bmatrix} w_1 \\ w_2 \end{bmatrix} = \begin{bmatrix} v_1 - w_1 \\ v_2 - w_2 \end{bmatrix} &&\why{vector subtraction}
\end{align}

\step{Step 3: Multiply a vector by a real number (scalar) $c$.}
Multiply every component of $\vect{v}$ by the scalar $c$:
\begin{equation}
c\vect{v} = \begin{bmatrix} cv_1 \\ cv_2 \end{bmatrix}, \qquad -\vect{v} = \begin{bmatrix} -v_1 \\ -v_2 \end{bmatrix}, \qquad \vect{0} = \begin{bmatrix} 0 \\ 0 \end{bmatrix}.
\end{equation}

\section{Linear Combinations\bk{1.1}}

\begin{goal}
Combine scalar multiplication and vector addition into one operation: the linear combination.
\end{goal}

\step{Step 1: Multiply $\vect{v}$ by $c$, multiply $\vect{w}$ by $d$, and add.}
\begin{equation}
c\vect{v} + d\vect{w} = c\begin{bmatrix} v_1 \\ v_2 \end{bmatrix} + d\begin{bmatrix} w_1 \\ w_2 \end{bmatrix} = \begin{bmatrix} cv_1 + dw_1 \\ cv_2 + dw_2 \end{bmatrix}.
\end{equation}
For three vectors $\vect{u}, \vect{v}, \vect{w}$ in three-dimensional space $\mathbf{R}^3$, a linear combination is:
\begin{equation}
c\vect{u} + d\vect{v} + e\vect{w} = c\begin{bmatrix} u_1 \\ u_2 \\ u_3 \end{bmatrix} + d\begin{bmatrix} v_1 \\ v_2 \\ v_3 \end{bmatrix} + e\begin{bmatrix} w_1 \\ w_2 \\ w_3 \end{bmatrix}.
\end{equation}

\begin{result}[ (Linear Combinations and Geometry in $\mathbf{R}^3$)]
For typical nonzero vectors $\vect{u}, \vect{v}, \vect{w}$ in $\mathbf{R}^3$:
\begin{enumerate}
\item All combinations $c\vect{u}$ fill a \term{line} through $(0,0,0)$.
\item All combinations $c\vect{u} + d\vect{v}$ fill a \term{plane} through $(0,0,0)$.
\item All combinations $c\vect{u} + d\vect{v} + e\vect{w}$ fill \term{three-dimensional space} $\mathbf{R}^3$.
\end{enumerate}
\end{result}

\begin{pic}
\begin{tikzpicture}[scale=0.72]
  % Left: Parallelogram v + w
  \begin{scope}[xshift=0cm]
    \draw[ax] (-1.5,0) -- (4.5,0) node[right] {$x$};
    \draw[ax] (0,-0.5) -- (0,4.5) node[above] {$y$};
    \draw[vec,teal] (0,0) -- (4,2) node[right] {$\vect{v}=\begin{bmatrix}4\\2\end{bmatrix}$};
    \draw[vec,amber] (0,0) -- (-1,2) node[left] {$\vect{w}=\begin{bmatrix}-1\\2\end{bmatrix}$};
    \draw[dashed,slate] (4,2) -- (3,4);
    \draw[dashed,slate] (-1,2) -- (3,4);
    \draw[vec,plum, very thick] (0,0) -- (3,4) node[above] {$\vect{v}+\vect{w}=\begin{bmatrix}3\\4\end{bmatrix}$};
  \end{scope}
  % Right: Difference v - w
  \begin{scope}[xshift=8.5cm]
    \draw[ax] (-1.5,0) -- (5.5,0) node[right] {$x$};
    \draw[ax] (0,-2.3) -- (0,2.7) node[above] {$y$};
    \draw[vec,teal] (0,0) -- (4,2) node[above] {$\vect{v}=\begin{bmatrix}4\\2\end{bmatrix}$};
    \draw[vec,amber] (0,0) -- (-1,2) node[left] {$\vect{w}$};
    \draw[vec,amber,dashed] (0,0) -- (1,-2) node[below] {$-\vect{w}=\begin{bmatrix}1\\-2\end{bmatrix}$};
    \draw[vec,rose,very thick] (0,0) -- (5,0) node[above right] {$\vect{v}-\vect{w}=\begin{bmatrix}5\\0\end{bmatrix}$};
    \draw[dashed,slate] (4,2) -- (5,0);
  \end{scope}
\end{tikzpicture}
\figcap{Vector addition $\vect{v}+\vect{w}=(3,4)$ along the parallelogram diagonal (left) and subtraction $\vect{v}-\vect{w}=(5,0)$ (right) \bookref{Fig.\ 1.1}.}
\end{pic}

\begin{example}{1.1 (p.\ 2--5)}{Vector Addition and Linear Combinations}
\begin{enumerate}
\item Check $\vect{v}+\vect{w} = \vect{w}+\vect{v}$ by an example:
\[
\vect{v}+\vect{w} = \begin{bmatrix} 1 \\ 5 \end{bmatrix} + \begin{bmatrix} 3 \\ 3 \end{bmatrix}, \qquad \vect{w}+\vect{v} = \begin{bmatrix} 3 \\ 3 \end{bmatrix} + \begin{bmatrix} 1 \\ 5 \end{bmatrix}.
\]
\item Compute the linear combination $\vect{u} + 4\vect{v} - 2\vect{w}$ for:
\[
\begin{bmatrix} 1 \\ 0 \\ 3 \end{bmatrix} + 4\begin{bmatrix} 1 \\ 2 \\ 1 \end{bmatrix} - 2\begin{bmatrix} 2 \\ 3 \\ -1 \end{bmatrix}.
\]
\end{enumerate}
\end{example}

\vfill
\begin{pageend_sol}{1.1 (p.\ 2--5)}
\textbf{1.} Add the first components and second components separately:
\[
\vect{v}+\vect{w} = \begin{bmatrix} 1+3 \\ 5+3 \end{bmatrix} = \begin{bmatrix} 4 \\ 8 \end{bmatrix}, \qquad \vect{w}+\vect{v} = \begin{bmatrix} 3+1 \\ 3+5 \end{bmatrix} = \begin{bmatrix} 4 \\ 8 \end{bmatrix}.
\]
\textbf{2.} Multiply the three vectors by $1, 4, -2$ and then add component by component:
\[
\begin{bmatrix} 1 \\ 0 \\ 3 \end{bmatrix} + \begin{bmatrix} 4 \\ 8 \\ 4 \end{bmatrix} + \begin{bmatrix} -4 \\ -6 \\ 2 \end{bmatrix} = \begin{bmatrix} 1 + 4 - 4 \\ 0 + 8 - 6 \\ 3 + 4 + 2 \end{bmatrix} = \begin{bmatrix} 1 \\ 2 \\ 9 \end{bmatrix}.
\]
\end{pageend_sol}
\clearpage

% ============================================================================
% CHAPTER 2: VECTOR SPACES AND SUBSPACES (3.1 - 3.6)
% ============================================================================
\chapter[Vector Spaces and Subspaces (3.1--3.6)]{Vector Spaces and Subspaces\bk{3.1--3.6}}

\begin{syll}
Syllabus topics: Vector spaces $\mathbf{R}^n$ and subspaces $\cdot$ Column space $C(A)$ $\cdot$ Nullspace $N(A)$ and special solutions $\cdot$ Echelon form $U$ and reduced row echelon form $R=\operatorname{rref}(A)$ $\cdot$ Rank $r$ $\cdot$ Complete solution $\vect{x} = \vect{x}_p + \vect{x}_n$ to $A\vect{x}=\vect{b}$ $\cdot$ Linear independence, spanning, basis, and dimension $\cdot$ Dimensions of the four fundamental subspaces (Fundamental Theorem of Linear Algebra, Part 1).
\end{syll}

\section{Spaces of Vectors and Subspaces\bk{3.1}}

\begin{goal}
Establish the algebraic test for a subspace and define the column space $C(A)$ of an $m \times n$ matrix $A$.
\end{goal}

\step{Step 1: State the subspace requirement.}
A subset of a vector space is a \term{subspace} if it is closed under vector addition and scalar multiplication:
\begin{align}
\text{(i)}\quad &\text{If } \vect{v} \text{ and } \vect{w} \text{ are in the subspace, then } \vect{v}+\vect{w} \text{ is in the subspace.} \\
\text{(ii)}\quad &\text{If } \vect{v} \text{ is in the subspace and } c \text{ is any scalar, then } c\vect{v} \text{ is in the subspace.}
\end{align}
Combine (i) and (ii) into one rule:
\begin{equation}
\eqbox{\text{A subspace containing } \vect{v} \text{ and } \vect{w} \text{ must contain all linear combinations } c\vect{v} + d\vect{w}.}
\end{equation}
Set $c=0$ in rule (ii) to get $0\vect{v} = \vect{0}$. Therefore, \term{every subspace must contain the zero vector $\vect{0}$}.

\step{Step 2: Define the column space $C(A)$ from $A\vect{x} = \vect{b}$.}
Let $A$ be an $m \times n$ matrix with columns $\vect{a}_1, \vect{a}_2, \dots, \vect{a}_n$ in $\mathbf{R}^m$. The product $A\vect{x}$ is a linear combination of the columns:
\begin{equation}
A\vect{x} = x_1\vect{a}_1 + x_2\vect{a}_2 + \dots + x_n\vect{a}_n.
\end{equation}
The \term{column space} $C(A)$ consists of all linear combinations of the columns of $A$. It is a subspace of $\mathbf{R}^m$.

\begin{result}[ (Solvability of $A\vect{x}=\vect{b}$)]
\begin{equation}
\text{The system } A\vect{x} = \vect{b} \text{ is solvable if and only if } \vect{b} \text{ is in the column space } C(A).
\end{equation}
\end{result}

\begin{example}{3.1.1--3.1.5 (p.\ 123--125)}{Subspace Tests and Column Spaces}
\begin{enumerate}
\item \textbf{Example 1:} Keep only the vectors $(x,y)$ in $\mathbf{R}^2$ whose components are positive or zero ($x\ge 0, y\ge 0$). Is this quarter-plane a subspace?
\item \textbf{Example 2:} Include also the vectors whose components are both negative. Is this union of two quarter-planes a subspace?
\item \textbf{Example 3:} Inside the vector space $\mathbf{M}$ of all $2$ by $2$ matrices, are the upper triangular matrices $\mathbf{U} = \left\{\begin{bmatrix}a & b\\0 & d\end{bmatrix}\right\}$ and diagonal matrices $\mathbf{D} = \left\{\begin{bmatrix}a & 0\\0 & d\end{bmatrix}\right\}$ subspaces?
\item \textbf{Example 4:} For $A = \begin{bmatrix}1 & 0\\4 & 3\\2 & 3\end{bmatrix}$, write $A\vect{x}$ as a combination of the columns and describe $C(A)$.
\item \textbf{Example 5:} Describe the column spaces (they are subspaces of $\mathbf{R}^2$) for:
\[
I = \begin{bmatrix} 1 & 0 \\ 0 & 1 \end{bmatrix} \quad \text{and} \quad A = \begin{bmatrix} 1 & 2 \\ 2 & 4 \end{bmatrix} \quad \text{and} \quad B = \begin{bmatrix} 1 & 2 & 3 \\ 0 & 0 & 4 \end{bmatrix}.
\]
\end{enumerate}
\end{example}

\vfill
\begin{pageend_sol}{3.1.1--3.1.5 (p.\ 123--125)}
\textbf{1.} \textbf{No.} The vector $(2,3)$ is included, but when you multiply by $c=-1$ you get $(-2,-3)$, which is outside the quarter-plane. Rule (ii) is violated.

\textbf{2.} \textbf{No.} Rule (ii) holds, but Rule (i) fails: the sum of $\vect{v}=(2,3)$ and $\vect{w}=(-3,-2)$ is $(-1,1)$, which is outside both quarter-planes.

\textbf{3.} \textbf{Yes.} The sum of two upper triangular matrices is upper triangular, and any scalar multiple stays upper triangular. The same holds for diagonal matrices $\mathbf{D}$ (and $\mathbf{D}$ is a subspace of $\mathbf{U}$).

\textbf{4.} $A\vect{x} = x_1\begin{bmatrix}1\\4\\2\end{bmatrix} + x_2\begin{bmatrix}0\\3\\3\end{bmatrix}$. The column space $C(A)$ is the plane in $\mathbf{R}^3$ through $(0,0,0)$ spanned by $(1,4,2)$ and $(0,3,3)$.

\textbf{5.} $C(I) = \mathbf{R}^2$ (the whole plane). $C(A)$ is the line through $(1,2)$ because column 2 is $2(\text{column 1})$. $C(B) = \mathbf{R}^2$ because combinations of $(1,0)$ and $(3,4)$ reach every $\vect{b}$ in $\mathbf{R}^2$.
\end{pageend_sol}
\clearpage

\section{The Nullspace of $A$: Solving $A\vect{x} = \vect{0}$\bk{3.2}}

\begin{goal}
Show that all solutions to $A\vect{x}=\vect{0}$ form a subspace $N(A)$ in $\mathbf{R}^n$, and construct the special solutions from the free variables.
\end{goal}

\step{Step 1: Verify that $N(A) = \{\vect{x} \in \mathbf{R}^n : A\vect{x} = \vect{0}\}$ is a subspace.}
Let $A\vect{x} = \vect{0}$ and $A\vect{y} = \vect{0}$. Use the distributive and associative rules of matrix multiplication:
\begin{align}
A(\vect{x} + \vect{y}) &= A\vect{x} + A\vect{y} = \vect{0} + \vect{0} = \vect{0} &&\why{closed under addition} \\
A(c\vect{x}) &= c(A\vect{x}) = c\vect{0} = \vect{0} &&\why{closed under scalar multiplication}
\end{align}
Thus the nullspace $N(A)$ is a subspace of $\mathbf{R}^n$.

\step{Step 2: Separate pivot variables from free variables via elimination.}
Forward elimination reduces $A$ to an upper triangular echelon matrix $U$, and upward elimination plus scaling reduces $U$ to the \term{reduced row echelon form} $R = \operatorname{rref}(A)$:
\begin{enumerate}
\item The $r$ columns with nonzero pivots are the \term{pivot columns}. Their variables are \term{pivot variables}.
\item The remaining $n - r$ columns without pivots are the \term{free columns}. Their variables are \term{free variables}.
\item To find each \term{special solution}, set one free variable to $1$, set all other free variables to $0$, and solve $U\vect{x} = \vect{0}$ (or $R\vect{x} = \vect{0}$) for the pivot variables.
\end{enumerate}

\begin{result}[ (Nullspace Dimension and Short Wide Matrices)]
If $A$ has $m$ rows and $n$ columns with $r$ pivots, then there are $n-r$ free variables and $n-r$ special solutions:
\begin{equation}
\vect{x}_{\text{complete}} = c_1\vect{s}_1 + c_2\vect{s}_2 + \dots + c_{n-r}\vect{s}_{n-r}.
\end{equation}
If $n > m$ (more unknowns than equations), there is at least $n-m \ge 1$ free variable, so $A\vect{x} = \vect{0}$ always has \term{nonzero solutions}.
\end{result}

\begin{example}{3.2.1--3.2.4 (p.\ 132--136)}{Finding Special Solutions and Nullspaces}
\begin{enumerate}
\item \textbf{Example 1:} Find two special solutions to $x + 2y + 3z = 0$ from $A = \begin{bmatrix}1 & 2 & 3\end{bmatrix}$.
\item \textbf{Example 2:} Describe the nullspace of $A = \begin{bmatrix}1 & 2\\3 & 6\end{bmatrix}$.
\item \textbf{Example 3:} Describe the nullspaces of $A = \begin{bmatrix}1 & 2\\3 & 8\end{bmatrix}$, $B = \begin{bmatrix}A\\2A\end{bmatrix}$, and $C = \begin{bmatrix}A & 2A\end{bmatrix} = \begin{bmatrix}1 & 2 & 2 & 4\\3 & 8 & 6 & 16\end{bmatrix}$.
\item \textbf{Example 4:} Find the nullspace of $U = \begin{bmatrix}1 & 5 & 7\\0 & 0 & 9\end{bmatrix}$.
\end{enumerate}
\end{example}

\vfill
\begin{pageend_sol}{3.2.1--3.2.4 (p.\ 132--136)}
\textbf{1.} Column 1 has the pivot ($x$ is the pivot variable); $y$ and $z$ are free. Set $(y,z)=(1,0)$ to get $x=-2$, and set $(y,z)=(0,1)$ to get $x=-3$:
\[
\vect{s}_1 = \begin{bmatrix}-2\\1\\0\end{bmatrix} \quad \text{and} \quad \vect{s}_2 = \begin{bmatrix}-3\\0\\1\end{bmatrix}.
\]
\textbf{2.} Subtract $3\times(\text{row 1})$ from row 2 to get $x_1 + 2x_2 = 0$. Set free variable $x_2=1$ to get $x_1=-2$. $N(A)$ contains all multiples of $\vect{s} = (-2,1)$.

\textbf{3.} $A$ reduces to $\begin{bmatrix}1 & 2\\0 & 2\end{bmatrix}$ (2 pivots, no free variables), so $N(A) = \mathbf{Z} = \{\vect{0}\}$. $B$ has the same 2 pivots, so $N(B) = \{\vect{0}\}$. $C$ reduces to $U = \begin{bmatrix}1 & 2 & 2 & 4\\0 & 2 & 0 & 4\end{bmatrix}$ and $R = \begin{bmatrix}1 & 0 & 2 & 0\\0 & 1 & 0 & 2\end{bmatrix}$. Free variables $x_3, x_4$ give $\vect{s}_1 = (-2,0,1,0)$ and $\vect{s}_2 = (0,-2,0,1)$.

\textbf{4.} Column 2 has no pivot, so $x_2$ is free. Set $x_2=1$. Then $9x_3=0 \implies x_3=0$, and $x_1+5(1)=0 \implies x_1=-5$. $N(U)$ is the line $\vect{x} = x_2(-5,1,0)$.
\end{pageend_sol}
\clearpage

\section{The Rank and the Row Reduced Form $R$\bk{3.3}}

\begin{goal}
Define the rank $r$ of a matrix and derive the block formula for the nullspace matrix $N$ from $R = \operatorname{rref}(A)$.
\end{goal}

\step{Step 1: Define rank $r$ and block-partition $R$.}
The \term{rank} $r$ of $A$ is the number of pivots. Suppose the first $r$ columns of $R$ are the pivot columns and the last $n-r$ columns are the free columns. Then $R$ has the block form:
\begin{equation}
R = \begin{bmatrix} I & F \\ \vect{0} & \vect{0} \end{bmatrix} \begin{array}{l} \scriptstyle r \text{ pivot rows} \\ \scriptstyle m-r \text{ zero rows} \end{array}
\end{equation}
where $I$ is the $r \times r$ identity matrix and $F$ is the $r \times (n-r)$ free part.

\step{Step 2: Solve $R\vect{x} = \vect{0}$ in block form.}
Write $\vect{x} = (\vect{x}_{\text{pivot}}, \vect{x}_{\text{free}})$. Then $R\vect{x} = \vect{0}$ becomes:
\begin{equation}
I\vect{x}_{\text{pivot}} + F\vect{x}_{\text{free}} = \vect{0} \implies \vect{x}_{\text{pivot}} = -F\vect{x}_{\text{free}}.
\end{equation}
Choose $\vect{x}_{\text{free}}$ to be the columns of the $(n-r)\times(n-r)$ identity matrix $I_{n-r}$. Then the $n-r$ special solutions form the columns of the \term{nullspace matrix} $N$:
\begin{result}[ (Nullspace Matrix $N$ from Reduced Echelon Form $R$)]
\begin{equation}
N = \begin{bmatrix} -F \\ I \end{bmatrix} \begin{array}{l} \scriptstyle r \text{ pivot variables} \\ \scriptstyle n-r \text{ free variables} \end{array} \qquad \text{satisfies} \qquad RN = \vect{0} \text{ and } AN = \vect{0}.
\end{equation}
Every rank-one matrix ($r=1$) has the special form $A = \vect{u}\vect{v}^\T = (\text{column})\times(\text{row})$.
\end{result}

\begin{example}{3.3.1--3.3.2 (p.\ 145--148)}{Rank One Matrices and Special Solutions from $R$}
\begin{enumerate}
\item \textbf{Example 1:} Find the rank and reduced row echelon form $R = \operatorname{rref}(A)$ for:
\[
\begin{bmatrix} 1 & 3 & 4 \\ 2 & 6 & 8 \end{bmatrix} \quad \text{and} \quad \begin{bmatrix} 0 & 3 \\ 0 & 5 \end{bmatrix} \quad \text{and} \quad \begin{bmatrix} 5 \\ 2 \end{bmatrix} \quad \text{and} \quad \begin{bmatrix} 6 \end{bmatrix}.
\]
\item \textbf{Example 2:} Find the special solutions of $R\vect{x} = x_1 + 2x_2 + 3x_3 = 0$ as the columns of $N$ for $R = \begin{bmatrix}1 & 2 & 3\end{bmatrix}$.
\end{enumerate}
\end{example}

\section{The Complete Solution to $A\vect{x} = \vect{b}$\bk{3.4}}

\begin{goal}
Derive the complete solution $\vect{x} = \vect{x}_p + \vect{x}_n$ to $A\vect{x} = \vect{b}$ and classify all linear systems by rank $r$.
\end{goal}

\step{Step 1: Combine one particular solution $\vect{x}_p$ with any nullspace solution $\vect{x}_n$.}
Suppose $\vect{x}_p$ is one particular solution ($A\vect{x}_p = \vect{b}$, found by setting all free variables to zero in $R\vect{x}=\vect{d}$), and $\vect{x}_n$ is any vector in the nullspace ($A\vect{x}_n = \vect{0}$). Add them and multiply by $A$:
\begin{equation}
A(\vect{x}_p + \vect{x}_n) = A\vect{x}_p + A\vect{x}_n = \vect{b} + \vect{0} = \vect{b}.
\end{equation}
Conversely, if $A\vect{x} = \vect{b}$, then $A(\vect{x} - \vect{x}_p) = \vect{b} - \vect{b} = \vect{0}$, so $\vect{x} - \vect{x}_p = \vect{x}_n$ is in the nullspace.

\begin{result}[ (Complete Solution and the Four Rank Cases)]
\begin{equation}
\vect{x}_{\text{complete}} = \vect{x}_p + \vect{x}_n = \vect{x}_p + c_1\vect{s}_1 + \dots + c_{n-r}\vect{s}_{n-r}.
\end{equation}
\begin{center}
\small
\begin{tabular}{@{}llll@{}}
\toprule
\textbf{Rank Case} & \textbf{Shape of $A$} & \textbf{Reduced Form $R$} & \textbf{Number of Solutions to $A\vect{x}=\vect{b}$} \\
\midrule
$r = m = n$ & Square and invertible & $R = [I]$ & $1$ solution for every $\vect{b}$ \\
$r = m < n$ & Short and wide (full row rank) & $R = [I \ \ F]$ & $\infty$ solutions for every $\vect{b}$ \\
$r = n < m$ & Tall and thin (full column rank) & $R = \begin{bmatrix}I\\\vect{0}\end{bmatrix}$ & $0$ or $1$ solution \\
$r < m$ and $r < n$ & Not full rank & $R = \begin{bmatrix}I & F\\\vect{0} & \vect{0}\end{bmatrix}$ & $0$ or $\infty$ solutions \\
\bottomrule
\end{tabular}
\end{center}
\end{result}

\vfill
\begin{pageend_sol}{3.3.1--3.3.2 (p.\ 145--148)}
\textbf{1.} All four matrices have \textbf{rank $r = 1$} because every row is a multiple of one pivot row. Their reduced row echelon forms are:
\[
R = \begin{bmatrix} 1 & 3 & 4 \\ 0 & 0 & 0 \end{bmatrix} \quad \text{and} \quad \begin{bmatrix} 0 & 1 \\ 0 & 0 \end{bmatrix} \quad \text{and} \quad \begin{bmatrix} 1 \\ 0 \end{bmatrix} \quad \text{and} \quad \begin{bmatrix} 1 \end{bmatrix}.
\]
\textbf{2.} Here $I = [1]$ and $F = \begin{bmatrix}2 & 3\end{bmatrix}$. The nullspace matrix is $N = \begin{bmatrix}-F\\I\end{bmatrix} = \begin{bmatrix}-2 & -3\\1 & 0\\0 & 1\end{bmatrix}$. The two special solutions are $(-2,1,0)$ and $(-3,0,1)$.
\end{pageend_sol}
\clearpage

\begin{example}{3.4.1--3.4.2 (p.\ 156--158)}{Full Column Rank and Full Row Rank Solutions}
\begin{enumerate}
\item \textbf{Example 1:} Find the condition on $(b_1, b_2, b_3)$ for $A\vect{x} = \vect{b}$ to be solvable, and find $\vect{x} = \vect{x}_p + \vect{x}_n$:
\[
A = \begin{bmatrix} 1 & 1 \\ 1 & 2 \\ -2 & -3 \end{bmatrix} \quad \text{and} \quad \vect{b} = \begin{bmatrix} b_1 \\ b_2 \\ b_3 \end{bmatrix}.
\]
\item \textbf{Example 2:} Find the complete solution $\vect{x} = \vect{x}_p + \vect{x}_n$ to the full row rank system ($r=m=2$):
\[
\begin{aligned}
x + y + z &= 3 \\
x + 2y - z &= 4.
\end{aligned}
\]
\end{enumerate}
\end{example}

\section{Independence, Basis, and Dimension\bk{3.5}}

\begin{goal}
Define linear independence, spanning, basis, and dimension, and prove that every basis for a vector space has the same number of vectors.
\end{goal}

\step{Step 1: Define linear independence.}
The columns of $A$ are \term{linearly independent} when the only solution to $A\vect{x} = \vect{0}$ is $\vect{x} = \vect{0}$. In general, a sequence of vectors $\vect{v}_1, \dots, \vect{v}_n$ is linearly independent if:
\begin{equation}
x_1\vect{v}_1 + x_2\vect{v}_2 + \dots + x_n\vect{v}_n = \vect{0} \quad \text{only happens when all } x_i = 0.
\end{equation}
If $n > m$, any $n$ vectors in $\mathbf{R}^m$ are automatically \term{linearly dependent} (because $A\vect{x}=\vect{0}$ has at least $n-m > 0$ free variables).

\step{Step 2: Define spanning and basis.}
\begin{itemize}
\item A set of vectors \term{spans} a space if their linear combinations fill the space.
\item The \term{row space} of $A$ is the subspace of $\mathbf{R}^n$ spanned by the rows of $A$. It is $C(A^\T)$.
\item A \term{basis} for a vector space is a sequence of vectors that are \term{linearly independent} and \term{span the space}.
\end{itemize}

\step{Step 3: Prove uniqueness of representation and equal size of all bases.}
If $\vect{v} = a_1\vect{v}_1 + \dots + a_n\vect{v}_n$ and $\vect{v} = b_1\vect{v}_1 + \dots + b_n\vect{v}_n$, subtract the two equations:
\begin{equation}
(a_1 - b_1)\vect{v}_1 + \dots + (a_n - b_n)\vect{v}_n = \vect{0} \implies a_i - b_i = 0 \implies a_i = b_i.
\end{equation}
Now suppose $\vect{v}_1, \dots, \vect{v}_m$ and $\vect{w}_1, \dots, \vect{w}_n$ are both bases for the same space with $n > m$. Each $\vect{w}_j$ is a combination of the $\vect{v}$'s, so $W = VA$ where $A$ is a short wide $m \times n$ matrix ($n > m$). Then $A\vect{x} = \vect{0}$ has a nonzero solution $\vect{x}$, which gives $W\vect{x} = VA\vect{x} = \vect{0}$. This contradicts the independence of the $\vect{w}$'s. By symmetry $m > n$ is also impossible, so $m = n$.

\begin{result}[ (Dimension of a Space)]
The \term{dimension} of a space is the number of vectors in every basis for that space.
\end{result}

\begin{example}{3.5.1--3.5.9 (p.\ 170--174)}{Independence, Row Space, and Bases}
\begin{enumerate}
\item \textbf{Example 1:} Show that the columns of $A = \begin{bmatrix}1 & 0 & 3\\2 & 1 & 5\\1 & 0 & 3\end{bmatrix}$ are dependent by solving $A\vect{x}=\vect{0}$.
\item \textbf{Examples 2--4:} Check spanning in $\mathbf{R}^2$ for $\vect{v}_1=\begin{bmatrix}1\\0\end{bmatrix}, \vect{v}_2=\begin{bmatrix}0\\1\end{bmatrix}, \vect{v}_3=\begin{bmatrix}4\\7\end{bmatrix}$ and $\vect{w}_1=\begin{bmatrix}1\\1\end{bmatrix}, \vect{w}_2=\begin{bmatrix}-1\\-1\end{bmatrix}$.
\item \textbf{Example 5:} Describe the column space and row space of $A = \begin{bmatrix}1 & 4\\2 & 7\\3 & 5\end{bmatrix}$.
\item \textbf{Examples 8--9:} Find bases for the column space and row space of $A = \begin{bmatrix}2 & 4\\3 & 6\end{bmatrix}$ and $R = \begin{bmatrix}1 & 2 & 0 & 3\\0 & 0 & 1 & 4\\0 & 0 & 0 & 0\end{bmatrix}$.
\end{enumerate}
\end{example}

\vfill
\begin{pageend_sol}{3.4.1--3.4.2 and 3.5.1--3.5.9}
\textbf{Ex.\ 3.4.1:} Reduce $[A \ \ \vect{b}] \to \begin{bmatrix}1 & 0 & 2b_1-b_2\\0 & 1 & b_2-b_1\\0 & 0 & b_3+b_1+b_2\end{bmatrix}$. Solvability condition: $b_1+b_2+b_3=0$. Unique solution: $\vect{x} = \vect{x}_p + \vect{x}_n = \begin{bmatrix}2b_1-b_2\\b_2-b_1\end{bmatrix} + \begin{bmatrix}0\\0\end{bmatrix}$.

\textbf{Ex.\ 3.4.2:} Reduce $[A \ \ \vect{b}] \to \begin{bmatrix}1 & 0 & 3 & 2\\0 & 1 & -2 & 1\end{bmatrix}$. Set $x_3=0$ for $\vect{x}_p=(2,1,0)$ and $x_3=1$ for $\vect{s}=(-3,2,1)$: $\vect{x} = \begin{bmatrix}2\\1\\0\end{bmatrix} + x_3\begin{bmatrix}-3\\2\\1\end{bmatrix}$.

\textbf{Ex.\ 3.5.1--3.5.9:} \textbf{(1)} $R = \begin{bmatrix}1 & 0 & 3\\0 & 1 & -1\\0 & 0 & 0\end{bmatrix}$ gives $\vect{x}=(-3,1,1) \neq \vect{0}$, so the columns are dependent ($r=2$). \textbf{(2--4)} $\vect{v}_1,\vect{v}_2$ span $\mathbf{R}^2$ (and are a basis); $\vect{v}_1,\vect{v}_2,\vect{v}_3$ span $\mathbf{R}^2$ (dependent); $\vect{w}_1,\vect{w}_2$ only span a line. \textbf{(5)} $C(A)$ is a plane in $\mathbf{R}^3$; the row space $C(A^\T)$ is all of $\mathbf{R}^2$. \textbf{(8)} For $\begin{bmatrix}2 & 4\\3 & 6\end{bmatrix}$, basis for $C(A)$ is $\begin{bmatrix}2\\3\end{bmatrix}$; basis for row space is $\begin{bmatrix}1\\2\end{bmatrix}$. \textbf{(9)} For $R$, pivot columns $\begin{bmatrix}1\\0\\0\end{bmatrix}, \begin{bmatrix}0\\1\\0\end{bmatrix}$ are a basis for $C(R)$; nonzero rows $(1,2,0,3)$ and $(0,0,1,4)$ are a basis for the row space.
\end{pageend_sol}
\clearpage

\section{Dimensions of the Four Subspaces\bk{3.6}}

\begin{goal}
Derive Part 1 of the Fundamental Theorem of Linear Algebra: the dimensions and bases of the four fundamental subspaces of an $m \times n$ matrix $A$.
\end{goal}

\step{Step 1: Define the four fundamental subspaces.}
\begin{enumerate}
\item \term{Row space} $C(A^\T)$: all combinations of the rows of $A$. It is a subspace of $\mathbf{R}^n$.
\item \term{Column space} $C(A)$: all combinations of the columns of $A$. It is a subspace of $\mathbf{R}^m$.
\item \term{Nullspace} $N(A)$: all solutions $\vect{x}$ to $A\vect{x} = \vect{0}$. It is a subspace of $\mathbf{R}^n$.
\item \term{Left nullspace} $N(A^\T)$: all solutions $\vect{y}$ to $A^\T\vect{y} = \vect{0}$ (or $\vect{y}^\T A = \vect{0}^\T$). It is a subspace of $\mathbf{R}^m$.
\end{enumerate}

\step{Step 2: Count the basis vectors from $E A = R$.}
Row reduction multiplies $A$ by an invertible $m \times m$ elimination matrix $E$ to produce $R$ with $r$ pivots:
\begin{itemize}
\item \textbf{Row space:} Row operations do not change the row space, so $C(A^\T) = C(R^\T)$. The $r$ nonzero pivot rows of $R$ are independent and span the row space. Thus $\dim C(A^\T) = r$.
\item \textbf{Column space:} $A\vect{x} = \vect{0}$ holds if and only if $R\vect{x} = \vect{0}$. Therefore, the $r$ pivot columns of $A$ (in the same column positions as the $r$ pivot columns of $R$) are independent and span $C(A)$. Thus $\dim C(A) = r$.
\item \textbf{Nullspace:} There are $n - r$ free variables, producing $n - r$ independent special solutions that span $N(A) = N(R)$. Thus $\dim N(A) = n - r$.
\item \textbf{Left nullspace:} The last $m - r$ rows of $R$ are zero rows. Since $EA = R$, the last $m - r$ rows of $E$ are the independent combinations of the rows of $A$ that produce $\vect{0}^\T$. Thus $\dim N(A^\T) = m - r$.
\end{itemize}

\begin{result}[ (Fundamental Theorem of Linear Algebra, Part 1)]
For any $m \times n$ matrix $A$ of rank $r$:
\begin{align}
\dim C(A) &= \dim C(A^\T) = r &&\why{column rank = row rank} \\
\dim C(A^\T) + \dim N(A) &= r + (n - r) = n &&\why{inside $\mathbf{R}^n$} \\
\dim C(A) + \dim N(A^\T) &= r + (m - r) = m &&\why{inside $\mathbf{R}^m$}
\end{align}
\end{result}

\begin{pic}
\begin{tikzpicture}[scale=0.75]
  % Left: R^n
  \draw[thick,teal,fill=tealsoft!40] (0,0) -- (-2,2) -- (0,4) -- (2,2) -- cycle;
  \draw[thick,amber,fill=ambersoft!50] (0,0) -- (-1.6,-1.6) -- (0,-3.2) -- (1.6,-1.6) -- cycle;
  \fill[ink] (0,0) circle (2.5pt);
  \node[font=\small\bfseries,teal] at (0,2.4) {row space $C(A^\T)$};
  \node[font=\footnotesize] at (0,1.7) {all $A^\T\vect{y}$ ($\dim = r$)};
  \node[font=\small\bfseries,amber!80!black] at (0,-1.4) {nullspace $N(A)$};
  \node[font=\footnotesize] at (0,-2.1) {$A\vect{x}=\vect{0}$ ($\dim = n-r$)};
  \node[font=\bfseries] at (-2.4,0) {$\mathbf{R}^n$};

  % Right: R^m
  \begin{scope}[xshift=9cm]
    \draw[thick,teal,fill=tealsoft!40] (0,0) -- (-1.6,2.4) -- (0.8,4) -- (2.4,1.6) -- cycle;
    \draw[thick,plum,fill=plumsoft!50] (0,0) -- (-2.4,-1) -- (-1.6,-2.4) -- (0.8,-1.4) -- cycle;
    \fill[ink] (0,0) circle (2.5pt);
    \node[font=\small\bfseries,teal] at (0.4,2.3) {column space $C(A)$};
    \node[font=\footnotesize] at (0.4,1.6) {all $A\vect{x}$ ($\dim = r$)};
    \node[font=\small\bfseries,plum] at (-0.8,-1.0) {left nullspace $N(A^\T)$};
    \node[font=\footnotesize] at (-0.8,-1.7) {$A^\T\vect{y}=\vect{0}$ ($\dim = m-r$)};
    \node[font=\bfseries] at (2.2,-0.2) {$\mathbf{R}^m$};
  \end{scope}
\end{tikzpicture}
\figcap{The dimensions of the Four Fundamental Subspaces in $\mathbf{R}^n$ and $\mathbf{R}^m$ \bookref{Fig.\ 3.5}.}
\end{pic}

\begin{example}{3.6.1--3.6.2 (p.\ 188)}{Dimensions and Bases of the Four Subspaces}
\begin{enumerate}
\item \textbf{Example 1:} Find the dimensions of the four subspaces for $A = \begin{bmatrix}1 & 2 & 3\end{bmatrix}$ ($m=1, n=3, r=1$).
\item \textbf{Example 2:} Find the four subspaces and their dimensions for $A = \begin{bmatrix}1 & 2 & 3\\2 & 4 & 6\end{bmatrix}$ ($m=2, n=3, r=1$).
\end{enumerate}
\end{example}

\vfill
\begin{pageend_sol}{3.6.1--3.6.2 (p.\ 188)}
\textbf{1.} For $A = \begin{bmatrix}1 & 2 & 3\end{bmatrix}$ ($m=1, n=3, r=1$): Row space is the line through $(1,2,3)$ in $\mathbf{R}^3$ ($\dim = 1$). Nullspace is the plane $x_1+2x_2+3x_3=0$ in $\mathbf{R}^3$ ($\dim = 3-1 = 2$). Column space is all of $\mathbf{R}^1$ ($\dim = 1$). Left nullspace $N(A^\T)$ is $\mathbf{Z} = \{0\}$ ($\dim = 1-1 = 0$).

\textbf{2.} For $A = \begin{bmatrix}1 & 2 & 3\\2 & 4 & 6\end{bmatrix}$ ($m=2, n=3, r=1$): Row space is the line through $(1,2,3)$ ($\dim = 1$). Nullspace is the plane $x_1+2x_2+3x_3=0$ ($\dim = 2$). Column space is the line through $\begin{bmatrix}1\\2\end{bmatrix}$ in $\mathbf{R}^2$ ($\dim = 1$). Left nullspace is the line through $\vect{y} = \begin{bmatrix}2\\-1\end{bmatrix}$ in $\mathbf{R}^2$ ($\dim = 2-1 = 1$).
\end{pageend_sol}
\clearpage

% ============================================================================
% END OF MODULE 1: ALL UNIQUE SOLVED PROBLEMS (WORKED EXAMPLES)
% ============================================================================
\chapter[Module 1: Unique Solved Problems]{Module 1: All Unique Solved Problems (Worked Examples)}

\begin{syll}
All unique Worked Examples from Module 1 (Sections 1.1 and 3.1--3.6), with problem statements copied verbatim and full step-by-step solutions.
\end{syll}

\begin{example}{1.1 A}{Describing a Plane of Linear Combinations}
The linear combinations of $\vect{v} = (1, 1, 0)$ and $\vect{w} = (0, 1, 1)$ fill a plane. \emph{Describe that plane.} Find a vector that is \emph{not} a combination of $\vect{v}$ and $\vect{w}$.
\end{example}
\textbf{Solution:} The combinations $c\vect{v} + d\vect{w}$ fill a plane in $\mathbf{R}^3$ for any $c$ and $d$:
\[
c\vect{v} + d\vect{w} = c\begin{bmatrix} 1 \\ 1 \\ 0 \end{bmatrix} + d\begin{bmatrix} 0 \\ 1 \\ 1 \end{bmatrix} = \begin{bmatrix} c \\ c + d \\ d \end{bmatrix}.
\]
In every vector in this plane, the second component $c+d$ equals the sum of the first component $c$ and the third component $d$ (so $x - y + z = 0$, perpendicular to $\vect{n}=(1,-1,1)$). Four vectors in the plane are $(0,0,0)$, $(2,3,1)$, $(5,7,2)$, and $(\pi,2\pi,\pi)$. The vector $(1,2,3)$ is \emph{not} in the plane because $2 \neq 1 + 3$.

\begin{example}{1.1 B}{Integer and Nonnegative Combinations}
For $\vect{v} = (1, 0)$ and $\vect{w} = (0, 1)$, describe all points $c\vect{v}$ with (1) \emph{whole numbers} $c$ (2) \emph{nonnegative} $c \ge 0$. Then add all vectors $d\vect{w}$ and describe all $c\vect{v} + d\vect{w}$.
\end{example}
\textbf{Solution:}
\begin{enumerate}
\item[(1)] The vectors $c\vect{v} = (c,0)$ with whole numbers $c$ are \textbf{equally spaced points} along the $x$ axis: $(-2,0), (-1,0), (0,0), (1,0), (2,0), \dots$
\item[(2)] The vectors $c\vect{v}$ with $c \ge 0$ fill a \textbf{half-line}: the positive $x$ axis starting at $(0,0)$.
\item[(1$'$)] Adding all $d\vect{w} = (0,d)$ puts a vertical line through each point $c\vect{v}$, giving infinitely many \textbf{parallel vertical lines} ($x = \text{whole number } c$, any $y = d$).
\item[(2$'$)] Adding all $d\vect{w}$ to every $c\vect{v}$ on the half-line produces the \textbf{right half-plane} ($x \ge 0$, any $y$).
\end{enumerate}

\begin{example}{1.1 C}{Modeling a Linear Combination by Scalar Equations}
Find two equations for the unknowns $c$ and $d$ so that the linear combination $c\vect{v} + d\vect{w}$ equals the vector $\vect{b}$:
\[
\vect{v} = \begin{bmatrix} 2 \\ -1 \end{bmatrix} \qquad \vect{w} = \begin{bmatrix} -1 \\ 2 \end{bmatrix} \qquad \vect{b} = \begin{bmatrix} 1 \\ 0 \end{bmatrix}.
\]
\end{example}
\textbf{Solution:} Write the vector equation $c\begin{bmatrix}2\\-1\end{bmatrix} + d\begin{bmatrix}-1\\2\end{bmatrix} = \begin{bmatrix}1\\0\end{bmatrix}$. Match the two components separately to get two scalar equations:
\[
\begin{aligned}
2c - d &= 1 \\
-c + 2d &= 0
\end{aligned}
\implies c = \frac{2}{3}, \quad d = \frac{1}{3}.
\]

\begin{example}{3.1 A}{Constructing a Matrix with Prescribed Column Space}
We are given three different vectors $\vect{b}_1, \vect{b}_2, \vect{b}_3$. Construct a matrix so that the equations $A\vect{x} = \vect{b}_1$ and $A\vect{x} = \vect{b}_2$ are solvable, but $A\vect{x} = \vect{b}_3$ is not solvable. How can you decide if this is possible? How could you construct $A$?
\end{example}
\textbf{Solution:} Make $\vect{b}_1$ and $\vect{b}_2$ the two columns of $A = \begin{bmatrix}\vect{b}_1 & \vect{b}_2\end{bmatrix}$. Then $A\vect{x}=\vect{b}_1$ is solved by $\vect{x}=(1,0)$ and $A\vect{x}=\vect{b}_2$ is solved by $\vect{x}=(0,1)$. Now test if $A\vect{x}=\vect{b}_3$ is solvable (that is, whether $\vect{b}_3$ is a linear combination of $\vect{b}_1$ and $\vect{b}_2$). If \emph{no}, then $A = \begin{bmatrix}\vect{b}_1 & \vect{b}_2\end{bmatrix}$ is the desired matrix. If \emph{yes}, it is \emph{not possible} to construct $A$, because any subspace containing $\vect{b}_1$ and $\vect{b}_2$ must also contain their combination $\vect{b}_3$.

\begin{example}{3.1 B}{Subspaces of Vector Spaces}
Describe a subspace $\mathbf{S}$ of each vector space $\mathbf{V}$, and then a subspace $\mathbf{SS}$ of $\mathbf{S}$.
\begin{itemize}
\item $\mathbf{V}_1 =$ all combinations of $(1, 1, 0, 0)$ and $(1, 1, 1, 0)$ and $(1, 1, 1, 1)$
\item $\mathbf{V}_2 =$ all vectors perpendicular to $\vect{u} = (1, 2, 1)$, so $\vect{u} \cdot \vect{v} = 0$
\item $\mathbf{V}_3 =$ all symmetric $2$ by $2$ matrices (a subspace of $\mathbf{M}$)
\item $\mathbf{V}_4 =$ all solutions to the equation $\diff^4 y / \diff x^4 = 0$ (a subspace of $\mathbf{F}$)
\end{itemize}
Describe each $\mathbf{V}$ two ways: \emph{All combinations of \dots, all solutions of the equations \dots}
\end{example}
\textbf{Solution:}
\begin{itemize}
\item For $\mathbf{V}_1$: $\mathbf{S} =$ all combinations of $(1,1,0,0)$ and $(1,1,1,0)$; $\mathbf{SS} =$ all multiples $(c,c,0,0)$ of $(1,1,0,0)$.\\
Two descriptions of $\mathbf{V}_1$: all combinations of the 3 given vectors; all solutions in $\mathbf{R}^4$ of $v_1 - v_2 = 0$.
\item For $\mathbf{V}_2$: $\mathbf{S} =$ the line through $(1,-1,1)$; $\mathbf{SS} = \mathbf{Z} = \{(0,0,0)\}$.\\
Two descriptions of $\mathbf{V}_2$: all combinations of $(1,0,-1)$ and $(1,-1,1)$; all solutions of $v_1 + 2v_2 + v_3 = 0$.
\item For $\mathbf{V}_3$: $\mathbf{S} =$ all diagonal $2\times 2$ matrices; $\mathbf{SS} =$ all multiples $cI$ of the identity matrix.\\
Two descriptions of $\mathbf{V}_3$: all combinations of $\begin{bmatrix}1&0\\0&0\end{bmatrix}, \begin{bmatrix}0&1\\1&0\end{bmatrix}, \begin{bmatrix}0&0\\0&1\end{bmatrix}$; all matrices $\begin{bmatrix}a&b\\c&d\end{bmatrix}$ with $b = c$.
\item For $\mathbf{V}_4$: $\mathbf{S} =$ all quadratic polynomials $y = a + bx + cx^2$; $\mathbf{SS} =$ all linear polynomials $y = a + bx$.\\
Two descriptions of $\mathbf{V}_4$: all combinations of $1, x, x^2, x^3$; all solutions to $\diff^4 y / \diff x^4 = 0$.
\end{itemize}

\begin{example}{3.2 A}{Constructing $R$ from Special Solutions}
Create a $3$ by $4$ matrix whose special solutions to $A\vect{x} = \vect{0}$ are $\vect{s}_1$ and $\vect{s}_2$:
\[
\vect{s}_1 = \begin{bmatrix} -3 \\ 1 \\ 0 \\ 0 \end{bmatrix} \quad \text{and} \quad \vect{s}_2 = \begin{bmatrix} -2 \\ 0 \\ -6 \\ 1 \end{bmatrix} \qquad \begin{array}{l} \text{pivot columns } 1 \text{ and } 3 \\ \text{free variables } x_2 \text{ and } x_4 \end{array}
\]
You could create the matrix $A$ in row reduced form $R$. Then describe \emph{all} possible matrices $A$ with the required nullspace $N(A) =$ all combinations of $\vect{s}_1$ and $\vect{s}_2$.
\end{example}
\textbf{Solution:} Put $1$'s in pivot columns $1$ and $3$, and put the negatives of the pivot entries of $\vect{s}_1, \vect{s}_2$ into free columns $2$ and $4$ (row 3 is all zeros because there are only 2 pivots):
\[
R = \begin{bmatrix} 1 & 3 & 0 & 2 \\ 0 & 0 & 1 & 6 \\ 0 & 0 & 0 & 0 \end{bmatrix} \quad \text{has} \quad R\vect{s}_1 = \vect{0} \quad \text{and} \quad R\vect{s}_2 = \vect{0}.
\]
Any matrix $A = E^{-1}R$ obtained by multiplying $R$ on the left by any invertible $3\times 3$ matrix has the exact same nullspace.

\begin{example}{3.2 B}{Special Solutions and Complete Nullspace Solutions}
Find the special solutions and describe the \emph{complete solution} to $A\vect{x} = \vect{0}$ for:
\[
A_1 = \begin{bmatrix} 0 & 0 & 0 & 0 \\ 0 & 0 & 0 & 0 \end{bmatrix} \qquad A_2 = \begin{bmatrix} 3 & 6 \\ 1 & 2 \end{bmatrix} \qquad A_3 = \begin{bmatrix} A_2 & A_2 \end{bmatrix}.
\]
Which are the pivot columns? Which are the free variables? What is $R$ in each case?
\end{example}
\textbf{Solution:}
\begin{itemize}
\item For $A_1$: No pivot columns; all four variables $x_1, x_2, x_3, x_4$ are free; $R_1 = A_1$. The four special solutions are the columns $\vect{s}_1, \vect{s}_2, \vect{s}_3, \vect{s}_4$ of $I_4$. Complete solution: any $\vect{x} \in \mathbf{R}^4$.
\item For $A_2$: Column 1 is the pivot column; $x_2$ is free; $R_2 = \begin{bmatrix}1 & 2\\0 & 0\end{bmatrix}$. Special solution: $\vect{s} = (-2,1)$. Complete solution: $\vect{x} = c(-2,1)$.
\item For $A_3$: Column 1 is the only pivot column; $x_2, x_3, x_4$ are free; $R_3 = \begin{bmatrix}1 & 2 & 1 & 2\\0 & 0 & 0 & 0\end{bmatrix}$. The three special solutions are $\vect{s}_1 = (-2,1,0,0)$, $\vect{s}_2 = (-1,0,1,0)$, $\vect{s}_3 = (-2,0,0,1)$. Complete solution: $\vect{x} = c_1\vect{s}_1 + c_2\vect{s}_2 + c_3\vect{s}_3$.
\end{itemize}

\begin{example}{3.3 A}{Row Reduced Form of a Singular Second-Difference Matrix}
Find the reduced echelon form of $A$. What is the rank? What is the special solution to $A\vect{x} = \vect{0}$?
\[
A = \begin{bmatrix} 1 & -1 & 0 & 0 \\ -1 & 2 & -1 & 0 \\ 0 & -1 & 2 & -1 \\ 0 & 0 & -1 & 1 \end{bmatrix}.
\]
\end{example}
\textbf{Solution:} Add row 1 to row 2, then row 2 to row 3, then row 3 to row 4 to get $U$. Then add row 3 to row 2, and row 2 to row 1 to get $R$:
\[
U = \begin{bmatrix} 1 & -1 & 0 & 0 \\ 0 & 1 & -1 & 0 \\ 0 & 0 & 1 & -1 \\ 0 & 0 & 0 & 0 \end{bmatrix} \longrightarrow R = \begin{bmatrix} 1 & 0 & 0 & -1 \\ 0 & 1 & 0 & -1 \\ 0 & 0 & 1 & -1 \\ 0 & 0 & 0 & 0 \end{bmatrix}.
\]
The rank is $r = 3$. There is $n-r = 1$ free variable ($x_4$). The special solution is $\vect{s} = (1,1,1,1)$.

\begin{example}{3.3 B}{Factoring Rank-One and Rank-Two Matrices}
Factor these rank one matrices into $A = \vect{u}\vect{v}^\T =$ column times row:
\[
A = \begin{bmatrix} 1 & 2 & 3 \\ 2 & 4 & 6 \\ 3 & 6 & 9 \end{bmatrix} \qquad A = \begin{bmatrix} a & b \\ c & d \end{bmatrix} \quad (\text{find } d \text{ from } a, b, c \text{ if } a \neq 0).
\]
Split this rank two matrix into $\vect{u}_1\vect{v}_1^\T + \vect{u}_2\vect{v}_2^\T = (3 \text{ by } 2)$ times $(2 \text{ by } 4)$ using $R$:
\[
A = \begin{bmatrix} 1 & 1 & 0 & 2 \\ 1 & 2 & 0 & 3 \\ 2 & 3 & 0 & 5 \end{bmatrix} = \begin{bmatrix} 1 & 1 & 0 \\ 1 & 2 & 0 \\ 2 & 3 & 1 \end{bmatrix} \begin{bmatrix} 1 & 0 & 0 & 1 \\ 0 & 1 & 0 & 1 \\ 0 & 0 & 0 & 0 \end{bmatrix} = E^{-1}R.
\]
\end{example}
\textbf{Solution:} For the $3\times 3$ matrix, $A = \begin{bmatrix}1\\2\\3\end{bmatrix}\begin{bmatrix}1 & 2 & 3\end{bmatrix}$. For the $2\times 2$ matrix, $\begin{bmatrix}a & b\\c & d\end{bmatrix} = \begin{bmatrix}1\\c/a\end{bmatrix}\begin{bmatrix}a & b\end{bmatrix} = \begin{bmatrix}a & b\\c & bc/a\end{bmatrix}$, so $d = bc/a$. For the rank-2 matrix, multiply the first 2 pivot columns of $A$ by the first 2 nonzero rows of $R$:
\[
\begin{bmatrix} 1 & 1 & 0 & 2 \\ 1 & 2 & 0 & 3 \\ 2 & 3 & 0 & 5 \end{bmatrix} = \begin{bmatrix} 1 \\ 1 \\ 2 \end{bmatrix}\begin{bmatrix} 1 & 0 & 0 & 1 \end{bmatrix} + \begin{bmatrix} 1 \\ 2 \\ 3 \end{bmatrix}\begin{bmatrix} 0 & 1 & 0 & 1 \end{bmatrix}.
\]

\begin{example}{3.3 C}{Rank and Special Solutions Depending on a Parameter $c$}
Find the row reduced form $R$ and the rank $r$ of $A$ and $B$ (those depend on $c$). Which are the pivot columns of $A$? What are the special solutions and the matrix $N$?
\[
A = \begin{bmatrix} 1 & 2 & 1 \\ 3 & 6 & 3 \\ 4 & 8 & c \end{bmatrix} \quad \text{and} \quad B = \begin{bmatrix} c & c \\ c & c \end{bmatrix}.
\]
\end{example}
\textbf{Solution:}
\begin{itemize}
\item For $A$: If $c \neq 4$, rank $r = 2$, pivot columns are $1$ and $3$, $R = \begin{bmatrix}1 & 2 & 0\\0 & 0 & 1\\0 & 0 & 0\end{bmatrix}$, and $N = \begin{bmatrix}-2\\1\\0\end{bmatrix}$. If $c = 4$, rank $r = 1$, pivot column is $1$, $R = \begin{bmatrix}1 & 2 & 1\\0 & 0 & 0\\0 & 0 & 0\end{bmatrix}$, and $N = \begin{bmatrix}-2 & -1\\1 & 0\\0 & 1\end{bmatrix}$.
\item For $B$: If $c \neq 0$, rank $r = 1$, $R = \begin{bmatrix}1 & 1\\0 & 0\end{bmatrix}$, and $N = \begin{bmatrix}-1\\1\end{bmatrix}$. If $c = 0$, rank $r = 0$, $R = \begin{bmatrix}0 & 0\\0 & 0\end{bmatrix}$, and $N = \begin{bmatrix}1 & 0\\0 & 1\end{bmatrix}$ (nullspace is $\mathbf{R}^2$).
\end{itemize}

\begin{example}{3.4 A}{Full Picture of $A\vect{x}=\vect{b}$ for a Rank-2 Matrix}
$A$ has rank $2$:
\[
A\vect{x} = \vect{b} \quad \text{is} \quad \begin{aligned}
x_1 + 2x_2 + 3x_3 + 5x_4 &= b_1 \\
2x_1 + 4x_2 + 8x_3 + 12x_4 &= b_2 \\
3x_1 + 6x_2 + 7x_3 + 13x_4 &= b_3
\end{aligned}
\]
1. Reduce $[A \ \ \vect{b}]$ to $[U \ \ \vect{c}]$. \quad 2. Find the condition on $b_1, b_2, b_3$ for $A\vect{x}=\vect{b}$ to have a solution. \quad 3. Describe $C(A)$. \quad 4. Describe $N(A)$. \quad 5. Find a particular solution to $A\vect{x} = (0, 6, -6)$ and the complete solution. \quad 6. Reduce $[U \ \ \vect{c}]$ to $[R \ \ \vect{d}]$.
\end{example}
\textbf{Solution:}
\begin{enumerate}
\item Subtract $2\times(\text{row 1})$ from row 2, $3\times(\text{row 1})$ from row 3, then add row 2 to row 3:
\[
[A \ \ \vect{b}] \to \begin{bmatrix} 1 & 2 & 3 & 5 & b_1 \\ 0 & 0 & 2 & 2 & b_2 - 2b_1 \\ 0 & 0 & 0 & 0 & b_3 + b_2 - 5b_1 \end{bmatrix} = [U \ \ \vect{c}].
\]
\item Solvability condition: $b_3 + b_2 - 5b_1 = 0$.
\item $C(A)$ is the plane in $\mathbf{R}^3$ spanned by pivot columns $(1,2,3)$ and $(3,8,7)$, which has equation $b_3 + b_2 - 5b_1 = 0$.
\item Free variables $x_2, x_4$ give $\vect{s}_1 = (-2,1,0,0)$ and $\vect{s}_2 = (-2,0,-1,1)$. $N(A)$ contains all $\vect{x}_n = c_1\vect{s}_1 + c_2\vect{s}_2$.
\item For $\vect{b} = (0,6,-6)$, set $x_2=x_4=0$ to get $\vect{x}_p = (-9,0,3,0)$. Complete solution: $\vect{x} = \vect{x}_p + c_1\vect{s}_1 + c_2\vect{s}_2$.
\item Divide row 2 by 2 and eliminate upward for $\vect{b}=(0,6,-6)$:
\[
[R \ \ \vect{d}] = \begin{bmatrix} 1 & 2 & 0 & 2 & -9 \\ 0 & 0 & 1 & 1 & 3 \\ 0 & 0 & 0 & 0 & 0 \end{bmatrix}.
\]
\end{enumerate}

\begin{example}{3.4 B}{Deducing Matrix Properties from Solution Sets}
If you have this information about the solutions to $A\vect{x} = \vect{b}$ for a \emph{specific} $\vect{b}$, what does that tell you about the shape of $A$ (and $A$ itself)?
1. There is exactly one solution. \quad 2. All solutions have the form $\vect{x} = \begin{bmatrix}2\\1\end{bmatrix} + c\begin{bmatrix}1\\1\end{bmatrix}$. \quad 3. There are no solutions. \quad 4. All solutions have the form $\vect{x} = \begin{bmatrix}1\\1\\0\end{bmatrix} + c\begin{bmatrix}1\\0\\1\end{bmatrix}$. \quad 5. There are infinitely many solutions.
\end{example}
\textbf{Solution:}
\textbf{1.} $A$ has full column rank $r = n$, $N(A) = \{\vect{0}\}$, and $m \ge n$.
\textbf{2.} $A$ has $n=2$ columns, rank $r=1$, column 2 $= -(\text{column 1})$, and $\vect{b} = 2(\text{col 1}) + (\text{col 2}) = \text{col 1}$.
\textbf{3.} $\vect{b}$ is not in $C(A)$, so $r < m$ and $\vect{b} \neq \vect{0}$.
\textbf{4.} $A$ has $n=3$ columns, rank $r = 3 - 1 = 2$, $m \ge 2$ rows, col 3 $= -(\text{col 1})$, and $\vect{b} = \text{col 1} + \text{col 2}$.
\textbf{5.} $r < n$ (not full column rank) and $\vect{b}$ is in $C(A)$.

\begin{example}{3.4 C}{Complete Solution and Left Nullspace Condition}
Find the complete solution $\vect{x} = \vect{x}_p + \vect{x}_n$ by forward elimination on $[A \ \ \vect{b}]$:
\[
\begin{bmatrix} 1 & 2 & 1 & 0 \\ 2 & 4 & 4 & 8 \\ 4 & 8 & 6 & 8 \end{bmatrix} \begin{bmatrix} x_1 \\ x_2 \\ x_3 \\ x_4 \end{bmatrix} = \begin{bmatrix} 4 \\ 2 \\ 10 \end{bmatrix}.
\]
Find numbers $y_1, y_2, y_3$ so that $y_1(\text{row 1}) + y_2(\text{row 2}) + y_3(\text{row 3}) = \text{zero row}$. Check $y_1b_1 + y_2b_2 + y_3b_3 = 0$.
\end{example}
\textbf{Solution:} Row reduction gives:
\[
[A \ \ \vect{b}] \to \begin{bmatrix} 1 & 2 & 1 & 0 & 4 \\ 0 & 0 & 2 & 8 & -6 \\ 0 & 0 & 0 & 0 & 0 \end{bmatrix} \to \begin{bmatrix} 1 & 2 & 0 & -4 & 7 \\ 0 & 0 & 1 & 4 & -3 \\ 0 & 0 & 0 & 0 & 0 \end{bmatrix} = [R \ \ \vect{d}].
\]
Thus $\vect{x}_p = (7,0,-3,0)$, $\vect{s}_1 = (-2,1,0,0)$, $\vect{s}_2 = (4,0,-4,1)$, and $\vect{x} = \vect{x}_p + c_1\vect{s}_1 + c_2\vect{s}_2$.
Since $2(\text{row 1}) + 1(\text{row 2}) - 1(\text{row 3}) = (0,0,0,0)$, we have $\vect{y} = (2,1,-1)$ in $N(A^\T)$, and $2(4) + 1(2) - 1(10) = 0$.

\begin{example}{3.5 A}{Basis and Dimension of a Subspace}
Start with the vectors $\vect{v}_1 = (1, 2, 0)$ and $\vect{v}_2 = (2, 3, 0)$. (a) Are they linearly independent? (b) Are they a basis for any space? (c) What space $\mathbf{V}$ do they span? (d) What is the dimension of $\mathbf{V}$? (e) Which matrices $A$ have $\mathbf{V}$ as their column space? (f) Which matrices have $\mathbf{V}$ as their nullspace? (g) Describe all vectors $\vect{v}_3$ that complete a basis $\vect{v}_1, \vect{v}_2, \vect{v}_3$ for $\mathbf{R}^3$.
\end{example}
\textbf{Solution:} (a) Yes, only $0\vect{v}_1 + 0\vect{v}_2 = \vect{0}$. (b) Yes, for the space they span. (c) The $xy$ plane in $\mathbf{R}^3$ of all vectors $(x,y,0)$. (d) Dimension is $2$. (e) Any $3\times n$ matrix $A$ of rank $2$ whose columns are combinations of $\vect{v}_1, \vect{v}_2$. (f) Any $m\times 3$ matrix $B$ of rank $1$ whose rows are nonzero multiples of $\begin{bmatrix}0 & 0 & 1\end{bmatrix}$. (g) Any vector $\vect{v}_3 = (a,b,c)$ with $c \neq 0$.

\begin{example}{3.5 B and 3.5 C}{Change of Basis Matrix $V = WM$ and Invertible $A$}
\textbf{3.5 B:} Start with three independent vectors $\vect{w}_1, \vect{w}_2, \vect{w}_3$. Take combinations $\begin{bmatrix}\vect{v}_1 & \vect{v}_2 & \vect{v}_3\end{bmatrix} = \begin{bmatrix}\vect{w}_1 & \vect{w}_2 & \vect{w}_3\end{bmatrix}\begin{bmatrix}1 & 1 & 0\\1 & 2 & 1\\0 & 1 & c\end{bmatrix} = WM$. If $c\neq 1$ show that $\vect{v}_1, \vect{v}_2, \vect{v}_3$ are independent. If $c=1$ show that they are dependent.\\
\textbf{3.5 C:} Suppose $\vect{v}_1, \dots, \vect{v}_n$ is a basis for $\mathbf{R}^n$ and the $n$ by $n$ matrix $A$ is invertible. Show that $A\vect{v}_1, \dots, A\vect{v}_n$ is also a basis for $\mathbf{R}^n$.
\end{example}
\textbf{Solution:}
\textbf{3.5 B:} If $c=1$, $\vect{v}_1 - \vect{v}_2 + \vect{v}_3 = \vect{0}$ (and $M\vect{x}=\vect{0}$ for $\vect{x}=(1,-1,1)$), so the $\vect{v}$'s are dependent. If $c\neq 1$, $M$ is invertible, so $V\vect{x} = WM\vect{x} = \vect{0} \implies M\vect{x} = \vect{0} \implies \vect{x} = \vect{0}$; thus the $\vect{v}$'s are independent.

\textbf{3.5 C:} If $c_1A\vect{v}_1 + \dots + c_nA\vect{v}_n = A(c_1\vect{v}_1 + \dots + c_n\vect{v}_n) = \vect{0}$, multiply by $A^{-1}$ to get $c_1\vect{v}_1 + \dots + c_n\vect{v}_n = \vect{0}$, which forces all $c_i = 0$ (independence). To solve $c_1A\vect{v}_1 + \dots + c_nA\vect{v}_n = \vect{b}$, solve $c_1\vect{v}_1 + \dots + c_n\vect{v}_n = A^{-1}\vect{b}$, which is solvable because the $\vect{v}$'s span $\mathbf{R}^n$.

\begin{example}{3.6 A and 3.6 B}{Bases of the Four Subspaces from $A = LU = E^{-1}R$}
\textbf{3.6 A:} Find bases and dimensions for all four fundamental subspaces if you know that:
\[
A = \begin{bmatrix} 1 & 0 & 0 \\ 2 & 1 & 0 \\ 5 & 0 & 1 \end{bmatrix} \begin{bmatrix} 1 & 3 & 0 & 5 \\ 0 & 0 & 1 & 6 \\ 0 & 0 & 0 & 0 \end{bmatrix} = LU = E^{-1}R.
\]
\textbf{3.6 B:} Put four $1$'s into a $5$ by $6$ matrix of zeros, keeping the dimension of its row space as small as possible. How to make the dimension of its nullspace as small as possible? What is the sum of the dimensions of all four subspaces?
\end{example}
\textbf{Solution:}
\textbf{3.6 A:} Rank $r=2$ (pivots in columns 1 and 3).
\begin{itemize}
\item \textbf{Row space ($\dim = 2$):} Basis $(1,3,0,5)$ and $(0,0,1,6)$ from the nonzero rows of $R$.
\item \textbf{Column space ($\dim = 2$):} Basis $(1,2,5)$ and $(0,1,0)$ from the first 2 columns of $E^{-1}$ (which are pivot columns 1 and 3 of $A$).
\item \textbf{Nullspace ($\dim = 2$):} Basis $(-3,1,0,0)$ and $(-5,0,-6,1)$ from $R$.
\item \textbf{Left nullspace $N(A^\T)$ ($\dim = 1$):} Row 3 of $E = L^{-1}$ is $(-5,0,1)$, which is a basis for $N(A^\T)$.
\end{itemize}
\textbf{3.6 B:} To make the row space (and column space) dimension $r=1$, put all four $1$'s in the same row, or in the same column, or in a $2\times 2$ block $\begin{bmatrix}1&1\\1&1\end{bmatrix}$. To make the nullspace dimension $6-r$ as small as possible ($6-4=2$), put the four $1$'s in four different rows and four different columns so $r=4$. The sum of all four dimensions is always $r + (n-r) + r + (m-r) = n + m = 6 + 5 = 11$.

% ============================================================================
% MODULE 2: SOLVING LINEAR EQUATIONS AND DETERMINANTS (PDF 4)
% ============================================================================
\part{Solving Linear Equations and Determinants}

\chapter[Linear Equations, Elimination, Inverses, and $A=LU$ (2.1--2.7)]{Linear Equations, Elimination, Inverses, and $A=LU$\bk{2.1--2.7}}

\begin{syll}
Syllabus topics: Row picture and column picture of $A\vect{x}=\vect{b}$ $\cdot$ Gaussian elimination, pivots, multipliers $\ell_{ij}$, back substitution $\cdot$ Elimination matrices $E_{ij}$ and permutation matrices $P_{ij}$ $\cdot$ Matrix multiplication (4 ways) and block multiplication $\cdot$ Inverse matrices $A^{-1}$ and Gauss-Jordan elimination $[A \ \ I] \to [I \ \ A^{-1}]$ $\cdot$ Factorization $A=LU$ and $A=LDU$ $\cdot$ Operation counts ($\frac{1}{3}n^3$ and $n^2$) $\cdot$ Transposes $A^\T$, symmetric matrices $A=LDL^\T$, and $PA=LU$.
\end{syll}

\section{Vectors and Linear Equations: Row and Column Pictures\bk{2.1}}

\begin{goal}
Represent a system of $n$ linear equations in $n$ unknowns in row form, column form, and matrix form $A\vect{x}=\vect{b}$.
\end{goal}

\step{Step 1: Compare the row picture and column picture.}
Consider the system:
\begin{equation}
\begin{aligned}
x - 2y &= 1 \\
3x + 2y &= 11
\end{aligned}
\quad \Longleftrightarrow \quad
x\begin{bmatrix} 1 \\ 3 \end{bmatrix} + y\begin{bmatrix} -2 \\ 2 \end{bmatrix} = \begin{bmatrix} 1 \\ 11 \end{bmatrix}
\quad \Longleftrightarrow \quad
\begin{bmatrix} 1 & -2 \\ 3 & 2 \end{bmatrix}\begin{bmatrix} x \\ y \end{bmatrix} = \begin{bmatrix} 1 \\ 11 \end{bmatrix}.
\end{equation}
\begin{itemize}
\item \term{Row picture:} Each equation defines a line (in $\mathbf{R}^2$) or a plane (in $\mathbf{R}^3$). They meet at the solution $(x,y)=(3,1)$.
\item \term{Column picture:} $A\vect{x} = \vect{b}$ combines the column vectors on the left side: $3(\text{column 1}) + 1(\text{column 2}) = \vect{b}$.
\end{itemize}

\begin{pic}
\begin{tikzpicture}[scale=0.58]
  % Row picture
  \begin{scope}[xshift=0cm]
    \draw[ax] (-1,0) -- (5,0) node[right] {$x$};
    \draw[ax] (0,-1) -- (0,4.5) node[above] {$y$};
    \draw[thick,teal] (-1,-1) -- (5,2) node[right,font=\footnotesize] {$x-2y=1$};
    \draw[thick,amber] (1,4) -- (4,-0.5) node[right,font=\footnotesize] {$3x+2y=11$};
    \fill[rose] (3,1) circle (3.5pt) node[above right,font=\footnotesize\bfseries] {$(3,1)$};
    \node[font=\small\bfseries] at (2,-1.8) {Row picture};
  \end{scope}
  % Column picture
  \begin{scope}[xshift=10cm,yshift=-1cm,scale=0.45]
    \draw[ax] (-3,0) -- (4,0) node[right] {$x$};
    \draw[ax] (0,-1) -- (0,12) node[above] {$y$};
    \draw[vec,teal] (0,0) -- (1,3) node[right,font=\footnotesize] {$\text{col } 1$};
    \draw[vec,teal,dashed] (1,3) -- (3,9) node[right,font=\footnotesize] {$3(\text{col } 1)=\begin{bmatrix}3\\9\end{bmatrix}$};
    \draw[vec,amber] (0,0) -- (-2,2) node[left,font=\footnotesize] {$\text{col } 2=\begin{bmatrix}-2\\2\end{bmatrix}$};
    \draw[vec,amber] (3,9) -- (1,11);
    \draw[vec,rose,very thick] (0,0) -- (1,11) node[above,font=\footnotesize\bfseries] {$\vect{b}=\begin{bmatrix}1\\11\end{bmatrix}$};
    \node[font=\small\bfseries] at (1,-2.2) {Column picture};
  \end{scope}
\end{tikzpicture}
\figcap{Row picture (lines meeting at $(3,1)$) and column picture ($3\vect{a}_1 + 1\vect{a}_2 = \vect{b}$) \bookref{Figs.\ 2.1--2.2}.}
\end{pic}

\begin{example}{2.1.1 (p.\ 37)}{Multiplication by a Singular Matrix $A$ and the Identity Matrix $I$}
Multiply $A\vect{x}$ and $I\vect{x}$ for:
\[
A\vect{x} = \begin{bmatrix} 1 & 0 & 0 \\ 1 & 0 & 0 \\ 1 & 0 & 0 \end{bmatrix} \begin{bmatrix} 4 \\ 5 \\ 6 \end{bmatrix} \quad \text{and} \quad I\vect{x} = \begin{bmatrix} 1 & 0 & 0 \\ 0 & 1 & 0 \\ 0 & 0 & 1 \end{bmatrix} \begin{bmatrix} 4 \\ 5 \\ 6 \end{bmatrix}.
\]
\end{example}

\vfill
\begin{pageend_sol}{2.1.1 (p.\ 37)}
By columns, $A\vect{x}$ is $4$ times the first column $(1,1,1)$ plus $5\vect{0} + 6\vect{0} = (4,4,4)$. The identity matrix $I$ leaves every vector unchanged:
\[
A\vect{x} = \begin{bmatrix} 4 \\ 4 \\ 4 \end{bmatrix} \quad \text{and} \quad I\vect{x} = \begin{bmatrix} 4 \\ 5 \\ 6 \end{bmatrix} = \vect{x}.
\]
\end{pageend_sol}
\clearpage

\section{The Idea of Elimination and Elimination Matrices\bk{2.2--2.3}}

\begin{goal}
Derive the multiplier rule $\ell_{ij} = a_{ij}/d_j$ and express each elimination step and row exchange as a matrix multiplication.
\end{goal}

\step{Step 1: Define the pivot and the multiplier.}
To eliminate the entry in row $i$, column $j$ (where $i > j$) below the nonzero pivot $d_j$ in row $j$, subtract $\ell_{ij}$ times row $j$ from row $i$:
\begin{equation}
\eqbox{\text{Multiplier } \ell_{ij} = \frac{\text{entry to eliminate in row } i}{\text{pivot in row } j}.} \qquad \text{(Pivots can never be zero!)}
\end{equation}
This converts $A\vect{x} = \vect{b}$ into an upper triangular system $U\vect{x} = \vect{c}$, which you solve from bottom to top by \term{back substitution}.

\step{Step 2: Represent elimination steps by elementary matrices $E_{ij}$ and permutations $P_{ij}$.}
\begin{itemize}
\item The \term{elimination matrix} $E_{ij}$ that subtracts $\ell_{ij}$ times row $j$ from row $i$ is the identity matrix $I$ with $-\ell_{ij}$ in the $(i,j)$ position.
\item The \term{row exchange matrix} $P_{ij}$ that exchanges row $i$ and row $j$ is the identity matrix $I$ with rows $i$ and $j$ exchanged.
\end{itemize}

\begin{example}{2.2.1--2.2.3 and 2.3.1--2.3.2 (p.\ 46--58)}{Breakdown of Elimination and Elimination Matrices}
\begin{enumerate}
\item \textbf{Examples 2.2.1--2.2.3:} Apply elimination to these three $2\times 2$ systems and classify the outcome:
\[
\text{(1) } \begin{aligned} x - 2y &= 1 \\ 3x - 6y &= 11 \end{aligned} \qquad
\text{(2) } \begin{aligned} x - 2y &= 1 \\ 3x - 6y &= 3 \end{aligned} \qquad
\text{(3) } \begin{aligned} 0x + 2y &= 4 \\ 3x - 2y &= 5 \end{aligned}
\]
\item \textbf{Example 2.3.1:} The matrix $A$ has $a_{ij} = 2i + j$. Compute $A\vect{x}$ for $\vect{x} = \begin{bmatrix}2\\1\end{bmatrix}$.
\item \textbf{Example 2.3.2:} Write the $3\times 3$ matrix $E_{31}$ with multiplier $\ell = 4$ and compute $E_{31}\begin{bmatrix}1\\3\\9\end{bmatrix}$.
\end{enumerate}
\end{example}

\section{Rules for Matrix Operations\bk{2.4}}

\begin{goal}
Derive the four equivalent ways to multiply an $m\times n$ matrix $A$ by an $n\times p$ matrix $B$.
\end{goal}

\begin{result}[ (Four Ways to Multiply $AB$ and Block Elimination)]
When $A$ is $m \times n$ and $B$ is $n \times p$, their product $AB$ is $m \times p$ and requires $mnp$ multiplications:
\begin{enumerate}
\item \textbf{Rows times columns (inner products):} $(AB)_{ij} = (\text{row } i \text{ of } A) \cdot (\text{column } j \text{ of } B) = \sum_{k=1}^n a_{ik}b_{kj}$.
\item \textbf{Matrix times columns:} $A\begin{bmatrix}\vect{b}_1 & \dots & \vect{b}_p\end{bmatrix} = \begin{bmatrix}A\vect{b}_1 & \dots & A\vect{b}_p\end{bmatrix}$. Each column of $AB$ is a combination of columns of $A$.
\item \textbf{Rows times matrix:} $\text{row } i \text{ of } AB = (\text{row } i \text{ of } A)B$. Each row of $AB$ is a combination of rows of $B$.
\item \textbf{Columns times rows (outer products):} $AB = \sum_{k=1}^n (\text{column } k \text{ of } A)(\text{row } k \text{ of } B)$.
\end{enumerate}
\textbf{Block Elimination and Schur Complement:}
\begin{equation}
\begin{bmatrix} I & \vect{0} \\ -CA^{-1} & I \end{bmatrix} \begin{bmatrix} A & B \\ C & D \end{bmatrix} = \begin{bmatrix} A & B \\ \vect{0} & D - CA^{-1}B \end{bmatrix}, \qquad S = D - CA^{-1}B = \text{Schur complement}.
\end{equation}
\end{result}

\vfill
\begin{pageend_sol}{2.2.1--2.2.3 and 2.3.1--2.3.2 (p.\ 46--58)}
\textbf{1.} \textbf{(1)} Subtract $3\times(\text{eqn 1})$ from eqn 2 to get $0y = 8$: \textbf{permanent failure, no solution} (parallel lines). \textbf{(2)} Subtract $3\times(\text{eqn 1})$ from eqn 2 to get $0y = 0$: \textbf{infinitely many solutions} $x = 1 + 2y$ (same line). \textbf{(3)} Zero in first pivot position: \textbf{temporary failure}; exchange the two equations to get $3x-2y=5$ and $2y=4 \implies y=2, x=3$ (pivots $3$ and $2$).

\textbf{2.} $A = \begin{bmatrix}3 & 4\\5 & 6\end{bmatrix}$, so $A\vect{x} = \begin{bmatrix}3 & 4\\5 & 6\end{bmatrix}\begin{bmatrix}2\\1\end{bmatrix} = \begin{bmatrix}3\cdot 2 + 4\cdot 1\\5\cdot 2 + 6\cdot 1\end{bmatrix} = \begin{bmatrix}10\\16\end{bmatrix}$.

\textbf{3.} $E_{31} = \begin{bmatrix}1 & 0 & 0\\0 & 1 & 0\\-4 & 0 & 1\end{bmatrix}$ and $E_{31}\begin{bmatrix}1\\3\\9\end{bmatrix} = \begin{bmatrix}1\\3\\9 - 4(1)\end{bmatrix} = \begin{bmatrix}1\\3\\5\end{bmatrix}$.
\end{pageend_sol}
\clearpage

\begin{example}{2.4.1--2.4.4 (p.\ 68--71)}{Matrix Multiplication: Inner, Outer, and Block Products}
\begin{enumerate}
\item \textbf{Example 1:} Multiply $\begin{bmatrix}1 & 1\\2 & -1\end{bmatrix}\begin{bmatrix}2 & 2\\3 & 4\end{bmatrix}$.
\item \textbf{Example 2:} Multiply the outer product $\begin{bmatrix}0\\1\\2\end{bmatrix}\begin{bmatrix}1 & 2 & 3\end{bmatrix}$.
\item \textbf{Example 3:} Multiply $\begin{bmatrix}1 & 4\\1 & 5\end{bmatrix}\begin{bmatrix}3 & 2\\1 & 0\end{bmatrix}$ using columns times rows.
\item \textbf{Example 4:} Find one block elimination matrix $E = E_{31}E_{21}$ that clears out $3$ and $4$ below the pivot $1$ in $\begin{bmatrix}1 & x & x\\3 & x & x\\4 & x & x\end{bmatrix}$.
\end{enumerate}
\end{example}

\section{Inverse Matrices and Gauss-Jordan Elimination\bk{2.5}}

\begin{goal}
Define the inverse matrix $A^{-1}$, derive the $2\times 2$ inverse, prove $(AB)^{-1} = B^{-1}A^{-1}$, and compute $A^{-1}$ by Gauss-Jordan elimination.
\end{goal}

\step{Step 1: Define $A^{-1}$ and prove uniqueness (left-inverse $=$ right-inverse).}
An $n \times n$ matrix $A$ is \term{invertible} if there exists a matrix $A^{-1}$ such that:
\begin{equation}
A^{-1}A = I \quad \text{and} \quad AA^{-1} = I.
\end{equation}
If $BA = I$ and $AC = I$, apply the associative law:
\begin{equation}
B(AC) = (BA)C \implies BI = IC \implies B = C.
\end{equation}

\step{Step 2: Derive the $2\times 2$ inverse and the reverse-order product rule.}
Multiply $\begin{bmatrix}a & b\\c & d\end{bmatrix}\begin{bmatrix}d & -b\\-c & a\end{bmatrix} = \begin{bmatrix}ad-bc & 0\\0 & ad-bc\end{bmatrix} = (ad-bc)I$. Divide by $\det A = ad-bc \neq 0$:
\begin{equation}
\begin{bmatrix} a & b \\ c & d \end{bmatrix}^{-1} = \frac{1}{ad - bc}\begin{bmatrix} d & -b \\ -c & a \end{bmatrix}.
\end{equation}
For a product $AB$, multiply $(AB)(B^{-1}A^{-1}) = A(BB^{-1})A^{-1} = AIA^{-1} = AA^{-1} = I$:
\begin{equation}
\eqbox{(AB)^{-1} = B^{-1}A^{-1} \qquad \text{and} \qquad (ABC)^{-1} = C^{-1}B^{-1}A^{-1}.}
\end{equation}

\step{Step 3: Compute $A^{-1}$ by Gauss-Jordan elimination.}
The equation $AA^{-1} = A\begin{bmatrix}\vect{x}_1 & \dots & \vect{x}_n\end{bmatrix} = \begin{bmatrix}\vect{e}_1 & \dots & \vect{e}_n\end{bmatrix} = I$ solves $n$ systems at once. Row-reduce the $n \times 2n$ block matrix $[A \ \ I]$ below and above the pivots, and divide by the pivots:
\begin{equation}
A^{-1}\begin{bmatrix} A & I \end{bmatrix} = \begin{bmatrix} I & A^{-1} \end{bmatrix}.
\end{equation}

\begin{example}{2.5.1--2.5.5 (p.\ 82--87)}{Inverses of Elimination, Tridiagonal, and Triangular Matrices}
\begin{enumerate}
\item \textbf{Example 1:} Show why $A = \begin{bmatrix}1 & 2\\1 & 2\end{bmatrix}$ is not invertible using three tests.
\item \textbf{Examples 2--3:} Find $E^{-1}, F^{-1}, FE$, and $(FE)^{-1} = E^{-1}F^{-1}$ for $E = \begin{bmatrix}1&0&0\\-5&1&0\\0&0&1\end{bmatrix}$ and $F = \begin{bmatrix}1&0&0\\0&1&0\\0&-4&1\end{bmatrix}$.
\item \textbf{Example 4:} Find $A^{-1}$ by Gauss-Jordan elimination starting from $A = \begin{bmatrix}2 & 3\\4 & 7\end{bmatrix}$.
\item \textbf{Example 5:} Use Gauss-Jordan elimination on $[L \ \ I]$ to invert $L = \begin{bmatrix}1 & 0 & 0\\3 & 1 & 0\\4 & 5 & 1\end{bmatrix}$.
\end{enumerate}
\end{example}

\vfill
\begin{pageend_sol}{2.4.1--2.4.4 and 2.5.1--2.5.5}
\textbf{Ex.\ 2.4.1--2.4.4:} \textbf{(1)} $\begin{bmatrix}5 & 6\\1 & 0\end{bmatrix}$. \textbf{(2)} $\begin{bmatrix}0 & 0 & 0\\1 & 2 & 3\\2 & 4 & 6\end{bmatrix}$. \textbf{(3)} $\begin{bmatrix}1\\1\end{bmatrix}\begin{bmatrix}3 & 2\end{bmatrix} + \begin{bmatrix}4\\5\end{bmatrix}\begin{bmatrix}1 & 0\end{bmatrix} = \begin{bmatrix}3 & 2\\3 & 2\end{bmatrix} + \begin{bmatrix}4 & 0\\5 & 0\end{bmatrix} = \begin{bmatrix}7 & 2\\8 & 2\end{bmatrix}$. \textbf{(4)} $E = \begin{bmatrix}1 & 0 & 0\\-3 & 1 & 0\\-4 & 0 & 1\end{bmatrix}$.

\textbf{Ex.\ 2.5.1--2.5.5:} \textbf{(1)} $ad-bc = 2-2=0$; $A\vect{x}=\vect{0}$ for $\vect{x}=(2,-1)$; elimination gives a zero row (only 1 pivot). \textbf{(2--3)} $E^{-1}=\begin{bmatrix}1&0&0\\5&1&0\\0&0&1\end{bmatrix}$, $F^{-1}=\begin{bmatrix}1&0&0\\0&1&0\\0&4&1\end{bmatrix}$, $FE = \begin{bmatrix}1&0&0\\-5&1&0\\20&-4&1\end{bmatrix}$, and $E^{-1}F^{-1} = \begin{bmatrix}1&0&0\\5&1&0\\0&4&1\end{bmatrix}$ (multipliers $5,4$ fall directly into place). \textbf{(4)} $\begin{bmatrix}2 & 3 & 1 & 0\\4 & 7 & 0 & 1\end{bmatrix} \to \begin{bmatrix}2 & 3 & 1 & 0\\0 & 1 & -2 & 1\end{bmatrix} \to \begin{bmatrix}1 & 0 & 7/2 & -3/2\\0 & 1 & -2 & 1\end{bmatrix}$, so $A^{-1} = \begin{bmatrix}7/2 & -3/2\\-2 & 1\end{bmatrix}$. \textbf{(5)} $L^{-1} = \begin{bmatrix}1 & 0 & 0\\-3 & 1 & 0\\11 & -5 & 1\end{bmatrix}$.
\end{pageend_sol}
\clearpage

\section{Elimination $=$ Factorization: $A = LU$ and $A = LDU$\bk{2.6}}

\begin{goal}
Derive the triangular factorization $A = LU$ (and $A = LDU$), explain why the multipliers $\ell_{ij}$ go directly into $L$ without mix-up, and calculate the operation count $\frac{1}{3}n^3$.
\end{goal}

\step{Step 1: Invert the sequence of elimination matrices.}
Without row exchanges, forward elimination multiplies $A$ by $E_{21}, E_{31}, E_{32}, \dots$ to reach $U$:
\begin{equation}
(E_{32}E_{31}E_{21})A = U \implies A = (E_{21}^{-1}E_{31}^{-1}E_{32}^{-1})U = LU.
\end{equation}
Why is there no mix-up of multipliers in $L$? Because row $i$ of $U$ is formed by subtracting $\ell_{ij}$ times \emph{final pivot rows of $U$}:
\begin{equation}
\text{Row 3 of } U = (\text{Row 3 of } A) - \ell_{31}(\text{Row 1 of } U) - \ell_{32}(\text{Row 2 of } U) \implies \text{Row 3 of } A = \begin{bmatrix}\ell_{31} & \ell_{32} & 1\end{bmatrix} U.
\end{equation}
Divide each row of $U$ by its pivot $d_i$ to get the balanced factorization $A = LDU$, where both $L$ and $U$ have $1$'s on the diagonal and $D = \operatorname{diag}(d_1, \dots, d_n)$.

\step{Step 2: Solve $A\vect{x} = \vect{b}$ in two triangular steps and count operations.}
\begin{align}
\text{\textbf{1. Factor } $A = LU$:} \quad &\sum_{k=1}^n (k^2 - k) \approx \int_0^n x^2\dd x = \frac{1}{3}n^3 \text{ multiplications and } \frac{1}{3}n^3 \text{ subtractions.} \\
\text{\textbf{2. Solve } $L\vect{c} = \vect{b}$ then $U\vect{x} = \vect{c}$:} \quad &[(n-1) + \dots + 1] + [1 + \dots + n] = n^2 \text{ multiplications and } n^2 \text{ subtractions.}
\end{align}
For a band matrix with half-bandwidth $w$, Factor takes $nw^2$ operations and Solve takes $2nw$ operations.

\begin{example}{2.6.1--2.6.3 (p.\ 96--98)}{Computing $A = LU$ and Solving $L\vect{c}=\vect{b}, U\vect{x}=\vect{c}$}
\begin{enumerate}
\item \textbf{Example 1:} Factor $A = \begin{bmatrix}2 & 1 & 0\\1 & 2 & 1\\0 & 1 & 2\end{bmatrix}$ into $LU$.
\item \textbf{Example 2:} Factor the $4\times 4$ matrix $A = \begin{bmatrix}1&1&0&0\\1&2&1&0\\0&1&2&1\\0&0&1&2\end{bmatrix}$ into $LU$.
\item \textbf{Example 3:} Solve $u + 2v = 5$, $4u + 9v = 21$ via $L\vect{c} = \vect{b}$ and $U\vect{x} = \vect{c}$.
\end{enumerate}
\end{example}

\section{Transposes, Symmetric Matrices, and Permutations\bk{2.7}}

\begin{goal}
Derive the rules for transposes $(AB)^\T = B^\T A^\T$ and $(A^{-1})^\T = (A^\T)^{-1}$, the inner product identity $(A\vect{x})^\T\vect{y} = \vect{x}^\T(A^\T\vect{y})$, the symmetric factorization $A = LDL^\T$, and $PA = LU$.
\end{goal}

\begin{result}[ (Transpose, Inner Product, Symmetric, and Permutation Identities)]
\begin{align}
(A^\T)_{ij} &= A_{ji}, \qquad (A+B)^\T = A^\T + B^\T, \qquad (AB)^\T = B^\T A^\T, \qquad (A^{-1})^\T = (A^\T)^{-1} \\
(A\vect{x})^\T\vect{y} &= \vect{x}^\T A^\T\vect{y} = \vect{x}^\T(A^\T\vect{y}) &&\why{defining property of $A^\T$} \\
(R^\T R)^\T &= R^\T(R^\T)^\T = R^\T R &&\why{$R^\T R$ and $RR^\T$ are always symmetric} \\
A = A^\T &\implies A = LDU = LDL^\T &&\why{symmetric factorization ($U = L^\T$)} \\
P^\T P &= I \implies P^{-1} = P^\T, \qquad PA = LU &&\why{for any permutation matrix $P$ and invertible $A$}
\end{align}
\end{result}

\begin{example}{2.7.1--2.7.5 (p.\ 108--111)}{Transposes, Inner Products, and $R^\T R$}
\begin{enumerate}
\item \textbf{Example 1:} Find $A^{-1}$, $A^\T$, $(A^{-1})^\T$, and $(A^\T)^{-1}$ for $A = \begin{bmatrix}1 & 0\\6 & 1\end{bmatrix}$.
\item \textbf{Example 2:} Use $(A\vect{x})^\T\vect{y} = \vect{x}^\T(A^\T\vect{y})$ to find $A^\T\vect{y}$ for $A = \begin{bmatrix}-1 & 1 & 0\\0 & -1 & 1\end{bmatrix}$.
\item \textbf{Example 4:} Multiply $RR^\T$ and $R^\T R$ for $R = \begin{bmatrix}-1 & 1 & 0\\0 & -1 & 1\end{bmatrix}$.
\end{enumerate}
\end{example}

\vfill
\begin{pageend_sol}{2.6.1--2.6.3 and 2.7.1--2.7.5}
\textbf{Ex.\ 2.6.1--2.6.3:} \textbf{(1)} $\ell_{21}=\frac{1}{2}, \ell_{31}=0, \ell_{32}=\frac{2}{3}$, so $A = \begin{bmatrix}1&0&0\\1/2&1&0\\0&2/3&1\end{bmatrix}\begin{bmatrix}2&1&0\\0&3/2&1\\0&0&4/3\end{bmatrix}$. \textbf{(2)} All pivots and multipliers are $1$: $L$ has $1$'s on the main diagonal and subdiagonal; $U$ has $1$'s on the main diagonal and superdiagonal. \textbf{(3)} $L\vect{c} = \begin{bmatrix}1&0\\4&1\end{bmatrix}\vect{c} = \begin{bmatrix}5\\21\end{bmatrix} \implies \vect{c} = \begin{bmatrix}5\\1\end{bmatrix}$; then $U\vect{x} = \begin{bmatrix}1&2\\0&1\end{bmatrix}\vect{x} = \begin{bmatrix}5\\1\end{bmatrix} \implies \vect{x} = \begin{bmatrix}3\\1\end{bmatrix}$.

\textbf{Ex.\ 2.7.1--2.7.5:} \textbf{(1)} $A^{-1} = \begin{bmatrix}1 & 0\\-6 & 1\end{bmatrix}$, $A^\T = \begin{bmatrix}1 & 6\\0 & 1\end{bmatrix}$, and $(A^{-1})^\T = (A^\T)^{-1} = \begin{bmatrix}1 & -6\\0 & 1\end{bmatrix}$. \textbf{(2)} $(A\vect{x})^\T\vect{y} = (x_2-x_1)y_1 + (x_3-x_2)y_2 = x_1(-y_1) + x_2(y_1-y_2) + x_3(y_2)$, so $A^\T\vect{y} = \begin{bmatrix}-y_1\\y_1-y_2\\y_2\end{bmatrix}$. \textbf{(4)} $RR^\T = \begin{bmatrix}2 & -1\\-1 & 2\end{bmatrix}$ and $R^\T R = \begin{bmatrix}1 & -1 & 0\\-1 & 2 & -1\\0 & -1 & 1\end{bmatrix}$.
\end{pageend_sol}
\clearpage

% ============================================================================
% CHAPTER 4: DETERMINANTS (5.1 - 5.3)
% ============================================================================
\chapter[Determinants: Properties, Cofactors, Cramer's Rule (5.1--5.3)]{Determinants: Properties, Cofactors, Cramer's Rule, and Volumes\bk{5.1--5.3}}

\begin{syll}
Syllabus topics: Rules 1--3 defining the determinant $\cdot$ Derived Rules 4--10 ($\det AB = (\det A)(\det B)$, $\det A^\T = \det A$) $\cdot$ Pivot formula $\det A = \pm(d_1d_2\cdots d_n)$ and $d_k = \det A_k / \det A_{k-1}$ $\cdot$ Big formula ($n!$ permutation terms) $\cdot$ Cofactor expansion $\det A = a_{i1}C_{i1} + \dots + a_{in}C_{in}$ $\cdot$ Cramer's Rule $x_j = \det B_j / \det A$ $\cdot$ Cofactor inverse $A^{-1} = C^\T / \det A$ $\cdot$ Area, volume, Jacobian $J$, and cross product $\vect{u}\times\vect{v}$.
\end{syll}

\section{The Ten Properties of Determinants\bk{5.1}}

\begin{goal}
State the three defining properties of $\det A$ and derive properties 4 through 10 from them.
\end{goal}

\step{Defining Rules 1, 2, 3:}
\begin{enumerate}
\item \textbf{Identity matrix:} $\det I = 1$.
\item \textbf{Row exchange (sign reversal):} $\det A$ changes sign when two rows are exchanged ($\det P = \pm 1$).
\item \textbf{Linearity in each row separately:}
\begin{equation}
\begin{vmatrix} ta & tb \\ c & d \end{vmatrix} = t\begin{vmatrix} a & b \\ c & d \end{vmatrix} \quad \text{and} \quad \begin{vmatrix} a+a' & b+b' \\ c & d \end{vmatrix} = \begin{vmatrix} a & b \\ c & d \end{vmatrix} + \begin{vmatrix} a' & b' \\ c & d \end{vmatrix}.
\end{equation}
\end{enumerate}

\step{Derived Rules 4 through 10:}
\begin{enumerate}[start=4]
\item \textbf{Two equal rows $\implies \det A = 0$:} Exchange the two equal rows. By Rule 2, $D = -D \implies D = 0$.
\item \textbf{Row elimination leaves $\det A$ unchanged:} Subtracting $\ell$ times row $j$ from row $i$ splits by Rule 3 into $\det A - \ell(\text{det with two equal rows } j) = \det A - 0 = \det A$.
\item \textbf{Row of zeros $\implies \det A = 0$:} Add another row to the zero row (Rule 5) to get two equal rows (Rule 4).
\item \textbf{Triangular matrix $A \implies \det A = a_{11}a_{22}\cdots a_{nn}$:} Eliminate off-diagonal entries (Rule 5), factor out $a_{11}, \dots, a_{nn}$ from each row (Rule 3), and use $\det I = 1$ (Rule 1).
\item \textbf{$\det A = 0$ exactly when $A$ is singular:} If $A$ is singular, $U$ has a zero row so $\det A = \det U = 0$. If $A$ is invertible, $\det A = \pm \det U = \pm(d_1 d_2 \cdots d_n) \neq 0$.
\item \textbf{Product rule $\det(AB) = (\det A)(\det B)$:} For $|B| \neq 0$, the ratio $D(A) = |AB|/|B|$ satisfies Rules 1, 2, 3 for $A$ ($D(I) = |B|/|B| = 1$, sign reversal, linearity in row 1 of $A$), so $D(A) = |A|$. Also $\det(A^{-1}) = 1/\det A$.
\item \textbf{Transpose rule $\det A^\T = \det A$:} From $PA = LU$, we have $A^\T P^\T = U^\T L^\T$. Since $\det L = \det L^\T = 1$, $\det U = \det U^\T$, and $\det P = \det P^\T = \pm 1$, we get $\det A^\T = \det A$.
\end{enumerate}

\section{Permutations, the Big Formula, and Cofactors\bk{5.2}}

\begin{goal}
Derive the pivot ratio formula, the $n!$ Big Formula, and the cofactor expansion for $\det A$.
\end{goal}

\begin{result}[ (Three Formulas for the Determinant)]
\textbf{1. Pivot Formula (and $k$th Pivot Ratio without Row Exchanges):}
\begin{equation}
\det A = (\det P)(d_1 d_2 \cdots d_n) = \pm(d_1 d_2 \cdots d_n), \qquad \text{and} \qquad d_k = \frac{d_1 d_2 \cdots d_k}{d_1 d_2 \cdots d_{k-1}} = \frac{\det A_k}{\det A_{k-1}}.
\end{equation}
\textbf{2. Big Formula (Sum over all $n!$ Column Permutations $P = (\alpha, \beta, \dots, \omega)$):}
\begin{equation}
\det A = \sum_{\text{all } n!\text{ permutations } P} (\det P)\, a_{1\alpha}a_{2\beta}\cdots a_{n\omega}.
\end{equation}
\textbf{3. Cofactor Formula (along Row $i$ or Column $j$):}
\begin{equation}
\det A = a_{i1}C_{i1} + a_{i2}C_{i2} + \dots + a_{in}C_{in}, \qquad \text{where } C_{ij} = (-1)^{i+j}\det M_{ij}.
\end{equation}
Here $M_{ij}$ is the $(n-1)\times(n-1)$ submatrix formed by deleting row $i$ and column $j$.
\end{result}

\begin{example}{5.2.1--5.2.7 (p.\ 255--261)}{Pivots, Big Formula, and Tridiagonal Cofactors}
\begin{enumerate}
\item \textbf{Example 1:} Find $\det A$ for $A = \begin{bmatrix}0&0&1\\0&2&3\\4&5&6\end{bmatrix}$.
\item \textbf{Examples 2, 6, 7:} Find the determinant $D_n$ of the $n\times n$ tridiagonal $-1,2,-1$ matrix $A_n$ by pivots and by cofactors ($D_n = 2D_{n-1} - D_{n-2}$). What is $\det B_4$ if the top-left entry $a_{11}=2$ changes to $b_{11}=1$?
\item \textbf{Example 4:} Find the determinant of $Z = \begin{bmatrix}1&0&a&0\\0&1&b&0\\0&0&c&0\\0&0&d&1\end{bmatrix}$.
\item \textbf{Example 5:} Find $\det A_4$ for $A_4 = \begin{bmatrix}0&1&0&0\\1&0&1&0\\0&1&0&1\\0&0&1&0\end{bmatrix}$.
\end{enumerate}
\end{example}

\vfill
\begin{pageend_sol}{5.2.1--5.2.7 (p.\ 255--261)}
\textbf{1.} Exchange rows 1 and 3 ($\det P = -1$) to get $PA = \begin{bmatrix}4&5&6\\0&2&3\\0&0&1\end{bmatrix}$. Then $\det A = -(4)(2)(1) = -8$.

\textbf{2.} By pivots: $d_k = \frac{k+1}{k}$, so $\det A_n = (2)\left(\frac{3}{2}\right)\left(\frac{4}{3}\right)\cdots\left(\frac{n+1}{n}\right) = n + 1$. By cofactors of row 1: $D_n = 2D_{n-1} - D_{n-2}$ with $D_1=2, D_2=3 \implies D_n = n+1$. For $B_4$ ($b_{11}=1$), $\det B_4 = 1\cdot D_3 - D_2 = 4 - 3 = 1$ (and $\det B_n = 1$ for all $n$).

\textbf{3.} $\det Z = (1)(1)(c)(1) = c$ (only the main diagonal term is nonzero).

\textbf{4.} Only the permutation $(\alpha,\beta,\gamma,\delta) = (2,1,4,3)$ picks nonzeros from $A_4$. It uses two row exchanges ($\det P_4 = +1$), so $\det A_4 = +1$.
\end{pageend_sol}
\clearpage

\section{Cramer's Rule, Inverses, Volumes, and Cross Product\bk{5.3}}

\begin{goal}
Derive Cramer's Rule for $A\vect{x}=\vect{b}$, the cofactor formula for $A^{-1}$, the determinant formulas for triangle area and box volume, and the 3D cross product $\vect{u}\times\vect{v}$.
\end{goal}

\step{Step 1: Derive Cramer's Rule from the product rule for determinants.}
To find $x_1$ in $A\vect{x} = \vect{b}$, replace column 1 of $I$ by $\vect{x}$ and multiply by $A$:
\begin{equation}
A\begin{bmatrix} x_1 & 0 & 0 \\ x_2 & 1 & 0 \\ x_3 & 0 & 1 \end{bmatrix} = \begin{bmatrix} \vect{b} & \vect{a}_2 & \vect{a}_3 \end{bmatrix} = B_1.
\end{equation}
Take determinants of both sides using $\det(\text{second matrix}) = x_1$:
\begin{equation}
(\det A)(x_1) = \det B_1 \implies x_1 = \frac{\det B_1}{\det A}.
\end{equation}
In general, replacing column $j$ of $A$ by $\vect{b}$ gives $B_j$:
\begin{result}[ (Cramer's Rule and Cofactor Formula for $A^{-1}$)]
\begin{equation}
x_j = \frac{\det B_j}{\det A}, \qquad \text{and for } AA^{-1}=I: \qquad (A^{-1})_{ij} = \frac{C_{ji}}{\det A} \quad \text{or} \quad A^{-1} = \frac{C^\T}{\det A}.
\end{equation}
Proof of $A C^\T = (\det A)I$: Row $i$ of $A$ times column $i$ of $C^\T$ is $a_{i1}C_{i1}+\dots+a_{in}C_{in} = \det A$. Row $i$ of $A$ times column $k \neq i$ of $C^\T$ is the determinant of a matrix with row $i$ equal to row $k$, which is $0$.
\end{result}

\step{Step 2: Area of a triangle, volume of a box, and cross product.}
A parallelogram with edges $(x_1,y_1)$ and $(x_2,y_2)$ from $(0,0)$ satisfies Rules 1, 2, 3 of the determinant, so its area is $|x_1y_2 - x_2y_1|$. A triangle is half of a parallelogram:
\begin{align}
\text{Area of triangle with corners } (x_1,y_1), (x_2,y_2), (x_3,y_3) &= \frac{1}{2}\left|\det\begin{bmatrix} x_1 & y_1 & 1 \\ x_2 & y_2 & 1 \\ x_3 & y_3 & 1 \end{bmatrix}\right| \\
\text{Volume of box with edges } \vect{u}, \vect{v}, \vect{w} &= |\det A| = |(\vect{u}\times\vect{v})\cdot\vect{w}| \\
\text{Cross product in } \mathbf{R}^3: \quad \vect{u}\times\vect{v} = \begin{vmatrix} \ihat & \jhat & \khat \\ u_1 & u_2 & u_3 \\ v_1 & v_2 & v_3 \end{vmatrix} &= (u_2v_3-u_3v_2)\ihat + (u_3v_1-u_1v_3)\jhat + (u_1v_2-u_2v_1)\khat.
\end{align}
The length is $\|\vect{u}\times\vect{v}\| = \|\vect{u}\|\|\vect{v}\||\sin\theta| = \text{area of the } (\vect{u},\vect{v})\text{ parallelogram}$, and $\vect{u}\cdot(\vect{u}\times\vect{v}) = 0$.

\begin{example}{5.3.1--5.3.8 (p.\ 269--276)}{Cramer's Rule, Cofactor Inverse, Jacobian, and Cross Product}
\begin{enumerate}
\item \textbf{Example 1:} Solve $3x_1 + 4x_2 = 2$ and $5x_1 + 6x_2 = 4$ by Cramer's Rule.
\item \textbf{Example 3:} Invert the sum matrix $A = \begin{bmatrix}1&0&0&0\\1&1&0&0\\1&1&1&0\\1&1&1&1\end{bmatrix}$ using $A^{-1} = C^\T / \det A$.
\item \textbf{Example 6:} Find the polar coordinate Jacobian $J$ for $x = r\cos\theta, y = r\sin\theta$.
\item \textbf{Examples 7--8:} Compute $\vect{u}\times\vect{v}$ for (7) $\vect{u}=(3,2,0), \vect{v}=(1,4,0)$ and (8) $\vect{u}=(1,1,1), \vect{v}=(1,1,2)$.
\end{enumerate}
\end{example}

\vfill
\begin{pageend_sol}{5.3.1--5.3.8 (p.\ 269--276)}
\textbf{1.} $\det A = 18-20 = -2$, $\det B_1 = \begin{vmatrix}2&4\\4&6\end{vmatrix} = -4$, $\det B_2 = \begin{vmatrix}3&2\\5&4\end{vmatrix} = 2$. Thus $x_1 = \frac{-4}{-2} = 2$ and $x_2 = \frac{2}{-2} = -1$.

\textbf{2.} $\det A = 1$. Cofactors $C_{ii} = 1$ on the diagonal and $C_{i,i+1} = -1$ on the superdiagonal of $C$. Transposing $C$ gives $A^{-1} = C^\T$ with $1$'s on the diagonal and $-1$'s on the subdiagonal.

\textbf{3.} $J = \begin{vmatrix}\partial x/\partial r & \partial x/\partial\theta\\\partial y/\partial r & \partial y/\partial\theta\end{vmatrix} = \begin{vmatrix}\cos\theta & -r\sin\theta\\\sin\theta & r\cos\theta\end{vmatrix} = r\cos^2\theta + r\sin^2\theta = r$.

\textbf{4.} \textbf{(7)} $\vect{u}\times\vect{v} = (3\cdot 4 - 2\cdot 1)\khat = 10\khat = (0,0,10)$, area $= 10$. \textbf{(8)} $\vect{u}\times\vect{v} = \ihat(2-1) - \jhat(2-1) + \khat(1-1) = (1,-1,0)$, area $= \sqrt{2}$.
\end{pageend_sol}
\clearpage

% ============================================================================
% END OF MODULE 2: ALL UNIQUE SOLVED PROBLEMS (WORKED EXAMPLES)
% ============================================================================
\chapter[Module 2: Unique Solved Problems]{Module 2: All Unique Solved Problems (Worked Examples)}

\begin{syll}
All unique Worked Examples from Module 2 (Sections 2.1--2.7 and 5.1--5.3), with problem statements copied verbatim and full step-by-step solutions.
\end{syll}

\begin{example}{2.1 A}{Solving by Inspection of the Column Picture}
Describe the column picture of these three equations $A\vect{x} = \vect{b}$. Solve by careful inspection of the columns (instead of elimination):
\[
\begin{aligned}
x + 3y + 2z &= -3 \\
2x + 2y + 2z &= -2 \\
3x + 5y + 6z &= -5
\end{aligned}
\quad \text{which is} \quad
\begin{bmatrix} 1 & 3 & 2 \\ 2 & 2 & 2 \\ 3 & 5 & 6 \end{bmatrix} \begin{bmatrix} x \\ y \\ z \end{bmatrix} = \begin{bmatrix} -3 \\ -2 \\ -5 \end{bmatrix}.
\]
\end{example}
\textbf{Solution:} In the column picture, $\vect{b} = (-3,-2,-5)$ is minus the second column $(3,2,5)$. Therefore, the solution is $(x,y,z) = (0,-1,0)$. If the right side changes to $\vect{b} = (4,4,8) =$ sum of the first two columns, the solution becomes $\vect{x} = (1,1,0)$.

\begin{example}{2.1 B}{Showing a System Has No Solution}
This system has \emph{no solution}. The planes in the row picture don't meet at a point. No combination of the three columns produces $\vect{b}$. How to show this?
\[
\begin{aligned}
x + 3y + 5z &= 4 \\
x + 2y - 3z &= 5 \\
2x + 5y + 2z &= 8
\end{aligned}
\qquad
\begin{bmatrix} 1 & 3 & 5 \\ 1 & 2 & -3 \\ 2 & 5 & 2 \end{bmatrix} \begin{bmatrix} x \\ y \\ z \end{bmatrix} = \begin{bmatrix} 4 \\ 5 \\ 8 \end{bmatrix} = \vect{b}
\]
(1) Multiply the equations by $1, 1, -1$ and add to get $0 = 1$. Are any two of the planes parallel? What are the equations of planes parallel to $x + 3y + 5z = 4$?\\
(2) Take the dot product of each column of $A$ (and also $\vect{b}$) with $\vect{y} = (1, 1, -1)$. How do those dot products show that $A\vect{x} = \vect{b}$ has no solution?\\
(3) Find three right side vectors $\vect{b}^*$ and $\vect{b}^{**}$ and $\vect{b}^{***}$ that \emph{do} allow solutions.
\end{example}
\textbf{Solution:}
\begin{enumerate}
\item[(1)] $(\text{eqn 1}) + (\text{eqn 2}) - (\text{eqn 3})$ gives $0x + 0y + 0z = 4 + 5 - 8 = 1$, which is impossible ($0=1$). No two planes are parallel. Planes parallel to $x+3y+5z=4$ have the form $x+3y+5z=c$ (or $2x+6y+10z=2c$).
\item[(2)] The dot product of each column of $A$ with $\vect{y}=(1,1,-1)$ is zero, so any combination $A\vect{x}$ must also have zero dot product with $\vect{y}$. However, $\vect{y}\cdot\vect{b} = 1(4)+1(5)-1(8) = 1 \neq 0$, so no solution exists.
\item[(3)] Choose $\vect{b}^* = \begin{bmatrix}1\\1\\2\end{bmatrix}$ ($\text{col 1}, \vect{x}^*=(1,0,0)$), $\vect{b}^{**} = \begin{bmatrix}9\\0\\9\end{bmatrix}$ ($\text{sum of cols}, \vect{x}^{**}=(1,1,1)$), and $\vect{b}^{***} = \begin{bmatrix}0\\0\\0\end{bmatrix}$ ($\vect{x}^{***}=(0,0,0)$).
\end{enumerate}

\begin{example}{2.2 A and 2.2 B}{Pivots and Multipliers of Band and Triangular Matrices}
\textbf{2.2 A:} When elimination is applied to $A = \begin{bmatrix}1 & -1 & 0\\-1 & 2 & -1\\0 & -1 & 2\end{bmatrix}$, what are the first and second pivots? What is $\ell_{21}$? What entry in the $(2,2)$ position would force a row exchange? Why is $\ell_{31}=0$? If $a_{33}=2$ changes to $a_{33}=1$, why does elimination fail?\\
\textbf{2.2 B:} Suppose $A$ is already a triangular matrix. Where do you see its pivots? When does $A\vect{x}=\vect{b}$ have a unique solution?
\end{example}
\textbf{Solution:}
\textbf{2.2 A:} The first pivot is $1$, the multiplier is $\ell_{21} = -1/1 = -1$, and the second pivot is $1$. Changing $a_{22}=2$ to $1$ would produce $0$ in the second pivot position and force a row exchange. $\ell_{31}=0$ because $a_{31}=0$. If $a_{33}=1$, the third pivot becomes $0$ and elimination fails.
\textbf{2.2 B:} The pivots of a triangular matrix sit directly on its main diagonal. $A\vect{x}=\vect{b}$ has a unique solution for every $\vect{b}$ when all diagonal entries are nonzero.

\begin{example}{2.2 C}{Success vs.\ Failure in $3\times 3$ Elimination}
Use elimination to reach upper triangular matrices $U$. Solve by back substitution or explain why this is impossible. What are the pivots?
\[
\textbf{Success:} \quad \begin{aligned} x+y+z &= 7 \\ x+y-z &= 5 \\ x-y+z &= 3 \end{aligned} \qquad\qquad
\textbf{Failure:} \quad \begin{aligned} x+y+z &= 7 \\ x+y-z &= 5 \\ -x-y+z &= 3 \end{aligned}
\]
\end{example}
\textbf{Solution:}
\begin{itemize}
\item \textbf{Success:} Subtract eqn 1 from eqns 2 and 3 to get $0y - 2z = -2$ and $-2y + 0z = -4$. Exchange equations 2 and 3 to get pivots $1, -2, -2$. Back substitution gives $z=1, y=2, x=4$.
\item \textbf{Failure:} Subtract eqn 1 from eqn 2 and add eqn 1 to eqn 3 to get $0y - 2z = -2$ and $0y + 2z = 10$. Adding those gives $0z = 8$: no pivot in column 2, no solution. (If the third right side were $-5$ instead of $3$, we would get $0=0$ and infinitely many solutions $z=1, x=6-y$.)
\end{itemize}

\begin{example}{2.3 A, 2.3 B, 2.3 C}{Elimination Matrices, Augmented Matrices, and Column-Row Products}
\textbf{2.3 A:} What $3\times 3$ matrix $E_{21}$ subtracts $4\times(\text{row 1})$ from row 2? What $P_{32}$ exchanges rows 2 and 3? What do $AE_{21}$ and $AP_{32}$ do?\\
\textbf{2.3 B:} Write $[A \ \ \vect{b}]$ for $x+2y+2z=1, 4x+8y+9z=3, 3y+2z=1$. Apply $E_{21}$ and $P_{32}$, solve, and find $P_{32}E_{21}$.\\
\textbf{2.3 C:} Multiply $AB = \begin{bmatrix}3 & 4\\1 & 5\\2 & 0\end{bmatrix}\begin{bmatrix}2 & 4\\1 & 1\end{bmatrix}$ by rows times columns and by columns times rows.
\end{example}
\textbf{Solution:}
\textbf{2.3 A:} $E_{21} = \begin{bmatrix}1&0&0\\-4&1&0\\0&0&1\end{bmatrix}$ and $P_{32} = \begin{bmatrix}1&0&0\\0&0&1\\0&1&0\end{bmatrix}$. On the right, $AE_{21}$ subtracts $4\times(\text{col 2})$ from col 1; $AP_{32}$ exchanges cols 2 and 3.
\textbf{2.3 B:} $[A \ \ \vect{b}] = \begin{bmatrix}1&2&2&1\\4&8&9&3\\0&3&2&1\end{bmatrix} \xrightarrow{E_{21}} \begin{bmatrix}1&2&2&1\\0&0&1&-1\\0&3&2&1\end{bmatrix} \xrightarrow{P_{32}} \begin{bmatrix}1&2&2&1\\0&3&2&1\\0&0&1&-1\end{bmatrix}$, giving $(x,y,z) = (1,1,-1)$ and $P_{32}E_{21} = \begin{bmatrix}1&0&0\\0&0&1\\-4&1&0\end{bmatrix}$.
\textbf{2.3 C:} $\begin{bmatrix}3\\1\\2\end{bmatrix}\begin{bmatrix}2 & 4\end{bmatrix} + \begin{bmatrix}4\\5\\0\end{bmatrix}\begin{bmatrix}1 & 1\end{bmatrix} = \begin{bmatrix}6 & 12\\2 & 4\\4 & 8\end{bmatrix} + \begin{bmatrix}4 & 4\\5 & 5\\0 & 0\end{bmatrix} = \begin{bmatrix}10 & 16\\7 & 9\\4 & 8\end{bmatrix}$ ($12$ multiplications).

\begin{example}{2.4 A, 2.4 B, 2.4 C}{Pascal Hypercube Matrix $H=L^2$, Commuting Matrices, and Graph Paths}
\textbf{2.4 A:} For the $4\times 4$ lower triangular Pascal matrix $L$ ($\ell_{ij} = \binom{i}{j}$), find the hypercube matrix $H = L^2$ and multiply $H$ times ones and powers $(1, x, x^2, x^3)$.\\
\textbf{2.4 B:} For $A = \begin{bmatrix}p&0\\q&r\end{bmatrix}, B = \begin{bmatrix}1&1\\0&1\end{bmatrix}, C = \begin{bmatrix}0&z\\0&0\end{bmatrix}$, when does $AB = BA$ and $BC = CB$?\\
\textbf{2.4 C:} For the graph adjacency matrix $A = \begin{bmatrix}1&1\\1&0\end{bmatrix}$, compute $A^2$ and $A^3$ and explain how $(A^k)_{ij}$ counts $k$-step paths.
\end{example}
\textbf{Solution:}
\textbf{2.4 A:} $H = L^2 = \begin{bmatrix}1&0&0&0\\2&1&0&0\\4&4&1&0\\8&12&6&1\end{bmatrix}$. Multiplying $H$ by $(1,1,1,1)$ gives $(1,3,9,27)$; multiplying by $(1,x,x^2,x^3)$ gives $(1, 2+x, (2+x)^2, (2+x)^3)$. Row 4 of $H$ $(8,12,6,1)$ counts corners, edges, faces, and box of a 3D cube; the next row $(16,32,24,8,1)$ counts those of a 4D hypercube.
\textbf{2.4 B:} $AB = \begin{bmatrix}p & p\\q & q+r\end{bmatrix}$ and $BA = \begin{bmatrix}p+q & r\\q & r\end{bmatrix}$, so $AB=BA \iff q=0$ and $p=r$. $BC = CB = \begin{bmatrix}0&z\\0&0\end{bmatrix}$ for all $z$.
\textbf{2.4 C:} $A^2 = \begin{bmatrix}2&1\\1&1\end{bmatrix}$ and $A^3 = \begin{bmatrix}3&2\\2&1\end{bmatrix} = \begin{bmatrix}F_4 & F_3\\F_3 & F_2\end{bmatrix}$. Each term $a_{ik}a_{kj}$ is $1$ if and only if a path goes $i \to k \to j$.

\begin{example}{2.5 A, 2.5 B, 2.5 C}{Inverting Difference, Pascal, and $3\times 3$ Matrices}
\textbf{2.5 A:} Invert the difference matrix $A = \begin{bmatrix}1&0&0\\-1&1&0\\0&-1&1\end{bmatrix}$ by Gauss-Jordan.\\
\textbf{2.5 B:} Test invertibility and find inverses when they exist for $A=\begin{bmatrix}4&3\\8&6\end{bmatrix}, B=\begin{bmatrix}4&3\\8&7\end{bmatrix}, C=\begin{bmatrix}6&6\\6&0\end{bmatrix}, D=\begin{bmatrix}6&6\\6&6\end{bmatrix}, S=\begin{bmatrix}1&0&0\\1&1&0\\1&1&1\end{bmatrix}, E=\begin{bmatrix}1&1&1\\1&1&0\\1&1&1\end{bmatrix}$.\\
\textbf{2.5 C:} Invert the $4\times 4$ lower triangular Pascal matrix $L$ by Gauss-Jordan.
\end{example}
\textbf{Solution:}
\textbf{2.5 A:} $A^{-1} = S = \begin{bmatrix}1&0&0\\1&1&0\\1&1&1\end{bmatrix}$ (the lower triangular sum matrix).
\textbf{2.5 B:} $A, D, E$ are singular ($\det A = 0$; $D$ has 1 pivot; $E\vect{x}=\vect{0}$ for $\vect{x}=(-1,1,0)$). The inverses of $B, C, S$ are:
\[
B^{-1} = \frac{1}{4}\begin{bmatrix}7 & -3\\-8 & 4\end{bmatrix}, \qquad C^{-1} = \frac{1}{36}\begin{bmatrix}0 & 6\\6 & -6\end{bmatrix}, \qquad S^{-1} = \begin{bmatrix}1 & 0 & 0\\-1 & 1 & 0\\0 & -1 & 1\end{bmatrix}.
\]
\textbf{2.5 C:} $L^{-1} = \begin{bmatrix}1&0&0&0\\-1&1&0&0\\1&-2&1&0\\-1&3&-3&1\end{bmatrix}$ (Pascal's triangle with alternating signs on odd diagonals).

\begin{example}{2.6 A and 2.6 B}{LU Factorization of the Symmetric Pascal Matrix}
\textbf{2.6 A:} Establish $P = LU$ for the $4\times 4$ symmetric Pascal matrix $P = \operatorname{pascal}(4) = \begin{bmatrix}1&1&1&1\\1&2&3&4\\1&3&6&10\\1&4&10&20\end{bmatrix}$.\\
\textbf{2.6 B:} Solve $P\vect{x} = \vect{b} = (1,0,0,0)$ using $L\vect{c} = \vect{b}$ and $U\vect{x} = \vect{c}$.
\end{example}
\textbf{Solution:}
\textbf{2.6 A:} Elimination on $P$ uses multipliers from Pascal's triangle and leaves $U = L^\T$ (all pivots $= 1$):
\[
P = LU = \begin{bmatrix}1&0&0&0\\1&1&0&0\\1&2&1&0\\1&3&3&1\end{bmatrix}\begin{bmatrix}1&1&1&1\\0&1&2&3\\0&0&1&3\\0&0&0&1\end{bmatrix} = LL^\T.
\]
\textbf{2.6 B:} Forward substitution on $L\vect{c} = (1,0,0,0)$ gives $\vect{c} = (1,-1,1,-1)$. Back substitution on $U\vect{x} = \vect{c}$ gives $\vect{x} = (4,-6,4,-1)$.

\begin{example}{2.7 A, 2.7 B, 2.7 C}{Symmetry Under Permutations, $LDL^\T$, and Saddle-Point Matrices}
\textbf{2.7 A and 2.7 B:} For $P = \begin{bmatrix}0&1&0\\0&0&1\\1&0&0\end{bmatrix}$ and $A = \begin{bmatrix}1&4&5\\4&2&6\\5&6&3\end{bmatrix}$, show that $PAQ$ is symmetric when $Q = P^\T$, and factor $A = LDL^\T$.\\
\textbf{2.7 C:} Apply block elimination to find $S = LDL^\T$ for the saddle-point matrix $S = \begin{bmatrix}I & A\\A^\T & \vect{0}\end{bmatrix}$.
\end{example}
\textbf{Solution:}
\textbf{2.7 A and 2.7 B:} $PA P^\T = \begin{bmatrix}2&6&4\\6&3&5\\4&5&1\end{bmatrix}$ is symmetric because $(PAP^\T)^\T = PA^\T P^\T = PAP^\T$. Eliminating on $A$ gives pivots $1, -14, -8$ and multipliers $\ell_{21}=4, \ell_{31}=5, \ell_{32}=1$:
\[
A = LDL^\T = \begin{bmatrix}1&0&0\\4&1&0\\5&1&1\end{bmatrix}\begin{bmatrix}1&0&0\\0&-14&0\\0&0&-8\end{bmatrix}\begin{bmatrix}1&4&5\\0&1&1\\0&0&1\end{bmatrix}.
\]
\textbf{2.7 C:} Subtract $A^\T$ times block row 1 from block row 2:
\[
S = \begin{bmatrix}I & A\\A^\T & \vect{0}\end{bmatrix} = \begin{bmatrix}I & \vect{0}\\A^\T & I\end{bmatrix}\begin{bmatrix}I & \vect{0}\\\vect{0} & -A^\T A\end{bmatrix}\begin{bmatrix}I & A\\\vect{0} & I\end{bmatrix} = LDL^\T.
\]
$S$ is invertible $\iff A^\T A$ is invertible $\iff A$ has independent columns.

\begin{example}{5.1 A and 5.1 B}{Row Operations and Eigenvalue Determinants}
\textbf{5.1 A:} How are the determinants of $M_1 = [(-1)^{i+j}a_{ij}]$, $M_2 = \begin{bmatrix}\text{row 1}-\text{row 3}\\\text{row 2}-\text{row 1}\\\text{row 3}-\text{row 2}\end{bmatrix}$, and $M_3 = \begin{bmatrix}\text{row 1}+\text{row 3}\\\text{row 2}+\text{row 1}\\\text{row 3}+\text{row 2}\end{bmatrix}$ related to $\det A$?\\
\textbf{5.1 B:} Explain by row operations why $\det\begin{bmatrix}1-a & 1 & 1\\1 & 1-a & 1\\1 & 1 & 1-a\end{bmatrix} = a^2(3-a)$.
\end{example}
\textbf{Solution:}
\textbf{5.1 A:} $M_1 = \operatorname{diag}(1,-1,1)A\operatorname{diag}(1,-1,1)$, so $\det M_1 = (-1)(\det A)(-1) = \det A$. In $M_2$, the three rows add to the zero row, so $\det M_2 = 0$. In $M_3$, splitting into 8 determinants by linearity leaves only the original rows and a 2-exchange permutation of the rows: $\det M_3 = \det A + \det A = 2\det A$.
\textbf{5.1 B:} Subtract row 3 from row 1 and row 2 to get $\begin{bmatrix}-a & 0 & a\\0 & -a & a\\1 & 1 & 1-a\end{bmatrix}$. Add columns 1 and 2 to column 3 to get the lower triangular matrix $\begin{bmatrix}-a & 0 & 0\\0 & -a & 0\\1 & 1 & 3-a\end{bmatrix}$, whose determinant is $(-a)(-a)(3-a) = a^2(3-a)$.

\begin{example}{5.2 A and 5.2 B}{Hessenberg Fibonacci Determinants and Permutation Signs}
\textbf{5.2 A:} Use cofactors of row 1 to show that the Hessenberg matrices $H_2=\begin{bmatrix}2&1\\1&2\end{bmatrix}, H_3=\begin{bmatrix}2&1&0\\1&2&1\\1&1&2\end{bmatrix}, H_4=\begin{bmatrix}2&1&0&0\\1&2&1&0\\1&1&2&1\\1&1&1&2\end{bmatrix}$ satisfy $|H_n| = |H_{n-1}| + |H_{n-2}| = F_{n+2}$.\\
\textbf{5.2 B:} (1) If $A$ is the $10\times 10$ all-ones matrix, how does the big formula give $\det A = 0$? (2) Is the product of all $n!$ permutations odd or even? (3) Why is $\det A$ unchanged if each $a_{ij}$ is multiplied by $i/j$?
\end{example}
\textbf{Solution:}
\textbf{5.2 A:} For $H_4$, $C_{11} = |H_3|$ and $C_{12} = -\begin{vmatrix}1&1&0\\1&2&1\\1&1&2\end{vmatrix} = -\begin{vmatrix}2&1&0\\1&2&1\\1&1&2\end{vmatrix} + \begin{vmatrix}1&0&0\\1&2&1\\1&1&2\end{vmatrix} = -|H_3| + |H_2|$. Thus $|H_4| = 2|H_3| - |H_3| + |H_2| = |H_3| + |H_2|$. With $|H_1|=2=F_3$ and $|H_2|=3=F_4$, we get $|H_n| = F_{n+2}$.
\textbf{5.2 B:} (1) All $10!$ products equal $1$; half have $+$ signs and half have $-$ signs, so they cancel to $0$. (2) Odd for $n=2,3$, and even for all $n \ge 4$ (since $n!/2$ is an even number of odd permutations). (3) Every term $a_{1\alpha}a_{2\beta}\cdots a_{n\omega}$ is multiplied by $1\cdot 2\cdots n$ from the rows and divided by $1\cdot 2\cdots n$ from the columns (or $DAD^{-1}$ has determinant $\det A$).

\begin{example}{5.3 A and 5.3 B}{Nullspace from Cofactors and Cramer's Rule}
\textbf{5.3 A:} Solve $A\vect{x} = \vect{0}$ using cofactors along a row for $A = \begin{bmatrix}1&4&7\\2&3&9\\2&2&8\end{bmatrix}$ and $A = \begin{bmatrix}1&1&2\\1&1&1\\1&1&1\end{bmatrix}$.\\
\textbf{5.3 B:} Use Cramer's Rule and $A^{-1} = C^\T/\det A$ to solve $\begin{bmatrix}2&6&2\\1&4&2\\5&9&0\end{bmatrix}\begin{bmatrix}x\\y\\z\end{bmatrix} = \begin{bmatrix}0\\0\\1\end{bmatrix}$ and find $A^{-1}$.
\end{example}
\textbf{Solution:}
\textbf{5.3 A:} When $\det A = 0$, $AC^\T = \vect{0}$, so any nonzero row of cofactors solves $A\vect{x}=\vect{0}$. For the first matrix, cofactors along row 1 give $\vect{x} = (6, 2, -2)$. For the second matrix, cofactors along row 2 give $\vect{x} = (1, -1, 0)$.
\textbf{5.3 B:} $|B_1| = C_{31} = 4$, $|B_2| = C_{32} = -2$, $|B_3| = C_{33} = 2$, and $\det A = 5(4)+9(-2)+0(2) = 2$. Thus $\vect{x} = \frac{1}{2}(4,-2,2) = (2,-1,1)$, which is column 3 of $A^{-1} = \frac{1}{2}\begin{bmatrix}-18 & 18 & 4\\10 & -10 & -2\\-11 & 12 & 2\end{bmatrix}$.

% ============================================================================
% MODULE 3: DOT PRODUCTS, ORTHOGONALITY, POSITIVE DEFINITE, STATISTICS (PDF 1)
% ============================================================================
\part{Dot Products, Orthogonality, Positive Definiteness, and Statistics}

\chapter[Lengths, Dot Products, and Orthogonality (1.2, 4.1--4.4)]{Lengths, Dot Products, and Orthogonality\bk{1.2, 4.1--4.4}}

\begin{syll}
Syllabus topics: Dot product $\vect{v}\cdot\vect{w} = \vect{v}^\T\vect{w}$, length $\|\vect{v}\| = \sqrt{\vect{v}\cdot\vect{v}}$, unit vectors $\vect{u} = \vect{v}/\|\vect{v}\|$ $\cdot$ Perpendicular vectors ($\vect{v}\cdot\vect{w}=0$), Cosine Formula, Schwarz Inequality, Triangle Inequality $\cdot$ Orthogonality of the Four Subspaces (Fundamental Theorem, Part 2) $\cdot$ Projections onto a line ($P = \vect{a}\vect{a}^\T/\vect{a}^\T\vect{a}$) and onto a subspace ($A^\T A\uv{x} = A^\T\vect{b}$, $P = A(A^\T A)^{-1}A^\T$) $\cdot$ Least squares approximations and line fitting $\cdot$ Orthonormal bases ($Q^\T Q = I$), orthogonal matrices ($Q^\T = Q^{-1}$), and Gram-Schmidt ($A = QR$).
\end{syll}

\section{Lengths, Dot Products, Angles, and Inequalities\bk{1.2}}

\begin{goal}
Define the dot product $\vect{v}\cdot\vect{w}$ and vector length $\|\vect{v}\|$, and derive the right-angle test, Cosine Formula, Schwarz Inequality, and Triangle Inequality from Plus-Two geometry.
\end{goal}

\begin{base}
From \tref{tool:pyth} and \tref{tool:trig}, a right triangle with sides $a, b, c$ satisfies $a^2 + b^2 = c^2$, and $\cos(\beta - \alpha) = \cos\beta\cos\alpha + \sin\beta\sin\alpha$.
\end{base}

\step{Step 1: Define the dot product, length, and unit vector.}
For vectors $\vect{v} = (v_1, \dots, v_n)$ and $\vect{w} = (w_1, \dots, w_n)$ in $\mathbf{R}^n$:
\begin{align}
\text{Dot product (inner product):}\quad \vect{v}\cdot\vect{w} &= \vect{v}^\T\vect{w} = v_1w_1 + v_2w_2 + \dots + v_nw_n = \sum_{i=1}^n v_iw_i \\
\text{Length (norm):}\quad \|\vect{v}\| &= \sqrt{\vect{v}\cdot\vect{v}} = \sqrt{v_1^2 + v_2^2 + \dots + v_n^2} \\
\text{Unit vector (length $1$):}\quad \vect{u} &= \frac{\vect{v}}{\|\vect{v}\|} \qquad (\text{satisfies } \vect{u}\cdot\vect{u} = \|\vect{u}\|^2 = 1).
\end{align}

\step{Step 2: Prove that $\vect{v}$ is perpendicular to $\vect{w}$ if and only if $\vect{v}\cdot\vect{w} = 0$.}
When $\vect{v}$ and $\vect{w}$ meet at $90^\circ$, they form the legs of a right triangle with hypotenuse $\vect{v} - \vect{w}$. Apply Pythagoras' Law $\|\vect{v}\|^2 + \|\vect{w}\|^2 = \|\vect{v}-\vect{w}\|^2$:
\begin{align}
(v_1^2 + v_2^2) + (w_1^2 + w_2^2) &= (v_1 - w_1)^2 + (v_2 - w_2)^2 &&\why{expand squares} \nonumber\\
\cancel{v_1^2} + \cancel{v_2^2} + \cancel{w_1^2} + \cancel{w_2^2} &= \cancel{v_1^2} - 2v_1w_1 + \cancel{w_1^2} + \cancel{v_2^2} - 2v_2w_2 + \cancel{w_2^2} &&\why{cancel square terms} \nonumber\\
0 &= -2(v_1w_1 + v_2w_2) \implies \eqbox{\vect{v}\cdot\vect{w} = 0.} &&\why{divide by $-2$}
\end{align}

\step{Step 3: Derive the Cosine Formula, Schwarz Inequality, and Triangle Inequality.}
Take two unit vectors $\vect{u} = (\cos\alpha, \sin\alpha)$ and $\vect{U} = (\cos\beta, \sin\beta)$ separated by angle $\theta = \beta - \alpha$:
\begin{equation}
\vect{u}\cdot\vect{U} = \cos\alpha\cos\beta + \sin\alpha\sin\beta = \cos(\beta - \alpha) = \cos\theta.
\end{equation}
For any nonzero vectors $\vect{v}$ and $\vect{w}$, substitute the unit vectors $\vect{u} = \vect{v}/\|\vect{v}\|$ and $\vect{U} = \vect{w}/\|\vect{w}\|$:
\begin{equation}
\frac{\vect{v}}{\|\vect{v}\|}\cdot\frac{\vect{w}}{\|\vect{w}\|} = \cos\theta \implies \eqbox{\cos\theta = \frac{\vect{v}\cdot\vect{w}}{\|\vect{v}\|\|\vect{w}\|}.}
\end{equation}
Since $|\cos\theta| \le 1$, multiply by $\|\vect{v}\|\|\vect{w}\|$ to get the \term{Schwarz Inequality}:
\begin{equation}
|\vect{v}\cdot\vect{w}| \le \|\vect{v}\|\|\vect{w}\|.
\end{equation}
Now expand $\|\vect{v}+\vect{w}\|^2$ and apply $\vect{v}\cdot\vect{w} \le \|\vect{v}\|\|\vect{w}\|$:
\begin{equation}
\begin{aligned}
\|\vect{v}+\vect{w}\|^2 &= (\vect{v}+\vect{w})\cdot(\vect{v}+\vect{w}) = \|\vect{v}\|^2 + 2\vect{v}\cdot\vect{w} + \|\vect{w}\|^2 \\
&\le \|\vect{v}\|^2 + 2\|\vect{v}\|\|\vect{w}\| + \|\vect{w}\|^2 = (\|\vect{v}\| + \|\vect{w}\|)^2.
\end{aligned}
\end{equation}
Take the square root of both sides to get the \term{Triangle Inequality}: $\|\vect{v}+\vect{w}\| \le \|\vect{v}\| + \|\vect{w}\|$.

\begin{pic}
\begin{tikzpicture}[scale=0.72]
  % Left: Right triangle v, w, v-w
  \begin{scope}[xshift=0cm]
    \draw[vec,teal] (0,0) -- (3,1.5) node[right] {$\vect{v}=\begin{bmatrix}4\\2\end{bmatrix}$};
    \draw[vec,amber] (0,0) -- (-0.75,1.5) node[left] {$\vect{w}=\begin{bmatrix}-1\\2\end{bmatrix}$};
    \draw[vec,rose,thick] (-0.75,1.5) -- (3,1.5) node[midway,above] {$\sqrt{25}=\|\vect{v}-\vect{w}\|$};
    \node[font=\footnotesize] at (1.1,-0.6) {$\vect{v}\cdot\vect{w}=0 \implies 5+20=25$};
  \end{scope}
  % Right: Unit circle cosine formula
  \begin{scope}[xshift=8.5cm]
    \draw[ax] (-1.5,0) -- (2.8,0) node[right] {$x$};
    \draw[ax] (0,-0.5) -- (0,2.6) node[above] {$y$};
    \draw[vec,teal] (0,0) -- (2.2,0.7) node[right] {$\begin{bmatrix}\cos\alpha\\\sin\alpha\end{bmatrix}$};
    \draw[vec,amber] (0,0) -- (0.8,2.15) node[above right] {$\begin{bmatrix}\cos\beta\\\sin\beta\end{bmatrix}$};
    \draw[slate] (0.8,0) arc (0:18:0.8) node[right,yshift=-2pt,font=\footnotesize] {$\alpha$};
    \draw[plum,thick] (0.95,0.3) arc (18:69:1.0) node[midway,above right,font=\footnotesize] {$\theta=\beta-\alpha$};
  \end{scope}
\end{tikzpicture}
\figcap{Perpendicular vectors satisfy $\|\vect{v}\|^2+\|\vect{w}\|^2=\|\vect{v}-\vect{w}\|^2$ (left) and unit vectors have $\vect{u}\cdot\vect{U}=\cos\theta$ (right) \bookref{Figs.\ 1.8--1.9}.}
\end{pic}

\begin{example}{1.2.1--1.2.6 (p.\ 11--16)}{Dot Products, Unit Vectors, Schwarz and AM-GM Inequalities}
\begin{enumerate}
\item \textbf{Examples 1--3:} Verify that $\vect{v}=(4,2)$ and $\vect{w}=(-1,2)$ have a zero dot product. Explain how $(4)(-1)+(2)(2)=0$ balances weights $4,2$ at distances $-1,2$ on a see-saw, and how $\text{Income} = \vect{q}\cdot\vect{p} = q_1p_1+q_2p_2+q_3p_3$.
\item \textbf{Example 4:} Since $(2,2,1)$ has length $3$, check that $\vect{u} = (\frac{2}{3}, \frac{2}{3}, \frac{1}{3})$ is a unit vector.
\item \textbf{Example 5:} Find $\cos\theta$ for $\vect{v} = \begin{bmatrix}2\\1\end{bmatrix}$ and $\vect{w} = \begin{bmatrix}1\\2\end{bmatrix}$ and check both inequalities (Schwarz and Triangle).
\item \textbf{Example 6:} Use the Schwarz inequality on $\vect{v}=(a,b)$ and $\vect{w}=(b,a)$ with $x=a^2, y=b^2$ to prove that the geometric mean $\sqrt{xy}$ is not larger than the arithmetic mean $\frac{1}{2}(x+y)$.
\end{enumerate}
\end{example}

\vfill
\begin{pageend_sol}{1.2.1--1.2.6 (p.\ 11--16)}
\textbf{1.} $\vect{v}\cdot\vect{w} = 4(-1) + 2(2) = -4 + 4 = 0$ ($90^\circ$ angle). On a see-saw, the moments $w_1v_1 + w_2v_2 = 0$ balance. In economics, $\vect{q}\cdot\vect{p} = 0$ means total sales equal total purchases ("the books balance").

\textbf{2.} $\vect{u}\cdot\vect{u} = \frac{4}{9} + \frac{4}{9} + \frac{1}{9} = 1$, so $\|\vect{u}\| = 1$.

\textbf{3.} $\vect{v}\cdot\vect{w} = 4$, $\|\vect{v}\| = \sqrt{5}$, $\|\vect{w}\| = \sqrt{5}$, so $\cos\theta = \frac{4}{\sqrt{5}\sqrt{5}} = \frac{4}{5}$. Schwarz: $|\vect{v}\cdot\vect{w}| = 4 \le \|\vect{v}\|\|\vect{w}\| = 5$. Triangle: $\|\vect{v}+\vect{w}\| = \|(3,3)\| = \sqrt{18} \le \sqrt{5}+\sqrt{5} = \sqrt{20}$.

\textbf{4.} $\vect{v}\cdot\vect{w} = 2ab$ and $\|\vect{v}\|\|\vect{w}\| = \sqrt{a^2+b^2}\sqrt{b^2+a^2} = a^2+b^2$. Schwarz gives $2ab \le a^2+b^2 \implies ab \le \frac{a^2+b^2}{2}$. Substitute $x=a^2, y=b^2$ to get $\sqrt{xy} \le \frac{x+y}{2}$.
\end{pageend_sol}
\clearpage

\section{Orthogonality of the Four Subspaces\bk{4.1}}

\begin{goal}
Derive Part 2 of the Fundamental Theorem of Linear Algebra: the row space and nullspace are orthogonal complements in $\mathbf{R}^n$, and the column space and left nullspace are orthogonal complements in $\mathbf{R}^m$.
\end{goal}

\step{Step 1: Prove that $N(A)$ and $C(A^\T)$ are orthogonal in $\mathbf{R}^n$.}
Take any $\vect{x}$ in the nullspace $N(A)$ ($A\vect{x}=\vect{0}$) and any row-space vector $A^\T\vect{y}$ in $C(A^\T)$:
\begin{equation}
\vect{x}^\T(A^\T\vect{y}) = (A\vect{x})^\T\vect{y} = \vect{0}^\T\vect{y} = 0.
\end{equation}
Visually, each row of $A\vect{x} = \vect{0}$ says $(\text{row } i \text{ of } A)\cdot\vect{x} = 0$. Therefore $\vect{x}$ is perpendicular to every row of $A$ and to every combination of the rows.

\step{Step 2: Prove that $N(A^\T)$ and $C(A)$ are orthogonal in $\mathbf{R}^m$.}
Take any $\vect{y}$ in the left nullspace $N(A^\T)$ ($A^\T\vect{y}=\vect{0}$) and any column-space vector $A\vect{x}$ in $C(A)$:
\begin{equation}
(A\vect{x})^\T\vect{y} = \vect{x}^\T(A^\T\vect{y}) = \vect{x}^\T\vect{0} = 0.
\end{equation}

\step{Step 3: Orthogonal complements and splitting $\vect{x} = \vect{x}_r + \vect{x}_n$.}
The \term{orthogonal complement} $V^\perp$ of a subspace $V$ contains \emph{every} vector perpendicular to $V$. Because $\dim C(A^\T) + \dim N(A) = r + (n-r) = n$, the row space and nullspace are \term{orthogonal complements}:
\begin{result}[ (Fundamental Theorem of Linear Algebra, Part 2)]
\begin{align}
N(A) &= C(A^\T)^\perp \quad \text{and} \quad C(A^\T) = N(A)^\perp \qquad (\text{in } \mathbf{R}^n) \\
N(A^\T) &= C(A)^\perp \quad \text{and} \quad C(A) = N(A^\T)^\perp \qquad (\text{in } \mathbf{R}^m).
\end{align}
Every vector $\vect{x} \in \mathbf{R}^n$ splits uniquely into a row space component $\vect{x}_r$ and a nullspace component $\vect{x}_n$:
\begin{equation}
\vect{x} = \vect{x}_r + \vect{x}_n, \qquad A\vect{x} = A\vect{x}_r + A\vect{x}_n = A\vect{x}_r + \vect{0} = \vect{b}.
\end{equation}
From the row space to the column space, $A$ is one-to-one and invertible.
\end{result}

\begin{pic}
\begin{tikzpicture}[scale=0.75]
  % Left: R^n
  \draw[thick,teal,fill=tealsoft!40] (0,0) -- (-2,2) -- (0,4) -- (2,2) -- cycle;
  \draw[thick,amber,fill=ambersoft!50] (0,0) -- (-1.6,-1.6) -- (0,-3.2) -- (1.6,-1.6) -- cycle;
  \draw[slate] (-0.25,0.25) -- (0,0.5) -- (0.25,0.25);
  \fill[teal] (0.5,2.2) circle (2.5pt) node[left,font=\footnotesize] {$\vect{x}_r$};
  \fill[amber!80!black] (0.3,-1.2) circle (2.5pt) node[left,font=\footnotesize] {$\vect{x}_n$};
  \fill[ink] (0.8,1.0) circle (2.5pt) node[right,font=\footnotesize] {$\vect{x}=\vect{x}_r+\vect{x}_n$};
  \draw[dashed,slate] (0.5,2.2) -- (0.8,1.0) -- (0.3,-1.2);
  \node[font=\small\bfseries,teal] at (-0.6,2.8) {row space};
  \node[font=\small\bfseries,amber!80!black] at (0,-2.2) {nullspace};
  \node[font=\bfseries] at (-2.2,0) {$\mathbf{R}^n$};

  % Right: R^m
  \begin{scope}[xshift=8.5cm]
    \draw[thick,teal,fill=tealsoft!40] (0,0) -- (-1.4,2.4) -- (0.8,3.8) -- (2.2,1.4) -- cycle;
    \draw[thick,plum,fill=plumsoft!50] (0,0) -- (-2.2,-1) -- (-1.4,-2.4) -- (0.8,-1.4) -- cycle;
    \fill[teal] (0.6,2.0) circle (2.5pt) node[above,font=\footnotesize] {$\vect{b}$};
    \fill[ink] (0,0) circle (2.5pt);
    \node[font=\small\bfseries,teal] at (1.1,2.6) {column space};
    \node[font=\small\bfseries,plum] at (-0.7,-1.4) {nullspace of $A^\T$};
    \node[font=\bfseries] at (2.0,-0.2) {$\mathbf{R}^m$};
  \end{scope}

  % Arrows
  \draw[vec,teal,thick] (0.6,2.2) -- (8.9,2.0) node[midway,above,font=\footnotesize] {$A\vect{x}_r = \vect{b}$};
  \draw[vec,ink] (0.9,1.0) -- (8.9,1.95) node[midway,below,font=\footnotesize] {$A\vect{x}$};
  \draw[vec,amber!80!black,thick] (0.4,-1.2) -- (8.3,0) node[midway,below,font=\footnotesize] {$A\vect{x}_n = \vect{0}$};
\end{tikzpicture}
\figcap{The true action of $A$ on $\vect{x} = \vect{x}_r + \vect{x}_n$: row space vector $\vect{x}_r$ goes to $\vect{b}$, nullspace vector $\vect{x}_n$ goes to $\vect{0}$ \bookref{Fig.\ 4.3}.}
\end{pic}

\begin{example}{4.1.3--4.1.5 (p.\ 197--200)}{Orthogonality to the Nullspace and Splitting $\vect{x} = \vect{x}_r + \vect{x}_n$}
\begin{enumerate}
\item \textbf{Example 3:} Verify that the rows of $A = \begin{bmatrix}1 & 3 & 4\\5 & 2 & 7\end{bmatrix}$ are perpendicular to $\vect{x} = (1,1,-1)$ in the nullspace.
\item \textbf{Example 4:} Identify the $r \times r$ invertible submatrices inside $A = \begin{bmatrix}3&0&0&0&0\\0&5&0&0&0\\0&0&0&0&0\end{bmatrix}$ and $B = \begin{bmatrix}1&2&3&4&5\\1&2&4&5&6\\1&2&4&5&6\end{bmatrix}$.
\item \textbf{Example 5:} For $A = \begin{bmatrix}1 & 2\\3 & 6\end{bmatrix}$, split $\vect{x} = \begin{bmatrix}4\\3\end{bmatrix}$ into $\vect{x}_r + \vect{x}_n$.
\end{enumerate}
\end{example}

\vfill
\begin{pageend_sol}{4.1.3--4.1.5 (p.\ 197--200)}
\textbf{1.} $A\vect{x} = \begin{bmatrix}1 & 3 & 4\\5 & 2 & 7\end{bmatrix}\begin{bmatrix}1\\1\\-1\end{bmatrix} = \begin{bmatrix}1+3-4\\5+2-7\end{bmatrix} = \begin{bmatrix}0\\0\end{bmatrix}$. Both dot products are zero.

\textbf{2.} Both matrices have rank $r=2$. Inside $A$ is $\begin{bmatrix}3&0\\0&5\end{bmatrix}$; inside $B$ (in pivot rows 1, 2 and pivot columns 1, 3) is $\begin{bmatrix}1&3\\1&4\end{bmatrix}$.

\textbf{3.} The row space is the line through $(1,2)$ and the nullspace is the line through $(-2,1)$. Splitting $\vect{x} = \begin{bmatrix}4\\3\end{bmatrix}$ gives $\vect{x}_r + \vect{x}_n = \begin{bmatrix}2\\4\end{bmatrix} + \begin{bmatrix}2\\-1\end{bmatrix}$. Notice $\vect{x}_r\cdot\vect{x}_n = 4 - 4 = 0$.
\end{pageend_sol}
\clearpage

\section{Projections Onto Lines and Subspaces\bk{4.2}}

\begin{goal}
Derive the projection vector $\vect{p}$ and projection matrix $P$ onto a line through $\vect{a}$ and onto the column space $C(A)$ of an $m \times n$ matrix $A$ with independent columns.
\end{goal}

\step{Step 1: Project $\vect{b}$ onto a line through $\vect{a}$.}
The projection $\vect{p} = \uv{x}\vect{a}$ is the multiple of $\vect{a}$ closest to $\vect{b}$. The error vector $\vect{e} = \vect{b} - \uv{x}\vect{a}$ must be perpendicular to $\vect{a}$:
\begin{equation}
\vect{a}\cdot(\vect{b} - \uv{x}\vect{a}) = 0 \implies \vect{a}\cdot\vect{b} - \uv{x}(\vect{a}\cdot\vect{a}) = 0 \implies \eqbox{\uv{x} = \frac{\vect{a}\cdot\vect{b}}{\vect{a}\cdot\vect{a}} = \frac{\vect{a}^\T\vect{b}}{\vect{a}^\T\vect{a}}.}
\end{equation}
Multiply $\vect{a}$ by $\uv{x}$ to get $\vect{p}$ and factor out $\vect{b}$ on the right to get the rank-one projection matrix $P$:
\begin{equation}
\vect{p} = \vect{a}\uv{x} = \vect{a}\frac{\vect{a}^\T\vect{b}}{\vect{a}^\T\vect{a}} = P\vect{b}, \qquad \text{where} \qquad \eqbox{P = \frac{\vect{a}\vect{a}^\T}{\vect{a}^\T\vect{a}}.}
\end{equation}

\step{Step 2: Project $\vect{b}$ onto the $n$-dimensional subspace $C(A)$.}
We look for the combination $\vect{p} = A\uv{\vect{x}}$ in the column space of $A$ closest to $\vect{b}$. The error vector $\vect{e} = \vect{b} - A\uv{\vect{x}}$ must be perpendicular to all columns $\vect{a}_1, \dots, \vect{a}_n$ of $A$ (so $\vect{e}$ lies in the left nullspace $N(A^\T)$):
\begin{equation}
\begin{bmatrix} -\vect{a}_1^\T- \\ \vdots \\ -\vect{a}_n^\T- \end{bmatrix}(\vect{b} - A\uv{\vect{x}}) = \vect{0} \implies A^\T(\vect{b} - A\uv{\vect{x}}) = \vect{0} \implies \eqbox{A^\T A\uv{\vect{x}} = A^\T\vect{b}.}
\end{equation}

\step{Step 3: Prove that $A^\T A$ is invertible when $A$ has independent columns.}
Suppose $A^\T A\vect{x} = \vect{0}$. Multiply on the left by $\vect{x}^\T$:
\begin{equation}
\vect{x}^\T A^\T A\vect{x} = 0 \implies (A\vect{x})^\T(A\vect{x}) = 0 \implies \|A\vect{x}\|^2 = 0 \implies A\vect{x} = \vect{0}.
\end{equation}
Thus $A^\T A$ has the exact same nullspace as $A$. When $A$ has independent columns, $N(A) = \{\vect{0}\}$, so the square symmetric matrix $A^\T A$ is invertible!

\begin{result}[ (Three Key Formulas for Projection Onto $C(A)$)]
\begin{equation}
\text{\textbf{(1) Coefficient vector:}}\quad \uv{\vect{x}} = (A^\T A)^{-1}A^\T\vect{b} \qquad
\text{\textbf{(2) Projection:}}\quad \vect{p} = A\uv{\vect{x}} = A(A^\T A)^{-1}A^\T\vect{b}
\end{equation}
\begin{equation}
\text{\textbf{(3) Projection matrix:}}\quad P = A(A^\T A)^{-1}A^\T \qquad (\text{satisfies } P^\T = P \text{ and } P^2 = P).
\end{equation}
The matrix $I - P$ projects onto the orthogonal complement $N(A^\T)$ and produces the error $\vect{e} = (I - P)\vect{b} = \vect{b} - \vect{p}$.
\end{result}

\begin{pic}
\begin{tikzpicture}[scale=0.78]
  % Left: Projection onto a line
  \begin{scope}[xshift=0cm]
    \draw[vec,teal,thick] (0,0) -- (4.2,1.4) node[right] {$\vect{a}$};
    \draw[vec,ink,thick] (0,0) -- (2.1,3.2) node[above left] {$\vect{b}$};
    \draw[vec,amber!90!black,very thick] (0,0) -- (2.7,0.9) node[below,yshift=-2pt] {$\vect{p} = \uv{x}\vect{a} = \frac{\vect{a}^\T\vect{b}}{\vect{a}^\T\vect{a}}\vect{a}$};
    \draw[dashed,rose,thick] (2.7,0.9) -- (2.1,3.2) node[midway,right] {$\vect{e}=\vect{b}-\vect{p}$};
    \draw[slate] (2.45,0.82) -- (2.37,1.07) -- (2.62,1.15);
  \end{scope}
  % Right: Projection onto a subspace S = C(A)
  \begin{scope}[xshift=7.5cm]
    \draw[thick,teal,fill=tealsoft!40] (-0.8,-0.5) -- (3.8,0.4) -- (4.8,2.2) -- (0.2,1.3) -- cycle;
    \node[teal,font=\small\bfseries] at (0.1,-0.1) {$S = C(A)$};
    \draw[vec,ink,thick] (1.0,0.4) -- (2.4,3.4) node[above left] {$\vect{b}$};
    \draw[vec,amber!90!black,very thick] (1.0,0.4) -- (3.2,1.1) node[right] {$\vect{p}=A\uv{\vect{x}}=P\vect{b}$};
    \draw[dashed,rose,thick] (3.2,1.1) -- (2.4,3.4) node[midway,right] {$\vect{e}=\vect{b}-A\uv{\vect{x}}$};
    \draw[slate] (2.98,1.03) -- (2.89,1.28) -- (3.11,1.35);
  \end{scope}
\end{tikzpicture}
\figcap{Projection $\vect{p}$ of $\vect{b}$ onto a line through $\vect{a}$ (left) and onto the column space $S = C(A)$ (right) \bookref{Fig.\ 4.5}.}
\end{pic}

\begin{example}{4.2.1--4.2.3 (p.\ 208--211)}{Projecting Onto a Line and Onto a Plane}
\begin{enumerate}
\item \textbf{Examples 1--2:} Project $\vect{b} = \begin{bmatrix}1\\1\\1\end{bmatrix}$ onto $\vect{a} = \begin{bmatrix}1\\2\\2\end{bmatrix}$ to find $\vect{p} = \uv{x}\vect{a}$, the error $\vect{e} = \vect{b}-\vect{p}$, and the projection matrix $P = \frac{\vect{a}\vect{a}^\T}{\vect{a}^\T\vect{a}}$.
\item \textbf{Example 3:} If $A = \begin{bmatrix}1 & 0\\1 & 1\\1 & 2\end{bmatrix}$ and $\vect{b} = \begin{bmatrix}6\\0\\0\end{bmatrix}$, find $\uv{\vect{x}}$ and $\vect{p}$ and $P$.
\end{enumerate}
\end{example}

\vfill
\begin{pageend_sol}{4.2.1--4.2.3 (p.\ 208--211)}
\textbf{1.} $\vect{a}^\T\vect{b} = 1+2+2 = 5$ and $\vect{a}^\T\vect{a} = 1+4+4 = 9$, so $\uv{x} = \frac{5}{9}$. Thus $\vect{p} = \frac{5}{9}\vect{a} = \left(\frac{5}{9}, \frac{10}{9}, \frac{10}{9}\right)$, $\vect{e} = \vect{b}-\vect{p} = \left(\frac{4}{9}, -\frac{1}{9}, -\frac{1}{9}\right)$ (check $\vect{e}^\T\vect{a} = 0$), and $P = \frac{1}{9}\begin{bmatrix}1\\2\\2\end{bmatrix}\begin{bmatrix}1 & 2 & 2\end{bmatrix} = \frac{1}{9}\begin{bmatrix}1&2&2\\2&4&4\\2&4&4\end{bmatrix}$.

\textbf{2.} $A^\T A = \begin{bmatrix}3 & 3\\3 & 5\end{bmatrix}$ and $A^\T\vect{b} = \begin{bmatrix}6\\0\end{bmatrix}$. Solve $\begin{bmatrix}3 & 3\\3 & 5\end{bmatrix}\uv{\vect{x}} = \begin{bmatrix}6\\0\end{bmatrix}$ to get $\uv{\vect{x}} = \begin{bmatrix}5\\-3\end{bmatrix}$. Then $\vect{p} = A\uv{\vect{x}} = 5\begin{bmatrix}1\\1\\1\end{bmatrix} - 3\begin{bmatrix}0\\1\\2\end{bmatrix} = \begin{bmatrix}5\\2\\-1\end{bmatrix}$, $\vect{e} = \vect{b}-\vect{p} = \begin{bmatrix}1\\-2\\1\end{bmatrix}$, and $P = A(A^\T A)^{-1}A^\T = \frac{1}{6}\begin{bmatrix}5 & 2 & -1\\2 & 2 & 2\\-1 & 2 & 5\end{bmatrix}$.
\end{pageend_sol}
\clearpage

\section{Least Squares Approximations\bk{4.3}}

\begin{goal}
Derive the normal equations $A^\T A\uv{\vect{x}} = A^\T\vect{b}$ that minimize $E = \|A\vect{x}-\vect{b}\|^2$ by geometry, algebra, and calculus, and apply them to straight-line and parabola fitting.
\end{goal}

\step{Step 1: Derive $A^\T A\uv{\vect{x}} = A^\T\vect{b}$ in three ways.}
When $A\vect{x} = \vect{b}$ has more equations than unknowns ($m > n$) and no exact solution, we minimize the squared error $E(\vect{x}) = \|A\vect{x} - \vect{b}\|^2$:
\begin{enumerate}
\item \textbf{By geometry:} The closest point to $\vect{b}$ in $C(A)$ is the projection $\vect{p} = A\uv{\vect{x}}$, where $\vect{e} = \vect{b} - A\uv{\vect{x}} \perp C(A)$, giving $A^\T(\vect{b} - A\uv{\vect{x}}) = \vect{0}$.
\item \textbf{By algebra (Pythagoras):} Split $\vect{b} = \vect{p} + \vect{e}$. Since $A\vect{x} - \vect{p} \in C(A)$ is perpendicular to $\vect{e} \in N(A^\T)$:
\begin{equation}
\|A\vect{x} - \vect{b}\|^2 = \|(A\vect{x} - \vect{p}) - \vect{e}\|^2 = \|A\vect{x} - \vect{p}\|^2 + \|\vect{e}\|^2.
\end{equation}
Choosing $\vect{x} = \uv{\vect{x}}$ makes $\|A\uv{\vect{x}} - \vect{p}\|^2 = 0$, leaving the minimum possible squared error $\|\vect{e}\|^2$.
\item \textbf{By calculus:} For fitting a line $b = C + Dt$ to $m$ points $(t_1,b_1), \dots, (t_m,b_m)$, set $\partial E/\partial C = 0$ and $\partial E/\partial D = 0$ in $E = \sum_{i=1}^m (C + Dt_i - b_i)^2$:
\begin{equation}
\frac{\partial E}{\partial C} = 2\sum_{i=1}^m (C + Dt_i - b_i) = 0, \qquad \frac{\partial E}{\partial D} = 2\sum_{i=1}^m (C + Dt_i - b_i)t_i = 0.
\end{equation}
\end{enumerate}

\begin{result}[ (Normal Equations for Straight-Line and Shifted-Time Fitting)]
For $A = \begin{bmatrix}1 & t_1\\\vdots & \vdots\\1 & t_m\end{bmatrix}$, the normal equations $A^\T A\uv{\vect{x}} = A^\T\vect{b}$ are:
\begin{equation}
\begin{bmatrix} m & \sum t_i \\ \sum t_i & \sum t_i^2 \end{bmatrix} \begin{bmatrix} C \\ D \end{bmatrix} = \begin{bmatrix} \sum b_i \\ \sum t_i b_i \end{bmatrix}.
\end{equation}
If the shifted times $T_i = t_i - \uv{t}$ add to $\sum T_i = 0$ (where $\uv{t} = \frac{1}{m}\sum t_i$), the columns of $A$ are orthogonal and $A^\T A$ is diagonal:
\begin{equation}
\text{Best line } b = C + D(t - \uv{t}) \quad \text{has} \quad C = \frac{b_1 + \dots + b_m}{m} \quad \text{and} \quad D = \frac{b_1T_1 + \dots + b_mT_m}{T_1^2 + \dots + T_m^2}.
\end{equation}
\end{result}

\begin{example}{4.3.1--4.3.3 (p.\ 218--224)}{Fitting a Straight Line and a Parabola}
\begin{enumerate}
\item \textbf{Example 1:} Find the closest line $b = C + Dt$ to the three points $(0,6), (1,0)$, and $(2,0)$.
\item \textbf{Example 2:} Fit $b = 1, 2, 4$ at times $t = -2, 0, 2$ (which add to zero) by $C + Dt$.
\item \textbf{Example 3:} Fit a parabola $b = C + Dt + Et^2$ through the three heights $b = 6, 0, 0$ at $t = 0, 1, 2$.
\end{enumerate}
\end{example}

\vfill
\begin{pageend_sol}{4.3.1--4.3.3 (p.\ 218--224)}
\textbf{1.} $A = \begin{bmatrix}1&0\\1&1\\1&2\end{bmatrix}$ and $\vect{b} = \begin{bmatrix}6\\0\\0\end{bmatrix}$. This is the exact system from Example 4.2.3: $A^\T A\uv{\vect{x}} = \begin{bmatrix}3&3\\3&5\end{bmatrix}\begin{bmatrix}C\\D\end{bmatrix} = \begin{bmatrix}6\\0\end{bmatrix} \implies C=5, D=-3$. The best line is $b = 5 - 3t$ with heights $\vect{p}=(5,2,-1)$ and vertical errors $\vect{e}=(1,-2,1)$.

\textbf{2.} $A = \begin{bmatrix}1&-2\\1&0\\1&2\end{bmatrix}$ has orthogonal columns: $A^\T A\uv{\vect{x}} = \begin{bmatrix}3&0\\0&8\end{bmatrix}\begin{bmatrix}C\\D\end{bmatrix} = \begin{bmatrix}7\\6\end{bmatrix} \implies C = \frac{7}{3}$ and $D = \frac{6}{8} = \frac{3}{4}$.

\textbf{3.} Three points give an exact $3\times 3$ system: $C=6$, $C+D+E=0$, $C+2D+4E=0 \implies (C,D,E) = (6,-9,3)$, so $b = 6 - 9t + 3t^2$ ($\vect{p}=\vect{b}$ and $\vect{e}=\vect{0}$).
\end{pageend_sol}
\clearpage

\section{Orthogonal Bases and Gram-Schmidt ($A = QR$)\bk{4.4}}

\begin{goal}
Define orthonormal columns ($Q^\T Q = I$) and orthogonal matrices ($Q^\T = Q^{-1}$), simplify projections using $Q$, and derive the Gram-Schmidt orthogonalization process $A = QR$.
\end{goal}

\step{Step 1: Orthonormal vectors and matrices $Q$.}
The vectors $\vect{q}_1, \dots, \vect{q}_n$ are \term{orthonormal} if $\vect{q}_i^\T\vect{q}_j = 0$ for $i \neq j$ and $\vect{q}_i^\T\vect{q}_i = \|\vect{q}_i\|^2 = 1$:
\begin{equation}
Q = \begin{bmatrix} \vect{q}_1 & \vect{q}_2 & \dots & \vect{q}_n \end{bmatrix} \quad \text{satisfies} \quad Q^\T Q = I_n.
\end{equation}
If $Q$ is square ($m=n$), then $Q^\T = Q^{-1}$ (and $QQ^\T = I$), and $Q$ is called an \term{orthogonal matrix}.
Multiplication by $Q$ preserves lengths and dot products:
\begin{equation}
\|Q\vect{x}\|^2 = (Q\vect{x})^\T(Q\vect{x}) = \vect{x}^\T Q^\T Q\vect{x} = \vect{x}^\T\vect{x} = \|\vect{x}\|^2, \qquad (Q\vect{x})^\T(Q\vect{y}) = \vect{x}^\T Q^\T Q\vect{y} = \vect{x}^\T\vect{y}.
\end{equation}

\step{Step 2: Projections and least squares when $A = Q$.}
Replace $A$ by $Q$ in the projection formulas and use $Q^\T Q = I$:
\begin{equation}
\uv{\vect{x}} = Q^\T\vect{b}, \qquad \vect{p} = Q\uv{\vect{x}} = QQ^\T\vect{b} = \vect{q}_1(\vect{q}_1^\T\vect{b}) + \dots + \vect{q}_n(\vect{q}_n^\T\vect{b}), \qquad P = QQ^\T.
\end{equation}
When $Q$ is square, $P = QQ^\T = I$, so every vector $\vect{b} \in \mathbf{R}^n$ is the sum of its 1D projections along the $\vect{q}$'s.

\step{Step 3: Derive the Gram-Schmidt process and $A = QR$.}
Start with independent vectors $\vect{a}, \vect{b}, \vect{c}$. Construct orthogonal vectors $\vect{A}, \vect{B}, \vect{C}$ by subtracting from each new vector its projections onto the directions already set:
\begin{align}
\text{\textbf{First step:}}\quad \vect{A} &= \vect{a}, \qquad \vect{B} = \vect{b} - \frac{\vect{A}^\T\vect{b}}{\vect{A}^\T\vect{A}}\vect{A} &&\why{check $\vect{A}^\T\vect{B} = 0$} \\
\text{\textbf{Next step:}}\quad \vect{C} &= \vect{c} - \frac{\vect{A}^\T\vect{c}}{\vect{A}^\T\vect{A}}\vect{A} - \frac{\vect{B}^\T\vect{c}}{\vect{B}^\T\vect{B}}\vect{B} &&\why{check $\vect{A}^\T\vect{C} = \vect{B}^\T\vect{C} = 0$} \\
\text{\textbf{Normalize:}}\quad \vect{q}_1 &= \frac{\vect{A}}{\|\vect{A}\|}, \qquad \vect{q}_2 = \frac{\vect{B}}{\|\vect{B}\|}, \qquad \vect{q}_3 = \frac{\vect{C}}{\|\vect{C}\|}.
\end{align}
Write $\vect{a}, \vect{b}, \vect{c}$ in terms of $\vect{q}_1, \vect{q}_2, \vect{q}_3$ (since later $\vect{q}$'s are orthogonal to earlier $\vect{a}$'s, $R = Q^\T A$ is upper triangular):
\begin{result}[ (Gram-Schmidt Factorization $A = QR$ and Least Squares)]
\begin{equation}
\begin{bmatrix} \vect{a} & \vect{b} & \vect{c} \end{bmatrix} = \begin{bmatrix} \vect{q}_1 & \vect{q}_2 & \vect{q}_3 \end{bmatrix} \begin{bmatrix} \vect{q}_1^\T\vect{a} & \vect{q}_1^\T\vect{b} & \vect{q}_1^\T\vect{c} \\ 0 & \vect{q}_2^\T\vect{b} & \vect{q}_2^\T\vect{c} \\ 0 & 0 & \vect{q}_3^\T\vect{c} \end{bmatrix} \quad \text{or} \quad A = QR \quad (R = Q^\T A).
\end{equation}
With $A = QR$, we have $A^\T A = R^\T Q^\T Q R = R^\T R$, so the least squares equation $A^\T A\uv{\vect{x}} = A^\T\vect{b}$ simplifies to:
\begin{equation}
R^\T R\uv{\vect{x}} = R^\T Q^\T\vect{b} \implies \eqbox{R\uv{\vect{x}} = Q^\T\vect{b} \quad \text{or} \quad \uv{\vect{x}} = R^{-1}Q^\T\vect{b}.}
\end{equation}
\end{result}

\begin{example}{4.4.1--4.4.5 (p.\ 231--236)}{Orthogonal Matrices and Gram-Schmidt $A = QR$}
\begin{enumerate}
\item \textbf{Examples 1--3:} Verify $Q^\T Q = I$ for (1) rotation $Q = \begin{bmatrix}\cos\theta & -\sin\theta\\\sin\theta & \cos\theta\end{bmatrix}$, (2) permutation matrices, and (3) reflection $Q = I - 2\vect{u}\vect{u}^\T$ ($\|\vect{u}\|=1$).
\item \textbf{Example 4:} Reconstruct $\vect{b} = (0,0,1)$ from its projections onto the columns of $Q = \frac{1}{3}\begin{bmatrix}-1&2&2\\2&-1&2\\2&2&-1\end{bmatrix}$.
\item \textbf{Example 5:} Apply Gram-Schmidt and find $A = QR$ for $\vect{a}=\begin{bmatrix}1\\-1\\0\end{bmatrix}, \vect{b}=\begin{bmatrix}2\\0\\-2\end{bmatrix}, \vect{c}=\begin{bmatrix}3\\-3\\3\end{bmatrix}$.
\end{enumerate}
\end{example}

\vfill
\begin{pageend_sol}{4.4.1--4.4.5 (p.\ 231--236)}
\textbf{1.} For reflection $Q = I - 2\vect{u}\vect{u}^\T$ with $\vect{u}^\T\vect{u}=1$: $Q^\T Q = (I-2\vect{u}\vect{u}^\T)^2 = I - 4\vect{u}\vect{u}^\T + 4\vect{u}(\vect{u}^\T\vect{u})\vect{u}^\T = I$.

\textbf{2.} $\vect{q}_1^\T\vect{b} = \frac{2}{3}, \vect{q}_2^\T\vect{b} = \frac{2}{3}, \vect{q}_3^\T\vect{b} = -\frac{1}{3}$, and $\frac{2}{3}\vect{q}_1 + \frac{2}{3}\vect{q}_2 - \frac{1}{3}\vect{q}_3 = \frac{1}{9}\begin{bmatrix}-2+4-2\\4-2-2\\4+4+1\end{bmatrix} = \begin{bmatrix}0\\0\\1\end{bmatrix} = \vect{b}$.

\textbf{3.} $\vect{A} = (1,-1,0)$ ($\vect{A}^\T\vect{A}=2$). $\vect{B} = \vect{b} - \frac{2}{2}\vect{A} = (1,1,-2)$ ($\vect{B}^\T\vect{B}=6$). $\vect{C} = \vect{c} - \frac{6}{2}\vect{A} - \frac{-6}{6}\vect{B} = (1,1,1)$ ($\vect{C}^\T\vect{C}=3$). Divide by $\sqrt{2}, \sqrt{6}, \sqrt{3}$:
\[
A = \begin{bmatrix}1&2&3\\-1&0&-3\\0&-2&3\end{bmatrix} = \begin{bmatrix}1/\sqrt{2} & 1/\sqrt{6} & 1/\sqrt{3}\\-1/\sqrt{2} & 1/\sqrt{6} & 1/\sqrt{3}\\0 & -2/\sqrt{6} & 1/\sqrt{3}\end{bmatrix} \begin{bmatrix}\sqrt{2} & \sqrt{2} & \sqrt{18}\\0 & \sqrt{6} & -\sqrt{6}\\0 & 0 & \sqrt{3}\end{bmatrix} = QR.
\]
\end{pageend_sol}
\clearpage

% ============================================================================
% CHAPTER 6: POSITIVE DEFINITE MATRICES AND STATISTICS (6.5, 8.6)
% ============================================================================
\chapter[Positive Definite Matrices and Statistics (6.5, 8.6)]{Positive Definite Matrices and Linear Algebra for Statistics\bk{6.5, 8.6}}

\begin{syll}
Syllabus topics: Symmetric positive definite matrices $\cdot$ Five equivalent tests (pivots, upper-left determinants, eigenvalues, energy $\vect{x}^\T A\vect{x} > 0$, and $A = R^\T R$) $\cdot$ Cholesky factorization $A = R^\T R$ $\cdot$ Semidefinite matrices $\cdot$ Ellipses $ax^2 + 2bxy + cy^2 = 1$ and Principal Axis Theorem $\cdot$ Weighted least squares $A^\T C A\uv{\vect{x}} = A^\T C\vect{b}$ $\cdot$ Mean, variance, Binomial, Poisson, Normal distribution, Covariance matrix $\Sigma$, and Principal Component Analysis (PCA).
\end{syll}

\section{Positive Definite Matrices\bk{6.5}}

\begin{goal}
Establish the five equivalent tests for a symmetric matrix $A = A^\T$ to be positive definite, and connect $\vect{x}^\T A\vect{x} = 1$ to the geometry of a tilted ellipse.
\end{goal}

\step{Step 1: Start with a $2\times 2$ symmetric matrix $A = \begin{bmatrix}a & b\\b & c\end{bmatrix}$.}
\begin{itemize}
\item \textbf{Eigenvalues:} $\lambda_1 > 0$ and $\lambda_2 > 0 \iff \lambda_1\lambda_2 = ac - b^2 > 0$ and $\lambda_1+\lambda_2 = a+c > 0 \iff a > 0 \text{ and } ac - b^2 > 0$.
\item \textbf{Pivots:} The first pivot is $a > 0$ and the second pivot is $c - \frac{b^2}{a} = \frac{ac - b^2}{a} > 0$.
\item \textbf{Energy $\vect{x}^\T A\vect{x}$ (completing the square):}
\begin{equation}
\vect{x}^\T A\vect{x} = \begin{bmatrix}x & y\end{bmatrix}\begin{bmatrix}a & b\\b & c\end{bmatrix}\begin{bmatrix}x\\y\end{bmatrix} = ax^2 + 2bxy + cy^2 = a\left(x + \frac{b}{a}y\right)^2 + \left(\frac{ac - b^2}{a}\right)y^2 > 0
\end{equation}
for all $(x,y) \neq (0,0)$ if and only if the two pivots $a$ and $(ac-b^2)/a$ are positive!
\end{itemize}

\begin{result}[ (Five Equivalent Statements of Positive Definiteness)]
When a symmetric matrix $A = A^\T$ has one of these five properties, it has them all:
\begin{enumerate}
\item All $n$ \term{pivots} are positive.
\item All $n$ \term{upper left determinants} ($1\times 1, 2\times 2, \dots, n\times n$) are positive.
\item All $n$ \term{eigenvalues} are positive.
\item $\vect{x}^\T A\vect{x}$ is positive except at $\vect{x} = \vect{0}$ (\term{energy-based definition}).
\item $A$ equals $R^\T R$ for a matrix $R$ with \term{independent columns} (including the Cholesky factor $R = \sqrt{D}L^\T$).
\end{enumerate}
If $A$ and $B$ are symmetric positive definite, so is $A + B$ because $\vect{x}^\T(A+B)\vect{x} = \vect{x}^\T A\vect{x} + \vect{x}^\T B\vect{x} > 0$.\\
For a \term{positive semidefinite} matrix, replace $>$ by $\ge$ ($\lambda_i \ge 0$, $\vect{x}^\T A\vect{x} \ge 0$, $A = R^\T R$ with possibly dependent columns).
\end{result}

\step{Step 2: Application to the tilted ellipse $\vect{x}^\T A\vect{x} = 1$.}
Diagonalize the positive definite matrix $A = Q\Lambda Q^\T$ and set $\vect{X} = \begin{bmatrix}X\\Y\end{bmatrix} = Q^\T\vect{x}$:
\begin{equation}
\vect{x}^\T A\vect{x} = \vect{x}^\T Q\Lambda Q^\T\vect{x} = \vect{X}^\T\Lambda\vect{X} = \lambda_1 X^2 + \lambda_2 Y^2 = 1.
\end{equation}
The axes of the tilted ellipse point along the \term{eigenvectors} of $A$ (columns of $Q$), and the half-lengths of the axes are $1/\sqrt{\lambda_1}$ and $1/\sqrt{\lambda_2}$.

\begin{example}{6.5.1--6.5.2 (p.\ 344--347)}{Testing Positive Definiteness and Finding Ellipse Axes}
\begin{enumerate}
\item \textbf{Example 1:} Test $A = \begin{bmatrix}2 & -1 & 0\\-1 & 2 & -1\\0 & -1 & 2\end{bmatrix}$ and $B = \begin{bmatrix}2 & -1 & b\\-1 & 2 & -1\\b & -1 & 2\end{bmatrix}$ for positive definiteness.
\item \textbf{Example 2:} Find the axes and half-lengths of the tilted ellipse $5x^2 + 8xy + 5y^2 = 1$.
\end{enumerate}
\end{example}

\vfill
\begin{pageend_sol}{6.5.1--6.5.2 (p.\ 344--347)}
\textbf{1.} For $A$: Pivots are $2, \frac{3}{2}, \frac{4}{3} > 0$; upper-left determinants are $2, 3, 4 > 0$; eigenvalues are $2-\sqrt{2}, 2, 2+\sqrt{2} > 0$; and $\vect{x}^\T A\vect{x} = 2(x_1 - \frac{1}{2}x_2)^2 + \frac{3}{2}(x_2 - \frac{2}{3}x_3)^2 + \frac{4}{3}x_3^2 > 0$. So $A$ is positive definite. For $B$: Upper-left $1\times 1$ and $2\times 2$ determinants are $2$ and $3$, and $\det B = 4 + 2b - 2b^2 = (1+b)(4-2b) > 0 \iff -1 < b < 2$.

\textbf{2.} $A = \begin{bmatrix}5 & 4\\4 & 5\end{bmatrix}$ has eigenvalues $\lambda_1 = 9$ (with $\vect{q}_1 = \frac{1}{\sqrt{2}}(1,1)$) and $\lambda_2 = 1$ (with $\vect{q}_2 = \frac{1}{\sqrt{2}}(1,-1)$). Thus $5x^2+8xy+5y^2 = 9\left(\frac{x+y}{\sqrt{2}}\right)^2 + 1\left(\frac{x-y}{\sqrt{2}}\right)^2 = 9X^2 + Y^2 = 1$. The shorter axis has half-length $1/\sqrt{9} = 1/3$ along $(1,1)/\sqrt{2}$; the longer axis has half-length $1/\sqrt{1} = 1$ along $(1,-1)/\sqrt{2}$.
\end{pageend_sol}
\clearpage

\section{Linear Algebra for Statistics and Probability\bk{8.6}}

\begin{goal}
Derive the weighted least squares normal equations, the mean and variance of Binomial, Poisson, and Normal distributions, the covariance matrix of the best estimate $\uv{\vect{x}}$, and Principal Component Analysis (PCA) via SVD.
\end{goal}

\step{Step 1: Weighted Least Squares.}
Multiply the $m$ equations $A\vect{x} = \vect{b}$ by a weight matrix $W$ to get $WA\vect{x} = W\vect{b}$. With $C = W^\T W$, the normal equations $(WA)^\T(WA)\uv{\vect{x}} = (WA)^\T W\vect{b}$ become:
\begin{equation}
\eqbox{A^\T C A\uv{\vect{x}} = A^\T C\vect{b}, \qquad \text{and for } n=1, A = \text{all ones:}\quad \uv{x}_W = \frac{w_1^2b_1 + \dots + w_m^2b_m}{w_1^2 + \dots + w_m^2}.}
\end{equation}
When independent errors $e_i$ have zero mean and variances $\sigma_i^2$, choose weights $w_i = 1/\sigma_i$ so all weighted equations have variance $1$. Then $C = W^\T W = \operatorname{diag}(1/\sigma_1^2, \dots, 1/\sigma_m^2) = \Sigma^{-1}$.

\step{Step 2: Mean, Variance, and Probability Distributions.}
\begin{align}
\text{\textbf{Discrete RV:}}\quad m &= \mathrm{E}[e] = \sum e_j p_j, & \sigma^2 &= \mathrm{E}[(e-m)^2] = \sum (e_j - m)^2 p_j \\
\text{\textbf{Continuous RV:}}\quad m &= \int_{-\infty}^\infty x p(x)\dd x, & \sigma^2 &= \int_{-\infty}^\infty (x-m)^2 p(x)\dd x \\
\text{\textbf{Binomial ($N$ fair coin flips):}}\quad p_M &= \frac{1}{2^N}\binom{N}{M} = \frac{1}{2^N}\frac{N!}{M!(N-M)!}, & m &= \frac{N}{2}, \quad \sigma^2 = \frac{N}{4} \\
\text{\textbf{Poisson ($\lambda = pN$, $p \ll 1$):}}\quad p_j &= \frac{\lambda^j}{j!}e^{-\lambda}, & m &= \lambda, \quad \sigma^2 = \lambda \\
\text{\textbf{Normal $\mathrm{N}(0,\sigma)$:}}\quad p(x) &= \frac{1}{\sqrt{2\pi}\,\sigma}e^{-x^2/2\sigma^2}, & \int_{-\infty}^\infty p(x)\dd x &= 1.
\end{align}

\step{Step 3: Covariance Matrix $\Sigma$, Gauss-Markov Estimate, and PCA.}
For $m$ zero-mean errors $\vect{e} = (e_1, \dots, e_m)$, the $(i,j)$ entry of the \term{covariance matrix} $\Sigma = \mathrm{E}[\vect{e}\vect{e}^\T]$ is $\sigma_{ij} = \mathrm{E}[e_i e_j]$:
\begin{align}
\text{\textbf{Multivariate Normal pdf:}}\quad p(\vect{b}) &= \frac{1}{(2\pi)^{m/2}|\Sigma|^{1/2}}e^{-\vect{b}^\T\Sigma^{-1}\vect{b}/2} \\
\text{\textbf{Best Linear Unbiased Estimate (BLUE):}}\quad A^\T\Sigma^{-1}A\uv{\vect{x}} &= A^\T\Sigma^{-1}\vect{b} \implies \uv{\vect{x}} = (A^\T\Sigma^{-1}A)^{-1}A^\T\Sigma^{-1}\vect{b} \\
\text{\textbf{Covariance of output error $\vect{x}-\uv{\vect{x}}$:}}\quad P &= \mathrm{E}[(\vect{x}-\uv{\vect{x}})(\vect{x}-\uv{\vect{x}})^\T] = (A^\T\Sigma^{-1}A)^{-1}.
\end{align}
\textbf{Principal Component Analysis (PCA):} Given an $m\times n$ mean-centered data matrix $A = U\Sigma_{\text{SVD}} V^\T$ with singular values $\sigma_1 \ge \sigma_2 \ge \dots \ge \sigma_r$, the leading singular vectors $\vect{u}_1$ and $\vect{v}_1$ give the combination with greatest variance. The fraction of total variance explained by the first component is $\sigma_1^2 / (\sigma_1^2 + \dots + \sigma_r^2)$.

\begin{example}{8.6.1--8.6.6 (p.\ 454--457)}{Coin Flips, Variance Reduction by Averaging, and PCA}
\begin{enumerate}
\item \textbf{Examples 1--2:} Find the mean and variance for $1$ flip and $3$ flips of a fair coin.
\item \textbf{Example 5:} Your pulse rate is measured $10$ times by independent doctors, each with variance $\sigma^2$ ($\Sigma = \sigma^2 I$). Find the variance $P$ of the average $\uv{x}$.
\item \textbf{Example 6:} In the grade matrix $A = U\Sigma V^\T$ with singular values $\sigma_1=3, \sigma_2=2, \sigma_3=1$ (or scaled $12, 8, 4$), what fraction of the variance is explained by the first eigencourse $\vect{u}_1$ and eigenstudent $\vect{v}_1$?
\end{enumerate}
\end{example}

\vfill
\begin{pageend_sol}{8.6.1--8.6.6 (p.\ 454--457)}
\textbf{1.} For 1 flip ($0$ or $1$ with $p_0=p_1=1/2$): $m = 0(\frac{1}{2})+1(\frac{1}{2}) = \frac{1}{2}$, and $\sigma^2 = (0-\frac{1}{2})^2\frac{1}{2} + (1-\frac{1}{2})^2\frac{1}{2} = \frac{1}{4}$. For $N=3$ flips: $p_0=\frac{1}{8}, p_1=\frac{3}{8}, p_2=\frac{3}{8}, p_3=\frac{1}{8}$, $m = \frac{3}{2}$, and $\sigma^2 = \frac{3}{4}$.

\textbf{2.} $A$ is the $10\times 1$ column of ones and $\Sigma^{-1} = \frac{1}{\sigma^2}I$. Thus $P = (A^\T\Sigma^{-1}A)^{-1} = \left(\frac{10}{\sigma^2}\right)^{-1} = \frac{\sigma^2}{10}$. Averaging $10$ independent measurements divides the variance by $10$.

\textbf{3.} The fraction of variance explained is $\frac{\sigma_1^2}{\sigma_1^2 + \sigma_2^2 + \sigma_3^2} = \frac{3^2}{3^2 + 2^2 + 1^2} = \frac{9}{9 + 4 + 1} = \frac{9}{14}$.
\end{pageend_sol}
\clearpage

% ============================================================================
% END OF MODULE 3: ALL UNIQUE SOLVED PROBLEMS (WORKED EXAMPLES)
% ============================================================================
\chapter[Module 3: Unique Solved Problems]{Module 3: All Unique Solved Problems (Worked Examples)}

\begin{syll}
All unique Worked Examples from Module 3 (Sections 1.2, 4.1--4.4, and 6.5), with problem statements copied verbatim and full step-by-step solutions.
\end{syll}

\begin{example}{1.2 A}{Testing Schwarz and Triangle Inequalities}
For the vectors $\vect{v} = (3, 4)$ and $\vect{w} = (4, 3)$ test the Schwarz inequality on $\vect{v}\cdot\vect{w}$ and the triangle inequality on $\|\vect{v} + \vect{w}\|$. Find $\cos\theta$ for the angle between $\vect{v}$ and $\vect{w}$. When will we have \emph{equality} $|\vect{v}\cdot\vect{w}| = \|\vect{v}\|\|\vect{w}\|$ and $\|\vect{v} + \vect{w}\| = \|\vect{v}\| + \|\vect{w}\|$?
\end{example}
\textbf{Solution:} $\vect{v}\cdot\vect{w} = (3)(4) + (4)(3) = 24$, $\|\vect{v}\| = \sqrt{9+16} = 5$, $\|\vect{w}\| = 5$, and $\vect{v}+\vect{w} = (7,7)$ has length $7\sqrt{2} < 10$:
\begin{itemize}
\item \textbf{Schwarz inequality:} $|\vect{v}\cdot\vect{w}| \le \|\vect{v}\|\|\vect{w}\|$ is $24 < 25$.
\item \textbf{Triangle inequality:} $\|\vect{v}+\vect{w}\| \le \|\vect{v}\| + \|\vect{w}\|$ is $7\sqrt{2} < 5 + 5$.
\item \textbf{Cosine of angle:} $\cos\theta = \frac{24}{25}$.
\end{itemize}
Equality $|\vect{v}\cdot\vect{w}| = \|\vect{v}\|\|\vect{w}\|$ holds when one vector is a multiple of the other ($\vect{w} = c\vect{v}$, angle $0^\circ$ or $180^\circ$). Equality $\|\vect{v}+\vect{w}\| = \|\vect{v}\| + \|\vect{w}\|$ holds when the angle is $0^\circ$ ($\vect{w} = c\vect{v}$ with $c \ge 0$).

\begin{example}{1.2 B}{Unit Vectors and Perpendicular Unit Vectors in $\mathbf{R}^2$}
Find a unit vector $\vect{u}$ in the direction of $\vect{v} = (3, 4)$. Find a unit vector $\vect{U}$ that is perpendicular to $\vect{u}$. How many possibilities for $\vect{U}$?
\end{example}
\textbf{Solution:} Divide $\vect{v}$ by $\|\vect{v}\| = 5$ to get $\vect{u} = \left(\frac{3}{5}, \frac{4}{5}\right)$. A perpendicular vector is $\vect{V} = (-4,3)$ since $(3)(-4)+(4)(3)=0$. Divide $\vect{V}$ by $\|\vect{V}\|=5$ to get $\vect{U} = \left(-\frac{4}{5}, \frac{3}{5}\right)$. There are \textbf{two} possibilities for $\vect{U}$ in $\mathbf{R}^2$: $\vect{U} = \left(-\frac{4}{5}, \frac{3}{5}\right)$ and $-\vect{U} = \left(\frac{4}{5}, -\frac{3}{5}\right)$.

\begin{example}{1.2 C}{Finding a Vector from Given Dot Products}
Find a vector $\vect{x} = (c, d)$ that has dot products $\vect{x}\cdot\vect{r} = 1$ and $\vect{x}\cdot\vect{s} = 0$ with the given vectors $\vect{r} = (2, -1)$ and $\vect{s} = (-1, 2)$. How is this question related to Example 1.1 C, which solved $c\vect{v} + d\vect{w} = \vect{b} = (1, 0)$?
\end{example}
\textbf{Solution:} The two dot products give the exact same linear equations as Worked Example 1.1 C:
\[
\begin{aligned}
\vect{x}\cdot\vect{r} &= 2c - d = 1 \\
\vect{x}\cdot\vect{s} &= -c + 2d = 0
\end{aligned}
\implies \vect{x} = (c,d) = \left(\frac{2}{3}, \frac{1}{3}\right).
\]
In 1.1 C, $\vect{v}$ and $\vect{w}$ were the columns of $A$; here $\vect{r}$ and $\vect{s}$ are the rows of $A$ in $A\vect{x} = \vect{b}$.

\begin{example}{4.1 A}{Dimensions of Orthogonal Subspaces and Complements}
Suppose $\mathbf{S}$ is a six-dimensional subspace of nine-dimensional space $\mathbf{R}^9$.
(a) What are the possible dimensions of subspaces orthogonal to $\mathbf{S}$?
(b) What are the possible dimensions of the orthogonal complement $\mathbf{S}^\perp$ of $\mathbf{S}$?
(c) What is the smallest possible size of a matrix $A$ that has row space $\mathbf{S}$?
(d) What is the shape of its nullspace matrix $N$?
\end{example}
\textbf{Solution:}
(a) Subspaces orthogonal to $\mathbf{S}$ can have dimensions $0, 1, 2, 3$.
(b) The orthogonal complement $\mathbf{S}^\perp$ is the largest orthogonal subspace, with dimension $9 - 6 = 3$.
(c) The smallest matrix $A$ with row space $\mathbf{S}$ is $6$ by $9$ (its $6$ rows are a basis for $\mathbf{S}$).
(d) Its nullspace matrix $N$ is $9$ by $3$ (its $3$ columns are a basis for $\mathbf{S}^\perp = N(A)$).

\begin{example}{4.1 B}{Splitting a Vector into Nullspace and Row Space Components}
The equation $x - 3y - 4z = 0$ describes a plane $\mathbf{P}$ in $\mathbf{R}^3$ (actually a subspace).
(a) The plane $\mathbf{P}$ is the nullspace $N(A)$ of what $1$ by $3$ matrix $A$?
(b) Find a basis $\vect{s}_1, \vect{s}_2$ of special solutions of $x - 3y - 4z = 0$.
(c) Also find a basis for the line $\mathbf{P}^\perp$ that is perpendicular to $\mathbf{P}$.
(d) Split $\vect{v} = (6, 4, 5)$ into its nullspace component $\vect{v}_n$ in $\mathbf{P}$ and its row space component $\vect{v}_r$ in $\mathbf{P}^\perp$.
\end{example}
\textbf{Solution:}
(a) $A = \begin{bmatrix}1 & -3 & -4\end{bmatrix}$.
(b) Set free variables $(y,z)=(1,0)$ and $(0,1)$ to get $\vect{s}_1 = (3,1,0)$ and $\vect{s}_2 = (4,0,1)$.
(c) The row space $\mathbf{P}^\perp = C(A^\T)$ is the line spanned by $\vect{z} = (1,-3,-4)$.
(d) Solve $c_1\vect{s}_1 + c_2\vect{s}_2 + c_3\vect{z} = \vect{v}$:
\[
\begin{bmatrix} 3 & 4 & 1 \\ 1 & 0 & -3 \\ 0 & 1 & -4 \end{bmatrix} \begin{bmatrix} c_1 \\ c_2 \\ c_3 \end{bmatrix} = \begin{bmatrix} 6 \\ 4 \\ 5 \end{bmatrix} \implies \begin{bmatrix} c_1 \\ c_2 \\ c_3 \end{bmatrix} = \begin{bmatrix} 1 \\ 1 \\ -1 \end{bmatrix} \implies \vect{v}_n = \vect{s}_1 + \vect{s}_2 = (7,1,1), \quad \vect{v}_r = -\vect{z} = (-1,3,4).
\]

\begin{example}{4.2 A}{Projecting Onto a Line and Onto a Plane}
Project the vector $\vect{b} = (3, 4, 4)$ onto the line through $\vect{a} = (2, 2, 1)$ and then onto the plane that also contains $\vect{a}^* = (1, 0, 0)$. Check that the first error vector $\vect{b} - \vect{p}$ is perpendicular to $\vect{a}$, and the second error vector $\vect{e}^* = \vect{b} - \vect{p}^*$ is also perpendicular to $\vect{a}^*$. Find the $3$ by $3$ projection matrix $P$ onto that plane of $\vect{a}$ and $\vect{a}^*$.
\end{example}
\textbf{Solution:}
\begin{itemize}
\item \textbf{Onto the line:} $\vect{p} = \frac{\vect{a}^\T\vect{b}}{\vect{a}^\T\vect{a}}\vect{a} = \frac{18}{9}(2,2,1) = (4,4,2)$. Error $\vect{e} = \vect{b}-\vect{p} = (-1,0,2)$ has $\vect{e}^\T\vect{a} = -2+0+2 = 0$.
\item \textbf{Onto the plane:} With $A = \begin{bmatrix}\vect{a} & \vect{a}^*\end{bmatrix} = \begin{bmatrix}2&1\\2&0\\1&0\end{bmatrix}$, we have $A^\T A = \begin{bmatrix}9&2\\2&1\end{bmatrix}$ and $(A^\T A)^{-1} = \frac{1}{5}\begin{bmatrix}1&-2\\-2&9\end{bmatrix}$:
\[
P = A(A^\T A)^{-1}A^\T = \begin{bmatrix} 1 & 0 & 0 \\ 0 & .8 & .4 \\ 0 & .4 & .2 \end{bmatrix}, \qquad \vect{p}^* = P\vect{b} = (3, 4.8, 2.4), \qquad \vect{e}^* = \vect{b} - \vect{p}^* = (0, -0.8, 1.6).
\]
Notice $\vect{e}^*$ is perpendicular to both $\vect{a}$ and $\vect{a}^*$, and $P\vect{e}^* = \vect{0}$.
\end{itemize}

\begin{example}{4.2 B and 4.2 C}{Least Squares Average and Recursive Least Squares (Kalman Update)}
\textbf{4.2 B:} Suppose your pulse is measured at $x = 70, 80, 120$. Minimize $E = (x-70)^2 + (x-80)^2 + (x-120)^2$ by calculus ($\diff E/\diff x = 0$) and by projecting $\vect{b} = (70,80,120)$ onto $\vect{a} = (1,1,1)$.\\
\textbf{4.2 C:} In recursive least squares, a fourth measurement $b_4 = 130$ changes $\uv{x}_{\text{old}}$ to $\uv{x}_{\text{new}}$. Verify the update formula $\uv{x}_{\text{new}} = \uv{x}_{\text{old}} + \frac{1}{4}(130 - \uv{x}_{\text{old}})$.
\end{example}
\textbf{Solution:}
\textbf{4.2 B:} $\frac{\diff E}{\diff x} = 2(x-70) + 2(x-80) + 2(x-120) = 0 \implies \uv{x} = \frac{70+80+120}{3} = 90$. By projection, $\uv{x} = \frac{\vect{a}^\T\vect{b}}{\vect{a}^\T\vect{a}} = \frac{270}{3} = 90$.
\textbf{4.2 C:} $\uv{x}_{\text{new}} = \frac{70+80+120+130}{4} = 100$, which matches $\uv{x}_{\text{old}} + \frac{1}{4}(b_4 - \uv{x}_{\text{old}}) = 90 + \frac{1}{4}(130 - 90) = 90 + 10 = 100$. In general, going from $k-1$ to $k$ measurements gives $\uv{x}_{\text{new}} = \uv{x}_{\text{old}} + \frac{1}{k}(b_k - \uv{x}_{\text{old}})$.

\begin{example}{4.3 A}{Comparing $E_2$ (Least Squares), $E_\infty$ (Minimax), and $E_1$ (Median) on an Outlier}
Start with nine measurements $b_1$ to $b_9$, all zero, at times $t = 1, \dots, 9$. The tenth measurement $b_{10} = 40$ is an outlier. Find the best horizontal line $y = C$ to fit the ten points using three measures for the error $E$:
(1) Least squares $E_2 = e_1^2 + \dots + e_{10}^2$ \quad (2) Least maximum error $E_\infty = |e_{\max}|$ \quad (3) Least sum of errors $E_1 = |e_1| + \dots + |e_{10}|$.
Then find the least squares straight line $C + Dt$ through those ten points.
\end{example}
\textbf{Solution:}
\begin{enumerate}
\item[(1)] \textbf{Least squares ($E_2$):} $C = \text{average of } b_i = \frac{40}{10} = 4$.
\item[(2)] \textbf{Least maximum error ($E_\infty$):} $C = 20$ (halfway between $0$ and $40$).
\item[(3)] \textbf{Least sum of absolute errors ($E_1$):} $C = \text{median} = 0$ (since $9|C| + |40-C|$ increases if $C$ moves away from $0$).
\end{enumerate}
For the least squares line $C + Dt$:
\[
A^\T A = \begin{bmatrix} 10 & \sum_{t=1}^{10} t \\ \sum_{t=1}^{10} t & \sum_{t=1}^{10} t^2 \end{bmatrix} = \begin{bmatrix} 10 & 55 \\ 55 & 385 \end{bmatrix}, \quad A^\T\vect{b} = \begin{bmatrix} 40 \\ 400 \end{bmatrix} \implies C = -8, \quad D = \frac{24}{11}.
\]

\begin{example}{4.3 B}{Uncoupled Parabola Fit with Symmetric Times}
Find the parabola $C + Dt + Et^2$ that comes closest (least squares error) to the values $\vect{b} = (0, 0, 1, 0, 0)$ at the times $t = -2, -1, 0, 1, 2$.
\end{example}
\textbf{Solution:} Because the times $-2,-1,0,1,2$ are symmetric around $t=0$, column 2 of $A$ is orthogonal to columns 1 and 3:
\[
A = \begin{bmatrix} 1 & -2 & 4 \\ 1 & -1 & 1 \\ 1 & 0 & 0 \\ 1 & 1 & 1 \\ 1 & 2 & 4 \end{bmatrix} \implies A^\T A\uv{\vect{x}} = \begin{bmatrix} 5 & 0 & 10 \\ 0 & 10 & 0 \\ 10 & 0 & 34 \end{bmatrix} \begin{bmatrix} C \\ D \\ E \end{bmatrix} = \begin{bmatrix} 1 \\ 0 \\ 0 \end{bmatrix} \implies \begin{aligned} C &= 34/70 \\ D &= 0 \\ E &= -10/70. \end{aligned}
\]

\begin{example}{4.4 A}{Hadamard Matrices and Orthogonal Projections}
Add two more columns with all entries $1$ or $-1$, so the columns of $H_4$ are orthogonal. How do you turn $H_4$ into an orthogonal matrix $Q$? What is $H_8^\T H_8$? If the projection of $\vect{b} = (6, 0, 0, 2)$ onto column 1 of $H_4$ is $\vect{p}_1 = (2, 2, 2, 2)$ and onto column 2 is $\vect{p}_2 = (1, -1, 1, -1)$, what is the projection $\vect{p}_{1,2}$ onto the 2D space spanned by the first two columns?
\end{example}
\textbf{Solution:}
\[
H_2 = \begin{bmatrix} 1 & 1 \\ 1 & -1 \end{bmatrix}, \qquad H_4 = \begin{bmatrix} H_2 & H_2 \\ H_2 & -H_2 \end{bmatrix} = \begin{bmatrix} 1 & 1 & 1 & 1 \\ 1 & -1 & 1 & -1 \\ 1 & 1 & -1 & -1 \\ 1 & -1 & -1 & 1 \end{bmatrix}, \qquad Q_4 = \frac{1}{2}H_4.
\]
For $H_8 = \begin{bmatrix}H_4 & H_4\\H_4 & -H_4\end{bmatrix}$, we have $H_8^\T H_8 = 8I_8$ and $Q_8 = H_8/\sqrt{8}$. Because the columns of $H_4$ are orthogonal, the 1D projections simply add without any matrix inversion:
\[
\vect{p}_{1,2} = \vect{p}_1 + \vect{p}_2 = (2, 2, 2, 2) + (1, -1, 1, -1) = (3, 1, 3, 1).
\]

\begin{example}{6.5 A and 6.5 B}{Energy Proofs $A = LDL^\T = Q\Lambda Q^\T$ and Block Positive Definiteness}
\textbf{6.5 A:} Show from $A = LDL^\T$ and $A = Q\Lambda Q^\T$ that $\vect{x}^\T A\vect{x} > 0$ for all $\vect{x} \neq \vect{0}$ exactly when the pivots and eigenvalues are positive.\\
\textbf{6.5 B:} When is the symmetric block matrix $M = \begin{bmatrix}A & B\\B^\T & C\end{bmatrix}$ positive definite?
\end{example}
\textbf{Solution:}
\textbf{6.5 A:} Put parentheses into $\vect{x}^\T LDL^\T\vect{x} = (L^\T\vect{x})^\T D(L^\T\vect{x}) = \vect{y}^\T D\vect{y} = d_1y_1^2 + \dots + d_ny_n^2 > 0$ and $\vect{x}^\T Q\Lambda Q^\T\vect{x} = (Q^\T\vect{x})^\T\Lambda(Q^\T\vect{x}) = \vect{z}^\T\Lambda\vect{z} = \lambda_1z_1^2 + \dots + \lambda_nz_n^2 > 0$. Since $L^\T$ and $Q^\T$ are invertible, $\vect{x} \neq \vect{0} \iff \vect{y} \neq \vect{0} \iff \vect{z} \neq \vect{0}$.
\textbf{6.5 B:} Block elimination gives:
\[
\begin{bmatrix} I & \vect{0} \\ -B^\T A^{-1} & I \end{bmatrix} \begin{bmatrix} A & B \\ B^\T & C \end{bmatrix} = \begin{bmatrix} A & B \\ \vect{0} & C - B^\T A^{-1}B \end{bmatrix} = \begin{bmatrix} A & B \\ \vect{0} & S \end{bmatrix}.
\]
Therefore $M$ is positive definite if and only if both the upper-left block $A$ and the Schur complement $S = C - B^\T A^{-1}B$ are positive definite.

\begin{example}{6.5 C and 6.5 D}{Second-Derivative Minimum Test and Sine Eigenvectors of $K_n$}
\textbf{6.5 C:} Does $F(x,y)$ have a minimum if $\partial F/\partial x = 0$ and $\partial F/\partial y = 0$ at $(x,y) = (0,0)$?\\
\textbf{6.5 D:} Find the eigenvalues and eigenvectors of the $n\times n$ tridiagonal $-1, 2, -1$ matrix $K$.
\end{example}
\textbf{Solution:}
\textbf{6.5 C:} $F(x,y)$ has a strict local minimum if the symmetric Hessian matrix of second derivatives is positive definite:
\[
H = \begin{bmatrix} \partial^2 F/\partial x^2 & \partial^2 F/\partial x\partial y \\ \partial^2 F/\partial y\partial x & \partial^2 F/\partial y^2 \end{bmatrix} = \begin{bmatrix} 2a & 2b \\ 2b & 2c \end{bmatrix} \quad \text{is positive definite} \quad (a > 0 \text{ and } ac - b^2 > 0).
\]
\textbf{6.5 D:} With step size $h = \frac{1}{n+1}$, the eigenvectors of $K$ are the orthogonal discrete sine vectors $\vect{s}_k = (\sin k\pi h, \sin 2k\pi h, \dots, \sin nk\pi h)$ for $k = 1, \dots, n$, and the eigenvalues are:
\[
\lambda_k = 2 - 2\cos k\pi h = 2 - 2\cos\frac{k\pi}{n+1} > 0.
\]

% ============================================================================
% MODULE 4: EIGENVALUES, SVD, MARKOV, AND COMPLEX MATRICES (PDF 3)
% ============================================================================
\part{Eigenvalues, SVD, Markov Matrices, and Complex Matrices}

\chapter[Eigenvalues, Diagonalization, Symmetric Matrices, SVD (6.1--6.7)]{Eigenvalues, Diagonalization, Symmetric Matrices, and SVD\bk{6.1, 6.2, 6.4, 6.6, 6.7}}

\begin{syll}
Syllabus topics: Eigenvalue equation $A\vect{x} = \lambda\vect{x}$ and characteristic equation $\det(A-\lambda I) = 0$ $\cdot$ Trace $\sum \lambda_i = \operatorname{trace}(A)$ and determinant $\prod \lambda_i = \det A$ $\cdot$ Diagonalization $A = S\Lambda S^{-1}$, powers $A^k = S\Lambda^k S^{-1}$, Fibonacci numbers, and difference equations $\vect{u}_{k+1} = A\vect{u}_k$ $\cdot$ Geometric vs.\ algebraic multiplicity (GM $\le$ AM) $\cdot$ Symmetric matrices and the Spectral Theorem $A = Q\Lambda Q^\T$ $\cdot$ Similar matrices $B = M^{-1}AM$ and Jordan form $J$ $\cdot$ Singular Value Decomposition $A = U\Sigma V^\T$.
\end{syll}

\section{Introduction to Eigenvalues and Eigenvectors\bk{6.1}}

\begin{goal}
Derive the characteristic equation $\det(A-\lambda I)=0$, the eigenvalues of $A^2, A^{-1}, A+cI$, and the trace and determinant checks.
\end{goal}

\step{Step 1: Move $\lambda\vect{x}$ to the left side of $A\vect{x} = \lambda\vect{x}$.}
An eigenvector $\vect{x} \neq \vect{0}$ keeps its direction (or reverses) when multiplied by $A$:
\begin{equation}
A\vect{x} = \lambda\vect{x} \iff (A - \lambda I)\vect{x} = \vect{0}.
\end{equation}
For $(A - \lambda I)\vect{x} = \vect{0}$ to have a nonzero solution $\vect{x}$ in the nullspace of $A - \lambda I$, the matrix $A - \lambda I$ must be singular:
\begin{equation}
\eqbox{\det(A - \lambda I) = 0.}
\end{equation}

\step{Step 2: Derive the eigenvalues of $A^2$, $A^{-1}$, and $A + cI$, and the trace/determinant rules.}
Start from $A\vect{x} = \lambda\vect{x}$:
\begin{align}
\text{\textbf{Square $A$:}}\quad A^2\vect{x} &= A(A\vect{x}) = A(\lambda\vect{x}) = \lambda(A\vect{x}) = \lambda^2\vect{x} &&\why{eigenvalue $\lambda^2$, same $\vect{x}$} \\
\text{\textbf{Invert $A$:}}\quad \vect{x} &= A^{-1}(\lambda\vect{x}) \implies A^{-1}\vect{x} = \lambda^{-1}\vect{x} &&\why{eigenvalue $1/\lambda$, same $\vect{x}$} \\
\text{\textbf{Shift $A$:}}\quad (A + cI)\vect{x} &= A\vect{x} + c\vect{x} = (\lambda + c)\vect{x} &&\why{eigenvalue $\lambda + c$, same $\vect{x}$}
\end{align}
For a $2\times 2$ matrix $A = \begin{bmatrix}a & b\\c & d\end{bmatrix}$, expand $\det(A - \lambda I) = (a-\lambda)(d-\lambda) - bc = \lambda^2 - (a+d)\lambda + (ad-bc) = (\lambda-\lambda_1)(\lambda-\lambda_2)$. For any $n\times n$ matrix:
\begin{result}[ (Trace and Determinant Eigenvalue Checks)]
\begin{align}
\lambda_1 + \lambda_2 + \dots + \lambda_n &= \operatorname{trace}(A) = a_{11} + a_{22} + \dots + a_{nn} &&\why{sum of eigenvalues = trace} \\
\lambda_1 \lambda_2 \cdots \lambda_n &= \det A &&\why{product of eigenvalues = determinant}
\end{align}
\end{result}

\begin{example}{6.1.1--6.1.5 (p.\ 284--289)}{Markov, Projection, Reflection, Singular, and Rotation Eigenvalues}
\begin{enumerate}
\item \textbf{Example 1:} Find the eigenvalues and eigenvectors of $A = \begin{bmatrix}.8 & .3\\.2 & .7\end{bmatrix}$ and compute $A^{99}\begin{bmatrix}.8\\.2\end{bmatrix}$.
\item \textbf{Examples 2--3:} Find the eigenvalues and eigenvectors of the projection $P = \begin{bmatrix}.5 & .5\\.5 & .5\end{bmatrix}$ and reflection $R = 2P - I = \begin{bmatrix}0 & 1\\1 & 0\end{bmatrix}$.
\item \textbf{Example 4:} Find $\lambda$'s and $\vect{x}$'s for the singular matrix $A = \begin{bmatrix}1 & 2\\2 & 4\end{bmatrix}$.
\item \textbf{Example 5:} Find $\lambda$'s and $\vect{x}$'s for the $90^\circ$ rotation matrix $Q = \begin{bmatrix}0 & -1\\1 & 0\end{bmatrix}$.
\end{enumerate}
\end{example}

\vfill
\begin{pageend_sol}{6.1.1--6.1.5 (p.\ 284--289)}
\textbf{1.} $\det(A-\lambda I) = \lambda^2 - \frac{3}{2}\lambda + \frac{1}{2} = (\lambda-1)(\lambda-\frac{1}{2}) = 0 \implies \lambda_1=1, \lambda_2=0.5$ with $\vect{x}_1=(.6,.4)$ and $\vect{x}_2=(1,-1)$. Since $\begin{bmatrix}.8\\.2\end{bmatrix} = \vect{x}_1 + 0.2\vect{x}_2$, we have $A^{99}\begin{bmatrix}.8\\.2\end{bmatrix} = \vect{x}_1 + 0.2\left(\frac{1}{2}\right)^{99}\vect{x}_2 \approx \begin{bmatrix}.6\\.4\end{bmatrix}$.

\textbf{2.} $P$ has $\lambda_1 = 1$ ($\vect{x}_1=(1,1)$) and $\lambda_2 = 0$ ($\vect{x}_2=(1,-1)$). $R = 2P - I$ has the same eigenvectors and shifted eigenvalues $\lambda = 2(1)-1 = 1$ and $2(0)-1 = -1$.

\textbf{3.} $\det(A-\lambda I) = \lambda^2 - 5\lambda = 0 \implies \lambda_1 = 0$ ($\vect{x}_1 = (2,-1)$) and $\lambda_2 = 5$ ($\vect{x}_2 = (1,2)$).

\textbf{4.} $\det(Q-\lambda I) = \lambda^2 + 1 = 0 \implies \lambda_1 = -i$ ($\vect{x}_1 = (1,i)$) and $\lambda_2 = i$ ($\vect{x}_2 = (i,1)$).
\end{pageend_sol}
\clearpage

\section{Diagonalizing a Matrix and Powers $A^k$\bk{6.2}}

\begin{goal}
Diagonalize $A = S\Lambda S^{-1}$ using $n$ independent eigenvectors, prove that distinct eigenvalues guarantee independent eigenvectors, and solve $\vect{u}_{k+1} = A\vect{u}_k$ (including Fibonacci numbers).
\end{goal}

\step{Step 1: Derive $AS = S\Lambda$ and $A = S\Lambda S^{-1}$.}
Put $n$ linearly independent eigenvectors $\vect{x}_1, \dots, \vect{x}_n$ of $A$ into the columns of the \term{eigenvector matrix} $S$:
\begin{equation}
AS = A\begin{bmatrix} \vect{x}_1 & \dots & \vect{x}_n \end{bmatrix} = \begin{bmatrix} \lambda_1\vect{x}_1 & \dots & \lambda_n\vect{x}_n \end{bmatrix} = \begin{bmatrix} \vect{x}_1 & \dots & \vect{x}_n \end{bmatrix} \begin{bmatrix} \lambda_1 & & \\ & \ddots & \\ & & \lambda_n \end{bmatrix} = S\Lambda.
\end{equation}
Because the $n$ eigenvectors are independent, $S$ is invertible:
\begin{equation}
\eqbox{S^{-1}AS = \Lambda \qquad \text{or} \qquad A = S\Lambda S^{-1} \qquad \text{and} \qquad A^k = S\Lambda^k S^{-1}.}
\end{equation}

\step{Step 2: Prove that eigenvectors for distinct eigenvalues $\lambda_1 \neq \lambda_2$ are linearly independent.}
Suppose $c_1\vect{x}_1 + c_2\vect{x}_2 = \vect{0}$. Multiply by $A$ and also multiply by $\lambda_2$:
\begin{equation}
\left.\begin{aligned}
c_1\lambda_1\vect{x}_1 + c_2\lambda_2\vect{x}_2 &= \vect{0} \\
c_1\lambda_2\vect{x}_1 + c_2\lambda_2\vect{x}_2 &= \vect{0}
\end{aligned}\right\} \implies (\lambda_1 - \lambda_2)c_1\vect{x}_1 = \vect{0} \implies c_1 = 0 \quad (\text{since } \lambda_1 \neq \lambda_2, \vect{x}_1 \neq \vect{0}).
\end{equation}
Then $c_2\vect{x}_2 = \vect{0} \implies c_2 = 0$. Therefore $\vect{x}_1$ and $\vect{x}_2$ are independent. (For $j$ distinct eigenvalues, multiplying by $(A-\lambda_2 I)\cdots(A-\lambda_j I)$ leaves $(\lambda_1-\lambda_2)\cdots(\lambda_1-\lambda_j)c_1\vect{x}_1 = \vect{0} \implies c_1 = 0$.)

\step{Step 3: Solve $\vect{u}_{k+1} = A\vect{u}_k$ in three steps and derive Binet's Fibonacci formula.}
\begin{result}[ (Solution to $\vect{u}_{k+1} = A\vect{u}_k$ and Fibonacci Numbers $F_k$)]
\begin{enumerate}
\item \textbf{Expand $\vect{u}_0$ in eigenvectors:} $\vect{u}_0 = c_1\vect{x}_1 + \dots + c_n\vect{x}_n = S\vect{c} \implies \vect{c} = S^{-1}\vect{u}_0$.
\item \textbf{Multiply each $\vect{x}_i$ by $\lambda_i^k$:} $\Lambda^k\vect{c} = \Lambda^k S^{-1}\vect{u}_0$.
\item \textbf{Combine into $\vect{u}_k = A^k\vect{u}_0$:} $\eqbox{\vect{u}_k = A^k\vect{u}_0 = S\Lambda^k S^{-1}\vect{u}_0 = c_1\lambda_1^k\vect{x}_1 + \dots + c_n\lambda_n^k\vect{x}_n.}$
\end{enumerate}
\textbf{Fibonacci Numbers ($F_{k+2} = F_{k+1} + F_k$ with $F_0=0, F_1=1$):}\\
Set $\vect{u}_k = \begin{bmatrix}F_{k+1}\\F_k\end{bmatrix}$ so $\vect{u}_{k+1} = \begin{bmatrix}1 & 1\\1 & 0\end{bmatrix}\vect{u}_k$. From $\lambda^2 - \lambda - 1 = 0$, the eigenvalues are $\lambda_1 = \frac{1+\sqrt{5}}{2}$ and $\lambda_2 = \frac{1-\sqrt{5}}{2}$ with $\vect{x}_i = \begin{bmatrix}\lambda_i\\1\end{bmatrix}$ and $\vect{u}_0 = \begin{bmatrix}1\\0\end{bmatrix} = \frac{\vect{x}_1 - \vect{x}_2}{\lambda_1 - \lambda_2}$. The second component of $\vect{u}_k$ is:
\begin{equation}
F_k = \frac{\lambda_1^k - \lambda_2^k}{\lambda_1 - \lambda_2} = \frac{1}{\sqrt{5}}\left[\left(\frac{1+\sqrt{5}}{2}\right)^k - \left(\frac{1-\sqrt{5}}{2}\right)^k\right].
\end{equation}
\end{result}

\begin{example}{6.2.1--6.2.3 (p.\ 299--303)}{Powers $A^k$ and Difference Equations $\vect{u}_{k+1} = A\vect{u}_k$}
\begin{enumerate}
\item \textbf{Example 1:} Diagonalize $A = \begin{bmatrix}1 & 5\\0 & 6\end{bmatrix}$ and compute $A^k = S\Lambda^k S^{-1}$.
\item \textbf{Example 3:} Start from $\vect{u}_0 = (1,0)$ and compute $\vect{u}_k = A^k\vect{u}_0$ for $A = \begin{bmatrix}1 & 2\\1 & 0\end{bmatrix}$ ($F_{k+2} = F_{k+1} + 2F_k$).
\end{enumerate}
\end{example}

\vfill
\begin{pageend_sol}{6.2.1--6.2.3 (p.\ 299--303)}
\textbf{1.} $\lambda_1=1$ has $\vect{x}_1 = \begin{bmatrix}1\\0\end{bmatrix}$ and $\lambda_2=6$ has $\vect{x}_2 = \begin{bmatrix}1\\1\end{bmatrix}$. Thus $S = \begin{bmatrix}1 & 1\\0 & 1\end{bmatrix}$ and $S^{-1} = \begin{bmatrix}1 & -1\\0 & 1\end{bmatrix}$:
\[
A^k = \begin{bmatrix}1 & 1\\0 & 1\end{bmatrix}\begin{bmatrix}1 & 0\\0 & 6^k\end{bmatrix}\begin{bmatrix}1 & -1\\0 & 1\end{bmatrix} = \begin{bmatrix}1 & 6^k - 1\\0 & 6^k\end{bmatrix}.
\]
\textbf{2.} $\det(A-\lambda I) = \lambda^2 - \lambda - 2 = (\lambda-2)(\lambda+1) = 0 \implies \lambda_1=2, \vect{x}_1=\begin{bmatrix}2\\1\end{bmatrix}$ and $\lambda_2=-1, \vect{x}_2=\begin{bmatrix}1\\-1\end{bmatrix}$. Step 1: $\vect{u}_0 = \begin{bmatrix}1\\0\end{bmatrix} = \frac{1}{3}\begin{bmatrix}2\\1\end{bmatrix} + \frac{1}{3}\begin{bmatrix}1\\-1\end{bmatrix}$ ($c_1=c_2=\frac{1}{3}$). Steps 2--3: $\vect{u}_k = \frac{1}{3}2^k\begin{bmatrix}2\\1\end{bmatrix} + \frac{1}{3}(-1)^k\begin{bmatrix}1\\-1\end{bmatrix}$, so $F_k = \frac{2^k - (-1)^k}{3}$.
\end{pageend_sol}
\clearpage

\section{Symmetric Matrices ($A = Q\Lambda Q^\T$) and Similar Matrices ($B = M^{-1}AM$)\bk{6.4, 6.6}}

\begin{goal}
Prove the Spectral Theorem $A = Q\Lambda Q^\T$ for real symmetric matrices, and show that similar matrices $B = M^{-1}AM$ share the same eigenvalues and Jordan form $J$.
\end{goal}

\step{Step 1: Prove that all eigenvalues of a real symmetric matrix $A = A^\T$ are real.}
Suppose $A\vect{x} = \lambda\vect{x}$. Take the complex conjugate $A\overline{\vect{x}} = \overline{\lambda}\overline{\vect{x}}$ and transpose to $\overline{\vect{x}}^\T A = \overline{\vect{x}}^\T\overline{\lambda}$ (since $A$ is real and $A^\T = A$). Take the dot product of the first equation with $\overline{\vect{x}}$ and the second with $\vect{x}$:
\begin{equation}
\overline{\vect{x}}^\T A\vect{x} = \lambda\overline{\vect{x}}^\T\vect{x} \quad \text{and} \quad \overline{\vect{x}}^\T A\vect{x} = \overline{\lambda}\overline{\vect{x}}^\T\vect{x} \implies (\lambda - \overline{\lambda})\|\vect{x}\|^2 = 0 \implies \lambda = \overline{\lambda} \quad (\text{real}).
\end{equation}

\step{Step 2: Prove that eigenvectors of $A = A^\T$ for distinct eigenvalues are orthogonal.}
Suppose $A\vect{x} = \lambda_1\vect{x}$ and $A\vect{y} = \lambda_2\vect{y}$ with $\lambda_1 \neq \lambda_2$:
\begin{equation}
\lambda_1\vect{x}^\T\vect{y} = (A\vect{x})^\T\vect{y} = \vect{x}^\T A^\T\vect{y} = \vect{x}^\T A\vect{y} = \lambda_2\vect{x}^\T\vect{y} \implies (\lambda_1 - \lambda_2)\vect{x}^\T\vect{y} = 0 \implies \vect{x}^\T\vect{y} = 0.
\end{equation}
By Schur's Theorem ($A = QTQ^{-1}$ with $T$ upper triangular), when $A = A^\T$ we have $T = Q^\T AQ = T^\T$, so $T = \Lambda$ is diagonal even when eigenvalues are repeated!

\begin{result}[ (Spectral Theorem, Sign of Pivots, and Similar Matrices)]
\begin{enumerate}
\item \textbf{Spectral Theorem:} Every real symmetric matrix $A = A^\T$ is diagonalized by an orthogonal eigenvector matrix $Q$ ($Q^{-1} = Q^\T$):
\begin{equation}
A = Q\Lambda Q^{-1} = Q\Lambda Q^\T = \lambda_1\vect{q}_1\vect{q}_1^\T + \lambda_2\vect{q}_2\vect{q}_2^\T + \dots + \lambda_n\vect{q}_n\vect{q}_n^\T.
\end{equation}
\item \textbf{Eigenvalues vs.\ Pivots:} For $A = A^\T = LDL^\T$, the number of positive eigenvalues equals the number of positive pivots.
\item \textbf{Similar Matrices:} $B = M^{-1}AM$ is similar to $A$. If $A\vect{x} = \lambda\vect{x}$, then $B(M^{-1}\vect{x}) = M^{-1}AM(M^{-1}\vect{x}) = \lambda(M^{-1}\vect{x})$, so $A$ and $B$ have the \term{same eigenvalues}, trace, determinant, rank, and number of independent eigenvectors.
\item \textbf{Jordan Form:} If $A$ has $s$ independent eigenvectors, it is similar to $J = \operatorname{diag}(J_1, \dots, J_s)$, where each Jordan block $J_i$ has $\lambda_i$ on the diagonal and $1$'s just above the diagonal.
\end{enumerate}
\end{result}

\begin{example}{6.4.1--6.4.4 and 6.6.1--6.6.4 (p.\ 331--357)}{Spectral Theorem $Q\Lambda Q^\T$, Pivots vs.\ Eigenvalues, and Jordan Form}
\begin{enumerate}
\item \textbf{Example 6.4.1:} Factor $A = \begin{bmatrix}1 & 2\\2 & 4\end{bmatrix}$ into $Q\Lambda Q^\T$.
\item \textbf{Example 6.4.4:} Compare the signs of the pivots and eigenvalues for $A = \begin{bmatrix}1 & 3\\3 & 1\end{bmatrix} = A^\T$ and $B = \begin{bmatrix}1 & 6\\-1 & -4\end{bmatrix} \neq B^\T$.
\item \textbf{Examples 6.6.3--6.6.4:} Solve $\frac{\diff\vect{u}}{\diff t} = J\vect{u}$ for the $3\times 3$ Jordan block $J = \begin{bmatrix}5&1&0\\0&5&1\\0&0&5\end{bmatrix}$.
\end{enumerate}
\end{example}

\vfill
\begin{pageend_sol}{6.4.1--6.4.4 and 6.6.1--6.6.4 (p.\ 331--357)}
\textbf{1.} $\lambda_1=0, \vect{q}_1 = \frac{1}{\sqrt{5}}\begin{bmatrix}2\\-1\end{bmatrix}$ and $\lambda_2=5, \vect{q}_2 = \frac{1}{\sqrt{5}}\begin{bmatrix}1\\2\end{bmatrix}$, so $Q = \frac{1}{\sqrt{5}}\begin{bmatrix}2&1\\-1&2\end{bmatrix}$ and $Q^\T AQ = \begin{bmatrix}0&0\\0&5\end{bmatrix} = \Lambda$.

\textbf{2.} Symmetric $A = \begin{bmatrix}1&3\\3&1\end{bmatrix}$ has pivots $1, -8$ and eigenvalues $4, -2$ (matching signs: one $+$, one $-$). Unsymmetric $B = \begin{bmatrix}1&6\\-1&-4\end{bmatrix}$ has pivots $1, 2$ (both $+$) but eigenvalues $-1, -2$ (both $-$).

\textbf{3.} Back substitution on $z' = 5z$, $y' = 5y + z$, $x' = 5x + y$ yields $z(t) = z(0)e^{5t}$, $y(t) = (y(0) + tz(0))e^{5t}$, and $x(t) = \left(x(0) + ty(0) + \frac{1}{2}t^2z(0)\right)e^{5t}$.
\end{pageend_sol}
\clearpage

\section{Singular Value Decomposition (SVD): $A = U\Sigma V^\T$\bk{6.7}}

\begin{goal}
Derive the Singular Value Decomposition $A = U\Sigma V^\T$ for any $m\times n$ matrix $A$ of rank $r$, and show how $U$ and $V$ provide orthonormal bases for all four fundamental subspaces.
\end{goal}

\step{Step 1: Find $\vect{v}_i$ and $\sigma_i^2$ from the symmetric positive semidefinite matrix $A^\T A$.}
Because $A^\T A$ is $n\times n$ symmetric positive semidefinite of rank $r$, it has $r$ positive eigenvalues $\sigma_1^2 \ge \sigma_2^2 \ge \dots \ge \sigma_r^2 > 0$ (and $n-r$ zero eigenvalues) with orthonormal eigenvectors $\vect{v}_1, \dots, \vect{v}_n$:
\begin{equation}
A^\T A\vect{v}_i = \sigma_i^2\vect{v}_i \quad (1 \le i \le r), \qquad A^\T A\vect{v}_i = \vect{0} \implies A\vect{v}_i = \vect{0} \quad (r+1 \le i \le n).
\end{equation}

\step{Step 2: Define $\vect{u}_i = A\vect{v}_i / \sigma_i$ for $1 \le i \le r$ and prove that the $\vect{u}_i$ are orthonormal eigenvectors of $AA^\T$.}
Take the inner product of $A\vect{v}_i$ with $A\vect{v}_j$:
\begin{equation}
(A\vect{v}_i)^\T(A\vect{v}_j) = \vect{v}_i^\T(A^\T A\vect{v}_j) = \vect{v}_i^\T(\sigma_j^2\vect{v}_j) = \sigma_j^2(\vect{v}_i^\T\vect{v}_j) = \begin{cases} \sigma_i^2 & \text{if } i = j \\ 0 & \text{if } i \neq j. \end{cases}
\end{equation}
Thus $\|A\vect{v}_i\| = \sigma_i$, and the unit vectors $\vect{u}_i = A\vect{v}_i / \sigma_i$ are orthonormal! Multiply $A^\T A\vect{v}_i = \sigma_i^2\vect{v}_i$ by $A$:
\begin{equation}
(AA^\T)(A\vect{v}_i) = \sigma_i^2(A\vect{v}_i) \implies AA^\T\vect{u}_i = \sigma_i^2\vect{u}_i.
\end{equation}
Complete $\vect{u}_1, \dots, \vect{u}_r$ with an orthonormal basis $\vect{u}_{r+1}, \dots, \vect{u}_m$ for $N(A^\T)$. Then $A\vect{v}_i = \sigma_i\vect{u}_i$ ($i \le r$) and $A\vect{v}_i = \vect{0}$ ($i > r$) combine into $AV = U\Sigma$:
\begin{result}[ (Singular Value Decomposition $A = U\Sigma V^\T$ and the Four Subspaces)]
\begin{equation}
A = U\Sigma V^\T = \sigma_1\vect{u}_1\vect{v}_1^\T + \sigma_2\vect{u}_2\vect{v}_2^\T + \dots + \sigma_r\vect{u}_r\vect{v}_r^\T,
\end{equation}
where $U$ ($m\times m$) and $V$ ($n\times n$) are orthogonal matrices ($U^\T U = I_m$, $V^\T V = I_n$):
\begin{itemize}
\item \textbf{First $r$ columns of $V$:} orthonormal basis for the \term{row space} $C(A^\T)$.
\item \textbf{Last $n-r$ columns of $V$:} orthonormal basis for the \term{nullspace} $N(A)$.
\item \textbf{First $r$ columns of $U$:} orthonormal basis for the \term{column space} $C(A)$.
\item \textbf{Last $m-r$ columns of $U$:} orthonormal basis for the \term{left nullspace} $N(A^\T)$.
\end{itemize}
\end{result}

\begin{example}{6.7.3--6.7.4 (p.\ 366--367)}{Computing the SVD of Invertible and Singular Matrices}
\begin{enumerate}
\item \textbf{Example 3:} Find the SVD $A = U\Sigma V^\T$ of $A = \begin{bmatrix}2 & 2\\-1 & 1\end{bmatrix}$.
\item \textbf{Example 4:} Find the full SVD $A = U\Sigma V^\T$ of the rank-one matrix $A = \begin{bmatrix}2 & 2\\1 & 1\end{bmatrix}$.
\end{enumerate}
\end{example}

\vfill
\begin{pageend_sol}{6.7.3--6.7.4 (p.\ 366--367)}
\textbf{1.} $A^\T A = \begin{bmatrix}5 & 3\\3 & 5\end{bmatrix}$ has eigenvalues $\sigma_1^2 = 8$ ($\sigma_1 = 2\sqrt{2}$) and $\sigma_2^2 = 2$ ($\sigma_2 = \sqrt{2}$), with unit eigenvectors $\vect{v}_1 = \begin{bmatrix}1/\sqrt{2}\\1/\sqrt{2}\end{bmatrix}$ and $\vect{v}_2 = \begin{bmatrix}-1/\sqrt{2}\\1/\sqrt{2}\end{bmatrix}$. Then $\vect{u}_1 = \frac{A\vect{v}_1}{2\sqrt{2}} = \begin{bmatrix}1\\0\end{bmatrix}$ and $\vect{u}_2 = \frac{A\vect{v}_2}{\sqrt{2}} = \begin{bmatrix}0\\1\end{bmatrix}$:
\[
A = \begin{bmatrix}2 & 2\\-1 & 1\end{bmatrix} = \begin{bmatrix}1 & 0\\0 & 1\end{bmatrix}\begin{bmatrix}2\sqrt{2} & 0\\0 & \sqrt{2}\end{bmatrix}\begin{bmatrix}1/\sqrt{2} & 1/\sqrt{2}\\-1/\sqrt{2} & 1/\sqrt{2}\end{bmatrix} = U\Sigma V^\T.
\]
\textbf{2.} Row space has $\vect{v}_1 = \frac{1}{\sqrt{2}}\begin{bmatrix}1\\1\end{bmatrix}$; nullspace has $\vect{v}_2 = \frac{1}{\sqrt{2}}\begin{bmatrix}1\\-1\end{bmatrix}$. Column space has $\vect{u}_1 = \frac{1}{\sqrt{5}}\begin{bmatrix}2\\1\end{bmatrix}$ with $A\vect{v}_1 = \sqrt{10}\vect{u}_1$ ($\sigma_1 = \sqrt{10}$); left nullspace has $\vect{u}_2 = \frac{1}{\sqrt{5}}\begin{bmatrix}1\\-2\end{bmatrix}$:
\[
\begin{bmatrix}2 & 2\\1 & 1\end{bmatrix} = \frac{1}{\sqrt{5}}\begin{bmatrix}2 & 1\\1 & -2\end{bmatrix}\begin{bmatrix}\sqrt{10} & 0\\0 & 0\end{bmatrix}\frac{1}{\sqrt{2}}\begin{bmatrix}1 & 1\\1 & -1\end{bmatrix} = U\Sigma V^\T.
\]
\end{pageend_sol}
\clearpage

% ============================================================================
% CHAPTER 8: MARKOV MATRICES AND COMPLEX MATRICES (8.3, 10.1--10.3)
% ============================================================================
\chapter[Markov Matrices and Complex Matrices (8.3, 10.1--10.3)]{Markov Matrices and Complex Vectors and Matrices\bk{8.3, 10.1--10.3}}

\begin{syll}
Syllabus topics: Markov matrices ($a_{ij} \ge 0$, columns add to $1$, $\lambda_{\max} = 1$, steady state $\vect{u}_\infty$) $\cdot$ Perron-Frobenius Theorem $\cdot$ Leslie population matrix and Leontief consumption matrix $(I - A)^{-1} = I + A + A^2 + \dots$ $\cdot$ Complex numbers $z = a + bi = re^{i\theta}$, Euler's formula, and $n$th roots of $1$ $\cdot$ Conjugate transpose $\vect{z}^\Herm = \overline{\vect{z}}^\T$, complex inner product $\vect{u}^\Herm\vect{v}$, and length $\|\vect{z}\|^2 = \vect{z}^\Herm\vect{z}$ $\cdot$ Hermitian matrices ($A^\Herm = A$), unitary matrices ($U^\Herm U = I$), and Fourier matrix $F_n$.
\end{syll}

\section{Markov Matrices, Population, and Economics\bk{8.3}}

\begin{goal}
Prove that every positive Markov matrix has $\lambda_1 = 1$ as its largest eigenvalue with a positive steady-state eigenvector $\vect{u}_\infty = \vect{x}_1$, and derive the geometric series $(I-A)^{-1} = I + A + A^2 + \dots$ for a productive economy.
\end{goal}

\step{Step 1: Properties of a Markov matrix $A$.}
A matrix $A$ is a \term{Markov matrix} if (1) every entry $a_{ij} \ge 0$, and (2) every column adds to $1$.
\begin{itemize}
\item \textbf{Why is $\lambda_1 = 1$ an eigenvalue?} Every column of $A - I$ adds to $1 - 1 = 0$, so the rows of $A - I$ add to the zero row. Thus $A - I$ is singular and $\lambda_1 = 1$ is an eigenvalue.
\item \textbf{Why can no eigenvalue have $|\lambda| > 1$?} Every power $A^k$ is also a Markov matrix with nonnegative entries between $0$ and $1$, so $A^k$ cannot blow up.
\item \textbf{Steady state (for $A > 0$ or $A^k > 0$):} All other eigenvalues satisfy $|\lambda_i| < 1$. For any starting probability vector $\vect{u}_0 = \vect{x}_1 + c_2\vect{x}_2 + \dots + c_n\vect{x}_n$:
\begin{equation}
\vect{u}_k = A^k\vect{u}_0 = \vect{x}_1 + c_2\lambda_2^k\vect{x}_2 + \dots + c_n\lambda_n^k\vect{x}_n \longrightarrow \vect{u}_\infty = \vect{x}_1 \quad \text{as } k \to \infty.
\end{equation}
\end{itemize}

\step{Step 2: Perron-Frobenius Theorem and the Consumption Matrix.}
\begin{result}[ (Perron-Frobenius Theorem and Leontief Input-Output Series)]
\begin{enumerate}
\item \textbf{Perron-Frobenius for $A > 0$:} The largest eigenvalue $\lambda_{\max}$ is real and strictly positive, and its eigenvector $\vect{x}$ has strictly positive entries: $A\vect{x} = \lambda_{\max}\vect{x} > \vect{0}$.
\item \textbf{Consumption Matrix ($p - Ap = y$):} To meet external demand $\vect{y} \ge \vect{0}$, production $\vect{p} = (I - A)^{-1}\vect{y}$ is nonnegative if and only if $\lambda_{\max}(A) < 1$, in which case the geometric series of matrices converges:
\begin{equation}
(I - A)^{-1} = I + A + A^2 + A^3 + \dots \ge 0 \qquad (\text{converges when all } |\lambda_i| < 1).
\end{equation}
\end{enumerate}
\end{result}

\begin{example}{8.3.1--8.3.5 (p.\ 432--437)}{Markov Chains and Consumption Matrices}
\begin{enumerate}
\item \textbf{Example 1:} Find $\vect{u}_k = A^k\vect{u}_0$ and the steady state $\vect{u}_\infty$ for $A = \begin{bmatrix}.80 & .05\\.20 & .95\end{bmatrix}$ starting from $\vect{u}_0 = \begin{bmatrix}.02\\.98\end{bmatrix}$.
\item \textbf{Example 3:} Find the eigenvalues and steady state of $A = \begin{bmatrix}0 & 1/2 & 1/2\\1/2 & 0 & 1/2\\1/2 & 1/2 & 0\end{bmatrix}$ starting from $\vect{u}_0 = (8,16,32)$.
\item \textbf{Examples 4--5:} Test whether the economy is productive ($\lambda_{\max} < 1$) for $A = \begin{bmatrix}.2&.3&.4\\.4&.4&.1\\.5&.1&.3\end{bmatrix}$ and $A = \begin{bmatrix}0&4\\1&0\end{bmatrix}$.
\end{enumerate}
\end{example}

\vfill
\begin{pageend_sol}{8.3.1--8.3.5 (p.\ 432--437)}
\textbf{1.} $\lambda_1 = 1$ has $\vect{x}_1 = \begin{bmatrix}.2\\.8\end{bmatrix}$ and $\lambda_2 = .75$ has $\vect{x}_2 = \begin{bmatrix}-1\\1\end{bmatrix}$. With $\vect{u}_0 = \begin{bmatrix}.02\\.98\end{bmatrix} = \begin{bmatrix}.2\\.8\end{bmatrix} + .18\begin{bmatrix}-1\\1\end{bmatrix}$, we get $\vect{u}_k = \begin{bmatrix}.2\\.8\end{bmatrix} + (.75)^k(.18)\begin{bmatrix}-1\\1\end{bmatrix} \to \vect{u}_\infty = \begin{bmatrix}.2\\.8\end{bmatrix}$.

\textbf{2.} $\lambda_1 = 1, \lambda_2 = \lambda_3 = -\frac{1}{2}$. Total population is $8+16+32 = 56$, so $\vect{u}_k \to 56\left(\frac{1}{3}, \frac{1}{3}, \frac{1}{3}\right) = \left(\frac{56}{3}, \frac{56}{3}, \frac{56}{3}\right)$.

\textbf{3.} In Example 4, every row adds to $.9$, so $\lambda_{\max} = .9 < 1$ and $(I-A)^{-1} \ge 0$ (productive). In Example 5, $\lambda_{\max} = 2 > 1$ and $(I-A)^{-1} = -\frac{1}{3}\begin{bmatrix}1&4\\1&1\end{bmatrix} < 0$ (not productive).
\end{pageend_sol}
\clearpage

\section{Complex Numbers, Hermitian Matrices, and Unitary Matrices\bk{10.1--10.3}}

\begin{goal}
Derive the polar form $z = re^{i\theta}$, the complex inner product $\vect{u}^\Herm\vect{v}$ and length $\|\vect{z}\|^2 = \vect{z}^\Herm\vect{z}$, the properties of Hermitian matrices ($A^\Herm = A$) and unitary matrices ($U^\Herm U = I$), and the unitary Fourier matrix $F_n$.
\end{goal}

\step{Step 1: Complex conjugate, polar form, Euler's formula, and $n$th roots of $1$.}
For $z = a + ib$ (with $i^2 = -1$), its conjugate is $\overline{z} = a - ib$ and its modulus is $r = |z| = \sqrt{a^2+b^2}$:
\begin{align}
z\overline{z} &= (a + ib)(a - ib) = a^2 + b^2 = |z|^2 = r^2, \qquad \frac{1}{a+ib} = \frac{a-ib}{a^2+b^2} \\
e^{i\theta} &= 1 + i\theta + \frac{(i\theta)^2}{2!} + \frac{(i\theta)^3}{3!} + \dots = \cos\theta + i\sin\theta \implies z = r(\cos\theta + i\sin\theta) = re^{i\theta} \\
z^n &= r^n(\cos n\theta + i\sin n\theta) = r^n e^{in\theta}, \qquad w = e^{2\pi i/n} \implies 1, w, w^2, \dots, w^{n-1} \text{ solve } z^n = 1.
\end{align}

\step{Step 2: Complex inner product $\vect{u}^\Herm\vect{v}$ and conjugate transpose $A^\Herm = \overline{A}^\T$.}
To make the squared length of a nonzero complex vector strictly positive, we must conjugate when we transpose ($\vect{z}^\Herm = \overline{\vect{z}}^\T$):
\begin{align}
\text{\textbf{Length squared:}}\quad \|\vect{z}\|^2 &= \vect{z}^\Herm\vect{z} = \overline{z}_1z_1 + \dots + \overline{z}_nz_n = |z_1|^2 + \dots + |z_n|^2 \\
\text{\textbf{Complex inner product:}}\quad \vect{u}^\Herm\vect{v} &= \overline{u}_1v_1 + \dots + \overline{u}_nv_n \\
(A\vect{u})^\Herm\vect{v} &= \vect{u}^\Herm(A^\Herm\vect{v}), \qquad (AB)^\Herm = B^\Herm A^\Herm.
\end{align}

\step{Step 3: Prove the three properties of Hermitian matrices ($A = A^\Herm$) and unitary matrices ($U^\Herm U = I$).}
\begin{enumerate}
\item \textbf{Real energy $\vect{z}^\Herm A\vect{z}$:} $(\vect{z}^\Herm A\vect{z})^\Herm = \vect{z}^\Herm A^\Herm(\vect{z}^\Herm)^\Herm = \vect{z}^\Herm A\vect{z}$, so $\vect{z}^\Herm A\vect{z}$ equals its own conjugate and is \term{real}.
\item \textbf{Real eigenvalues of $A = A^\Herm$:} If $A\vect{z} = \lambda\vect{z}$ ($\vect{z}\neq\vect{0}$), multiply by $\vect{z}^\Herm$: $\lambda = \frac{\vect{z}^\Herm A\vect{z}}{\vect{z}^\Herm\vect{z}} = \frac{\text{real}}{\text{real} > 0} = \text{real}$.
\item \textbf{Orthogonal eigenvectors of $A = A^\Herm$:} If $A\vect{z} = \lambda\vect{z}$ and $A\vect{y} = \beta\vect{y}$ ($\lambda \neq \beta$), then $\vect{y}^\Herm A\vect{z} = \lambda\vect{y}^\Herm\vect{z}$ and $(\vect{y}^\Herm A^\Herm)\vect{z} = \beta\vect{y}^\Herm\vect{z} \implies (\lambda - \beta)\vect{y}^\Herm\vect{z} = 0 \implies \vect{y}^\Herm\vect{z} = 0$.
\item \textbf{Unitary matrix $U$ ($U^\Herm U = I$, $U^{-1} = U^\Herm$):} Preserves length $\|U\vect{z}\|^2 = \vect{z}^\Herm U^\Herm U\vect{z} = \|\vect{z}\|^2$. Thus $U\vect{z} = \lambda\vect{z} \implies \|U\vect{z}\| = |\lambda|\|\vect{z}\| = \|\vect{z}\| \implies \eqbox{|\lambda| = 1.}$
\end{enumerate}

\begin{result}[ (Spectral Theorem for Hermitian Matrices and the Fourier Matrix $F_n$)]
Every Hermitian matrix $A = A^\Herm$ is diagonalized by a unitary matrix $U$: $A = U\Lambda U^{-1} = U\Lambda U^\Herm$.\\
The $n\times n$ \term{Fourier matrix} $F_n$ with $(j,k)$ entry $\frac{1}{\sqrt{n}}w^{jk}$ ($w = e^{2\pi i/n}$, $0 \le j,k \le n-1$) is unitary ($F_n^{-1} = F_n^\Herm$), and the Fast Fourier Transform (FFT) reduces the multiplication count for $n = 2^\ell$ from $n^2$ to $\frac{1}{2}n\log_2 n$.
\end{result}

\begin{example}{10.1.1 and 10.2.1--10.2.3 (p.\ 496--505)}{Polar Form, Hermitian Eigenvalues, and the $3\times 3$ Fourier Matrix $F_3$}
\begin{enumerate}
\item \textbf{Example 10.1.1:} Find $r$ and $\theta$ for $z = 1 + i$ and $\overline{z} = 1 - i$, and compute $(1+i)^8$.
\item \textbf{Example 10.2.1:} Find the complex inner product $\vect{u}^\Herm\vect{v}$ of $\vect{u} = \begin{bmatrix}1\\i\end{bmatrix}$ and $\vect{v} = \begin{bmatrix}i\\1\end{bmatrix}$.
\item \textbf{Example 10.2.2:} Find the eigenvalues and unit eigenvectors (unitary $U$) of $A = \begin{bmatrix}2 & 3-3i\\3+3i & 5\end{bmatrix} = A^\Herm$.
\item \textbf{Example 10.2.3:} Verify that the $3\times 3$ Fourier matrix $F_3 = \frac{1}{\sqrt{3}}\begin{bmatrix}1 & 1 & 1\\1 & e^{2\pi i/3} & e^{4\pi i/3}\\1 & e^{4\pi i/3} & e^{2\pi i/3}\end{bmatrix}$ is unitary and find $F_3^{-1}$.
\end{enumerate}
\end{example}

\vfill
\begin{pageend_sol}{10.1.1 and 10.2.1--10.2.3 (p.\ 496--505)}
\textbf{1.} $r = \sqrt{1^2+1^2} = \sqrt{2}$, $\theta = \pi/4$ for $1+i$ (and $-\pi/4$ or $7\pi/4$ for $1-i$). Thus $(1+i)^8 = (\sqrt{2})^8 e^{i8(\pi/4)} = 16e^{2\pi i} = 16$.

\textbf{2.} $\vect{u}^\Herm\vect{v} = \begin{bmatrix}1 & -i\end{bmatrix}\begin{bmatrix}i\\1\end{bmatrix} = i - i = 0$ (orthogonal!).

\textbf{3.} $\det(A-\lambda I) = (2-\lambda)(5-\lambda) - |3+3i|^2 = \lambda^2 - 7\lambda - 8 = (\lambda-8)(\lambda+1) = 0 \implies \lambda_1=8, \lambda_2=-1$. Eigenvectors are $\vect{z} = \begin{bmatrix}1\\1+i\end{bmatrix}$ and $\vect{y} = \begin{bmatrix}1-i\\-1\end{bmatrix}$ ($\vect{y}^\Herm\vect{z} = 0$), giving $U = \frac{1}{\sqrt{3}}\begin{bmatrix}1 & 1-i\\1+i & -1\end{bmatrix}$.

\textbf{4.} Each column of $F_3$ has length $\sqrt{\frac{1}{3}(1+1+1)} = 1$, and distinct columns have inner product $\frac{1}{3}(1 + e^{2\pi i/3} + e^{4\pi i/3}) = 0$. Thus $F_3$ is unitary, and $F_3^{-1} = F_3^\Herm = \overline{F}_3$ (replace $i$ by $-i$).
\end{pageend_sol}
\clearpage
% ============================================================================
% CHAPTER 12: MODULE 4 UNIQUE SOLVED PROBLEMS (WORKED EXAMPLES)
% ============================================================================
\chapter[Module 4: Unique Solved Problems]{Module 4: All Unique Solved Problems (Worked Examples)}

\begin{syll}
All unique Worked Examples from Module 4 (Sections 6.1, 6.2, 6.4, 6.6, and 6.7), with problem statements copied verbatim and full step-by-step solutions.
\end{syll}

\begin{example}{6.1 A}{Eigenvalues and Trace / Determinant of Matrix Powers}
Find the eigenvalues and eigenvectors of $A$ and $A^2$ and $A^{-1}$ and $A+4I$:
\[
A = \begin{bmatrix} 2 & -1 \\ -1 & 2 \end{bmatrix} \quad \text{and} \quad A^2 = \begin{bmatrix} 5 & -4 \\ -4 & 5 \end{bmatrix}.
\]
Check the trace $\lambda_1 + \lambda_2$ and the determinant $\lambda_1\lambda_2$ for $A$ and also $A^2$.
\end{example}
\textbf{Solution:} Compute the characteristic polynomial $\det(A - \lambda I) = 0$:
\[
\det(A - \lambda I) = \begin{vmatrix} 2-\lambda & -1 \\ -1 & 2-\lambda \end{vmatrix} = (2-\lambda)^2 - 1 = \lambda^2 - 4\lambda + 3 = (\lambda - 1)(\lambda - 3) = 0.
\]
The eigenvalues of $A$ are $\lambda_1 = 1$ and $\lambda_2 = 3$.
\begin{itemize}
\item \textbf{Trace check:} $\lambda_1 + \lambda_2 = 1 + 3 = 4 = 2 + 2 = \operatorname{trace}(A)$.
\item \textbf{Determinant check:} $\lambda_1\lambda_2 = (1)(3) = 3 = (2)(2) - (-1)(-1) = \det A$.
\end{itemize}
Find the eigenvectors by solving $(A - \lambda I)\vect{x} = \vect{0}$:
\begin{itemize}
\item For $\lambda_1 = 1$: $(A - I)\vect{x} = \begin{bmatrix} 1 & -1 \\ -1 & 1 \end{bmatrix}\begin{bmatrix} x_1 \\ x_2 \end{bmatrix} = \begin{bmatrix} 0 \\ 0 \end{bmatrix} \implies x_1 - x_2 = 0 \implies \vect{x}_1 = \begin{bmatrix} 1 \\ 1 \end{bmatrix}$.
\item For $\lambda_2 = 3$: $(A - 3I)\vect{x} = \begin{bmatrix} -1 & -1 \\ -1 & -1 \end{bmatrix}\begin{bmatrix} x_1 \\ x_2 \end{bmatrix} = \begin{bmatrix} 0 \\ 0 \end{bmatrix} \implies x_1 + x_2 = 0 \implies \vect{x}_2 = \begin{bmatrix} -1 \\ 1 \end{bmatrix}$.
\end{itemize}
The matrices $A^2, A^{-1}$, and $A + 4I$ keep the same eigenvectors $\vect{x}_1$ and $\vect{x}_2$:
\begin{itemize}
\item $A^2$ has eigenvalues $\lambda_1^2 = 1^2 = 1$ and $\lambda_2^2 = 3^2 = 9$. Trace: $1+9 = 10 = 5+5$. Determinant: $(1)(9) = 9 = 25 - 16$.
\item $A^{-1}$ has eigenvalues $1/\lambda_1 = 1/1 = 1$ and $1/\lambda_2 = 1/3$.
\item $A + 4I$ has eigenvalues $\lambda_1 + 4 = 1 + 4 = 5$ and $\lambda_2 + 4 = 3 + 4 = 7$.
\end{itemize}

\begin{example}{6.1 B}{Eigenvalues and Orthogonal Eigenvectors of a 3 by 3 Singular Matrix}
Find the eigenvalues and eigenvectors of this $3$ by $3$ symmetric, singular matrix:
\[
A = \begin{bmatrix} 1 & -1 & 0 \\ -1 & 2 & -1 \\ 0 & -1 & 1 \end{bmatrix}.
\]
\end{example}
\textbf{Solution:} Notice that all rows of $A$ add to zero: $1 + (-1) + 0 = 0$. Thus the vector $\vect{x}_1 = (1, 1, 1)$ satisfies $A\vect{x}_1 = \vect{0}$. This confirms that $\lambda_1 = 0$ is an eigenvalue.
Compute the determinant $\det(A - \lambda I)$:
\begin{align*}
\det(A - \lambda I) &= \begin{vmatrix} 1-\lambda & -1 & 0 \\ -1 & 2-\lambda & -1 \\ 0 & -1 & 1-\lambda \end{vmatrix} = (1-\lambda)[(2-\lambda)(1-\lambda) - 1] - (-1)[-(1-\lambda)] \\
&= (1-\lambda)[\lambda^2 - 3\lambda + 1 - 1] = (1-\lambda)(\lambda^2 - 3\lambda) = -\lambda(1-\lambda)(3-\lambda) = 0.
\end{align*}
The eigenvalues are $\lambda_1 = 0, \lambda_2 = 1$, and $\lambda_3 = 3$.
\begin{itemize}
\item \textbf{Trace check:} $\lambda_1 + \lambda_2 + \lambda_3 = 0 + 1 + 3 = 4 = 1 + 2 + 1 = \operatorname{trace}(A)$.
\item \textbf{Determinant check:} $\lambda_1\lambda_2\lambda_3 = (0)(1)(3) = 0 = \det A$.
\end{itemize}
Find the eigenvectors:
\begin{itemize}
\item For $\lambda_1 = 0$: $\vect{x}_1 = (1, 1, 1)$.
\item For $\lambda_2 = 1$: $(A - I)\vect{x} = \begin{bmatrix} 0 & -1 & 0 \\ -1 & 1 & -1 \\ 0 & -1 & 0 \end{bmatrix}\begin{bmatrix} x_1 \\ x_2 \\ x_3 \end{bmatrix} = \begin{bmatrix} 0 \\ 0 \\ 0 \end{bmatrix} \implies x_2 = 0, x_1 + x_3 = 0 \implies \vect{x}_2 = \begin{bmatrix} 1 \\ 0 \\ -1 \end{bmatrix}$.
\item For $\lambda_3 = 3$: $(A - 3I)\vect{x} = \begin{bmatrix} -2 & -1 & 0 \\ -1 & -1 & -1 \\ 0 & -1 & -2 \end{bmatrix}\begin{bmatrix} x_1 \\ x_2 \\ x_3 \end{bmatrix} = \begin{bmatrix} 0 \\ 0 \\ 0 \end{bmatrix} \implies \vect{x}_3 = \begin{bmatrix} 1 \\ -2 \\ 1 \end{bmatrix}$.
\end{itemize}
Because $A$ is symmetric ($A = A^\T$), the eigenvectors are mutually orthogonal:
\[
\vect{x}_1\cdot\vect{x}_2 = 1(1) + 1(0) + 1(-1) = 0, \quad \vect{x}_1\cdot\vect{x}_3 = 1 - 2 + 1 = 0, \quad \vect{x}_2\cdot\vect{x}_3 = 1 + 0 - 1 = 0.
\]

\begin{example}{6.2 A}{Powers and Lucas Numbers $L_{100} = \lambda_1^{100} + \lambda_2^{100}$}
The Lucas numbers follow the Fibonacci rule $L_{k+2} = L_{k+1} + L_k$ starting with $L_1 = 1$ and $L_2 = 3$. Following this rule, the next numbers are $4, 7, 11, 18$. Show that the Lucas number $L_{100}$ is $\lambda_1^{100} + \lambda_2^{100}$ where $\lambda_1, \lambda_2 = (1\pm\sqrt{5})/2$.
\end{example}
\textbf{Solution:} Define the state vector $\vect{u}_k = \begin{bmatrix} L_{k+1} \\ L_k \end{bmatrix}$. The recurrence relation becomes the matrix equation:
\[
\vect{u}_{k+1} = \begin{bmatrix} L_{k+2} \\ L_{k+1} \end{bmatrix} = \begin{bmatrix} 1 & 1 \\ 1 & 0 \end{bmatrix}\begin{bmatrix} L_{k+1} \\ L_k \end{bmatrix} = A\vect{u}_k.
\]
The eigenvalues of $A$ come from $\det(A-\lambda I) = \lambda^2 - \lambda - 1 = 0$:
\[
\lambda_1 = \frac{1+\sqrt{5}}{2}, \qquad \lambda_2 = \frac{1-\sqrt{5}}{2}, \qquad \text{with eigenvectors } \vect{x}_1 = \begin{bmatrix} \lambda_1 \\ 1 \end{bmatrix}, \quad \vect{x}_2 = \begin{bmatrix} \lambda_2 \\ 1 \end{bmatrix}.
\]
Expand the starting vector $\vect{u}_1 = \begin{bmatrix} L_2 \\ L_1 \end{bmatrix} = \begin{bmatrix} 3 \\ 1 \end{bmatrix}$ in the basis of eigenvectors:
\[
c_1\vect{x}_1 + c_2\vect{x}_2 = c_1\begin{bmatrix} \lambda_1 \\ 1 \end{bmatrix} + c_2\begin{bmatrix} \lambda_2 \\ 1 \end{bmatrix} = \begin{bmatrix} 3 \\ 1 \end{bmatrix}.
\]
Notice that $\lambda_1 + \lambda_2 = 1$ and $\lambda_1^2 + \lambda_2^2 = \operatorname{trace}(A^2) = 3$. Setting $c_1 = \lambda_1$ and $c_2 = \lambda_2$ satisfies both equations:
\[
\lambda_1\vect{x}_1 + \lambda_2\vect{x}_2 = \begin{bmatrix} \lambda_1^2 + \lambda_2^2 \\ \lambda_1 + \lambda_2 \end{bmatrix} = \begin{bmatrix} 3 \\ 1 \end{bmatrix} = \vect{u}_1.
\]
Multiply by $A^{99}$ to obtain $\vect{u}_{100}$:
\[
\vect{u}_{100} = A^{99}\vect{u}_1 = \lambda_1(A^{99}\vect{x}_1) + \lambda_2(A^{99}\vect{x}_2) = \lambda_1(\lambda_1^{99}\vect{x}_1) + \lambda_2(\lambda_2^{99}\vect{x}_2) = \lambda_1^{100}\begin{bmatrix} \lambda_1 \\ 1 \end{bmatrix} + \lambda_2^{100}\begin{bmatrix} \lambda_2 \\ 1 \end{bmatrix}.
\]
The second component gives the required formula: $L_{100} = \lambda_1^{100} + \lambda_2^{100}$.

\begin{example}{6.2 B}{Inverse, Eigenvalues, and Determinant of All-Ones Shift $5I - \operatorname{ones}(4)$}
Find the inverse, the eigenvalues, and the determinant of:
\[
A = 5I - \operatorname{ones}(4) = \begin{bmatrix} 4 & -1 & -1 & -1 \\ -1 & 4 & -1 & -1 \\ -1 & -1 & 4 & -1 \\ -1 & -1 & -1 & 4 \end{bmatrix}.
\]
Describe an eigenvector matrix $S$ that satisfies $S^{-1}AS = \Lambda$.
\end{example}
\textbf{Solution:} The all-ones matrix $B = \operatorname{ones}(4)$ has rank $1$. Its trace is $1+1+1+1 = 4$, so its eigenvalues are $4, 0, 0, 0$.
Subtracting $B$ from $5I$ shifts the eigenvalues:
\[
\lambda_1 = 5 - 4 = 1, \qquad \lambda_2 = \lambda_3 = \lambda_4 = 5 - 0 = 5.
\]
The determinant of $A$ is the product of its eigenvalues: $\det A = (1)(5)(5)(5) = 125$.
The eigenvector for $\lambda_1 = 1$ is $\vect{x}_1 = (1, 1, 1, 1)$. The other eigenvectors are orthogonal to $\vect{x}_1$. A complete set of orthonormal eigenvectors forms the normalized Hadamard matrix $S = \frac{1}{2}H_4$:
\[
S = \frac{1}{2}\begin{bmatrix} 1 & 1 & 1 & 1 \\ 1 & -1 & 1 & -1 \\ 1 & 1 & -1 & -1 \\ 1 & -1 & -1 & 1 \end{bmatrix} \implies S^\T S = I \quad \text{and} \quad S^{-1}AS = \operatorname{diag}(1, 5, 5, 5).
\]
The eigenvalues of $A^{-1}$ are $1, 1/5, 1/5, 1/5$. The inverse is:
\[
A^{-1} = \frac{1}{5}\left(I + \operatorname{ones}(4)\right) = \frac{1}{5}\begin{bmatrix} 2 & 1 & 1 & 1 \\ 1 & 2 & 1 & 1 \\ 1 & 1 & 2 & 1 \\ 1 & 1 & 1 & 2 \end{bmatrix}.
\]

\begin{example}{6.4 A}{Reflection Matrix from Orthonormal Eigenvectors}
What matrix $A$ has eigenvalues $\lambda = 1, -1$ and eigenvectors $\vect{x}_1 = (\cos\theta, \sin\theta)$ and $\vect{x}_2 = (-\sin\theta, \cos\theta)$? Which properties of $A$ can be predicted before computation?
\end{example}
\textbf{Solution:} Real eigenvalues in $\Lambda = \operatorname{diag}(1,-1)$ and orthonormal eigenvectors in $Q = \begin{bmatrix} \cos\theta & -\sin\theta \\ \sin\theta & \cos\theta \end{bmatrix}$ guarantee that $A = Q\Lambda Q^\T$ is symmetric ($A = A^\T$).
Properties predicted before computation:
\begin{enumerate}
\item $A^2 = I$ because $\lambda_i^2 = (\pm 1)^2 = 1$. Thus $A^{-1} = A$.
\item $\det A = (1)(-1) = -1$.
\item The pivots of $A$ have opposite signs ($d_1 > 0$ and $d_2 < 0$), matching the signs of the eigenvalues.
\item $A$ is a reflection matrix across the line at angle $\theta$.
\end{enumerate}
Multiply $A = Q\Lambda Q^\T$ using $c = \cos\theta$ and $s = \sin\theta$:
\[
A = \begin{bmatrix} c & -s \\ s & c \end{bmatrix}\begin{bmatrix} 1 & 0 \\ 0 & -1 \end{bmatrix}\begin{bmatrix} c & s \\ -s & c \end{bmatrix} = \begin{bmatrix} c^2 - s^2 & 2cs \\ 2cs & s^2 - c^2 \end{bmatrix} = \begin{bmatrix} \cos 2\theta & \sin 2\theta \\ \sin 2\theta & -\cos 2\theta \end{bmatrix}.
\]
Multiplying $(1,0)$ by $A$ gives $(\cos 2\theta, \sin 2\theta)$, which is the reflection of $(1,0)$ across the $\theta$-line.

\begin{example}{6.4 B}{Discrete Sines and Cosines from Tridiagonal Matrices}
Find the eigenvalues of $A_3$ and $B_4$, and check the orthogonality of their eigenvectors:
\[
A_3 = \begin{bmatrix} 2 & -1 & 0 \\ -1 & 2 & -1 \\ 0 & -1 & 2 \end{bmatrix} \qquad \text{and} \qquad B_4 = \begin{bmatrix} 1 & -1 & 0 & 0 \\ -1 & 2 & -1 & 0 \\ 0 & -1 & 2 & -1 \\ 0 & 0 & -1 & 1 \end{bmatrix}.
\]
\end{example}
\textbf{Solution:}
\begin{itemize}
\item \textbf{For $A_3$ (fixed boundaries):} The eigenvalues are $\lambda_k = 2 - 2\cos\frac{k\pi}{4}$ for $k = 1, 2, 3$:
\[
\lambda_1 = 2 - \sqrt{2}, \qquad \lambda_2 = 2, \qquad \lambda_3 = 2 + \sqrt{2}.
\]
Trace: $(2-\sqrt{2}) + 2 + (2+\sqrt{2}) = 6 = 2+2+2$. Determinant: $(2-\sqrt{2})(2)(2+\sqrt{2}) = 4$.
The eigenvectors are discrete sines $\vect{s}_k = (\sin\frac{k\pi}{4}, \sin\frac{2k\pi}{4}, \sin\frac{3k\pi}{4})$, forming the Discrete Sine Transform (DST).
\item \textbf{For $B_4$ (free boundaries):} All rows add to zero, so $\lambda_1 = 0$ with constant eigenvector $\vect{x}_1 = (1, 1, 1, 1)$.
The other three eigenvalues match $A_3$: $\lambda_2 = 2 - \sqrt{2}, \lambda_3 = 2, \lambda_4 = 2 + \sqrt{2}$.
The determinant is $\det B_4 = 0$ and the trace is $6$.
The eigenvectors are discrete cosines evaluated at midpoints $\frac{\pi}{8}, \frac{3\pi}{8}, \frac{5\pi}{8}, \frac{7\pi}{8}$, forming the Discrete Cosine Transform (DCT) used in JPEG compression.
\end{itemize}

\begin{example}{6.6 A}{Similar Matrices: Triangular Pascal Matrix and Its Inverse}
The $4$ by $4$ triangular Pascal matrix $A$ and its inverse $A^{-1}$ are:
\[
A = \begin{bmatrix} 1 & 0 & 0 & 0 \\ 1 & 1 & 0 & 0 \\ 1 & 2 & 1 & 0 \\ 1 & 3 & 3 & 1 \end{bmatrix} \qquad \text{and} \qquad A^{-1} = \begin{bmatrix} 1 & 0 & 0 & 0 \\ -1 & 1 & 0 & 0 \\ 1 & -2 & 1 & 0 \\ -1 & 3 & -3 & 1 \end{bmatrix}.
\]
Check that $A$ and $A^{-1}$ have the same eigenvalues. Find a diagonal matrix $D$ with alternating signs that gives $A^{-1} = D^{-1}AD$. What is their Jordan form $J$?
\end{example}
\textbf{Solution:} Both triangular matrices have diagonal entries equal to $1$, so both have eigenvalues $\lambda = 1, 1, 1, 1$.
Choose $D = \operatorname{diag}(1, -1, 1, -1)$. Because $D^2 = I$, we have $D^{-1} = D$.
Multiplying $D^{-1}AD$ multiplies row $i$ by $(-1)^{i-1}$ and column $j$ by $(-1)^{j-1}$. Thus entry $(i,j)$ is multiplied by $(-1)^{i-j}$. On subdiagonals where $i-j$ is odd, the sign reverses:
\[
D^{-1}AD = \begin{bmatrix} 1 & 0 & 0 & 0 \\ 0 & -1 & 0 & 0 \\ 0 & 0 & 1 & 0 \\ 0 & 0 & 0 & -1 \end{bmatrix}\begin{bmatrix} 1 & 0 & 0 & 0 \\ 1 & 1 & 0 & 0 \\ 1 & 2 & 1 & 0 \\ 1 & 3 & 3 & 1 \end{bmatrix}\begin{bmatrix} 1 & 0 & 0 & 0 \\ 0 & -1 & 0 & 0 \\ 0 & 0 & 1 & 0 \\ 0 & 0 & 0 & -1 \end{bmatrix} = \begin{bmatrix} 1 & 0 & 0 & 0 \\ -1 & 1 & 0 & 0 \\ 1 & -2 & 1 & 0 \\ -1 & 3 & -3 & 1 \end{bmatrix} = A^{-1}.
\]
This proves $A$ is similar to $A^{-1}$. Because $A - I$ has rank $3$, there is only $4-3 = 1$ independent eigenvector $(0,0,0,1)$. The Jordan form has a single $4\times 4$ block:
\[
J = \begin{bmatrix} 1 & 1 & 0 & 0 \\ 0 & 1 & 1 & 0 \\ 0 & 0 & 1 & 1 \\ 0 & 0 & 0 & 1 \end{bmatrix}.
\]

\begin{example}{6.6 B}{Eigenvalue Computation: QR Algorithm and Cubic Convergence}
Explain how the QR algorithm computes eigenvalues without solving $\det(A-\lambda I) = 0$. For $A = \begin{bmatrix} c & s \\ s & 0 \end{bmatrix}$ (with $c=\cos\theta, s=\sin\theta$), show how one step $A_1 = RQ$ cubes the small off-diagonal entry $s$.
\end{example}
\textbf{Solution:} Factor $A = QR$ by Gram-Schmidt, where $Q$ is orthogonal and $R$ is upper triangular. Then reverse the factors:
\[
A_1 = RQ = Q^{-1}(QR)Q = Q^{-1}AQ.
\]
Because $A_1$ is similar to $A$, both matrices have the same eigenvalues.
For $A = \begin{bmatrix} c & s \\ s & 0 \end{bmatrix}$:
\[
A = QR = \begin{bmatrix} c & -s \\ s & c \end{bmatrix}\begin{bmatrix} 1 & cs \\ 0 & -s^2 \end{bmatrix} \implies A_1 = RQ = \begin{bmatrix} 1 & cs \\ 0 & -s^2 \end{bmatrix}\begin{bmatrix} c & -s \\ s & c \end{bmatrix} = \begin{bmatrix} c + cs^2 & -s + c^2s \\ -s^3 & -cs^2 \end{bmatrix}.
\]
The entry below the diagonal has dropped from $s$ to $-s^3$. Each iteration of the shifted QR method cubes the error, producing very rapid convergence to triangular form.

\begin{example}{6.7 A}{Three Fundamental Matrix Factorizations: SVD, QR, and LU}
Identify by name these decompositions $A = \vect{c}_1\vect{b}_1^\T + \dots + \vect{c}_r\vect{b}_r^\T$ of an $m$ by $n$ matrix of rank $r$:
\begin{enumerate}
\item Orthogonal columns $\vect{c}_1, \dots, \vect{c}_r$ and orthogonal rows $\vect{b}_1^\T, \dots, \vect{b}_r^\T$.
\item Orthogonal columns $\vect{c}_1, \dots, \vect{c}_r$ and triangular rows $\vect{b}_1^\T, \dots, \vect{b}_r^\T$.
\item Triangular columns $\vect{c}_1, \dots, \vect{c}_r$ and triangular rows $\vect{b}_1^\T, \dots, \vect{b}_r^\T$.
\end{enumerate}
Where do singular values, pivots, and heights enter?
\end{example}
\textbf{Solution:}
\begin{enumerate}
\item \textbf{Singular Value Decomposition (SVD):} $A = U\Sigma V^\T$. Orthonormal columns in $U$, orthonormal rows in $V^\T$, and singular values $\sigma_i$ on the diagonal of $\Sigma$.
\item \textbf{Gram-Schmidt Factorization:} $A = QR = QHR_0$. Orthonormal columns in $Q$, upper triangular rows in $R_0$, and column heights $h_i$ on the diagonal of $H$.
\item \textbf{Gaussian Elimination:} $A = LU = LDU_0$. Lower triangular columns in $L$, upper triangular rows in $U_0$, and pivots $d_i$ on the diagonal of $D$.
\end{enumerate}
For a square invertible matrix ($r = m = n$), the volume of the box is:
\[
|\det A| = \sigma_1\sigma_2\cdots\sigma_n = |d_1 d_2 \cdots d_n| = h_1 h_2 \cdots h_n.
\]

\begin{example}{6.7 B}{Rank-One SVD vs.\ Diagonalization and Jordan Form}
For a rank-one $2$ by $2$ matrix $A = \vect{x}\vect{y}^\T$, compare the Singular Value Decomposition $A = U\Sigma V^\T$ with diagonalization $A = S\Lambda S^{-1}$ and Jordan form.
\end{example}
\textbf{Solution:} Multiply $A$ by $\vect{x}$: $A\vect{x} = (\vect{x}\vect{y}^\T)\vect{x} = (\vect{y}^\T\vect{x})\vect{x}$.
\begin{itemize}
\item If $\vect{y}^\T\vect{x} \neq 0$: The eigenvalues are $\lambda_1 = \vect{y}^\T\vect{x} = \operatorname{trace}(A)$ and $\lambda_2 = 0$. The eigenvectors are $\vect{x}_1 = \vect{x}$ and any vector $\vect{x}_2$ perpendicular to $\vect{y}$ ($\vect{y}^\T\vect{x}_2 = 0$). The matrix is diagonalizable: $A = S\Lambda S^{-1}$.
\item If $\vect{y}^\T\vect{x} = 0$: Both eigenvalues are zero ($\lambda_1 = \lambda_2 = 0$). Since $A \neq \vect{0}$, there is only one independent eigenvector. The matrix cannot be diagonalized; its Jordan form is $J = \begin{bmatrix} 0 & 1 \\ 0 & 0 \end{bmatrix}$.
\item In \textbf{both} cases, SVD always exists without restriction:
\[
A^\T A = (\vect{y}\vect{x}^\T)(\vect{x}\vect{y}^\T) = \|\vect{x}\|^2\vect{y}\vect{y}^\T \implies \sigma_1^2 = \|\vect{x}\|^2\|\vect{y}\|^2 \implies \sigma_1 = \|\vect{x}\|\|\vect{y}\|.
\]
The unit singular vectors are $\vect{v}_1 = \vect{y}/\|\vect{y}\|$ and $\vect{u}_1 = A\vect{v}_1/\sigma_1 = \vect{x}/\|\vect{x}\|$, with $\sigma_2 = 0$:
\[
A = \sigma_1\vect{u}_1\vect{v}_1^\T = (\|\vect{x}\|\|\vect{y}\|)\left(\frac{\vect{x}}{\|\vect{x}\|}\right)\left(\frac{\vect{y}}{\|\vect{y}\|}\right)^\T = \vect{x}\vect{y}^\T.
\]
\end{itemize}

\begin{example}{6.7 C}{Singular Values Dominating Eigenvalues and Matrix Entries}
Show that $\sigma_1 \ge |\lambda|_{\max}$ (the largest singular value dominates all eigenvalues). Show that $\sigma_1 \ge |a_{ij}|_{\max}$ (the largest singular value dominates all entries of $A$). Estimate $\sigma_1$ and $\sigma_2$ for $A = \begin{bmatrix} 1 & C \\ 0 & 1 \end{bmatrix}$.
\end{example}
\textbf{Solution:} Start from $A = U\Sigma V^\T$. Orthogonal matrices preserve length: $\|U\vect{z}\| = \|\vect{z}\|$ and $\|V^\T\vect{z}\| = \|\vect{z}\|$. For any vector $\vect{x}$:
\[
\|A\vect{x}\| = \|U\Sigma V^\T\vect{x}\| = \|\Sigma(V^\T\vect{x})\| \le \sigma_1\|V^\T\vect{x}\| = \sigma_1\|\vect{x}\|.
\]
\begin{enumerate}
\item If $\vect{x}$ is a unit eigenvector ($A\vect{x} = \lambda\vect{x}, \|\vect{x}\|=1$), then $|\lambda| = \|A\vect{x}\| \le \sigma_1$. Thus $\sigma_1 \ge |\lambda|_{\max}$.
\item If $\vect{x} = \vect{e}_j$ (unit coordinate vector), $A\vect{e}_j$ is column $j$ of $A$. Its length is $\sqrt{\sum_i a_{ij}^2} \le \sigma_1$. Thus every entry satisfies $|a_{ij}| \le \|A\vect{e}_j\| \le \sigma_1$.
\item For $A = \begin{bmatrix} 1 & C \\ 0 & 1 \end{bmatrix}$:
Column 2 has length $\sqrt{1+C^2}$, so $\sigma_1 \ge \sqrt{1+C^2} > |C|$.
Since $\det A = 1$, we have $\det(A^\T A) = \sigma_1^2\sigma_2^2 = 1 \implies \sigma_1\sigma_2 = 1 \implies \sigma_2 = 1/\sigma_1 < 1/|C|$.
For $C = 1000$, $\sigma_1 \approx 1000$ and $\sigma_2 \approx 1/1000$. The matrix is ill-conditioned with condition number $\sigma_1/\sigma_2 \approx 10^6$, even though both eigenvalues are $\lambda = 1$.
\end{enumerate}

% ============================================================================
% FORMULA SHEETS FOR MODULES 1 TO 4
% ============================================================================
\formulasheet{Module 1 Formula Sheet: Vectors, Vector Spaces, and Subspaces}

\begin{fs}{1. Vector Operations, Linear Combinations, and Subspaces}
\begin{align*}
  \vect{v}+\vect{w} &= (v_1+w_1, \dots, v_n+w_n), \qquad c\vect{v} = (cv_1, \dots, cv_n), \qquad \text{Combinations: } c\vect{v}+d\vect{w}\\
  \text{Subspace Test:}\quad & c\vect{v}+d\vect{w} \text{ stays in } \mathbf{S} \quad (\text{must contain } \vect{0})\\
  \text{Column Space } C(A):\quad & \text{all combinations of columns of } A, \quad A\vect{x}=\vect{b} \text{ solvable } \iff \vect{b} \in C(A)\\
  \text{Nullspace } N(A):\quad & \text{all solutions to } A\vect{x}=\vect{0}, \quad \dim N(A) = n - r \text{ (free variables)}
\end{align*}
\end{fs}

\begin{fs}{2. Reduced Echelon Form $R$, Complete Solution, and Four Subspaces}
\begin{align*}
  R &= \operatorname{rref}(A) = \begin{bmatrix} I & F \\ \vect{0} & \vect{0} \end{bmatrix}, \qquad \text{Nullspace Matrix: } N = \begin{bmatrix} -F \\ I \end{bmatrix} \implies RN = \vect{0}, \ AN = \vect{0}\\
  \text{Complete Solution:}\quad & \vect{x}_{\text{complete}} = \vect{x}_p + \vect{x}_n = \vect{x}_p + c_1\vect{s}_1 + \dots + c_{n-r}\vect{s}_{n-r}\\
  \text{FTLA Part 1:}\quad & \dim C(A) = \dim C(A^\T) = r \quad (\text{rank}), \qquad \dim N(A) = n - r, \qquad \dim N(A^\T) = m - r
\end{align*}
\end{fs}

\formulasheet{Module 2 Formula Sheet: Linear Equations, Inverses, LU, and Determinants}

\begin{fs}{1. Elimination, Matrix Multiplication, Inverses, and $A = LU$}
\begin{align*}
  \text{Multipliers:}\quad & \ell_{ij} = \frac{a_{ij}}{d_j}, \qquad (AB)_{ij} = \sum_{k=1}^n a_{ik}b_{kj}, \qquad A = \sum_{k=1}^n (\text{col } k)(\text{row } k), \qquad (AB)^{-1} = B^{-1}A^{-1}\\
  \text{Gauss-Jordan:}\quad & [A \ \ I] \longrightarrow [I \ \ A^{-1}], \qquad \begin{bmatrix} a & b \\ c & d \end{bmatrix}^{-1} = \frac{1}{ad-bc}\begin{bmatrix} d & -b \\ -c & a \end{bmatrix}\\
  \text{Factorization:}\quad & A = LU, \quad A = LDU, \quad A = LDL^\T \ (A=A^\T), \quad PA = LU \ (P^{-1}=P^\T)\\
  \text{Operation Counts:}\quad & \text{Factor: } \frac{1}{3}n^3 \text{ mults/subs}, \qquad \text{Solve } (L\vect{c}=\vect{b}, U\vect{x}=\vect{c}): n^2 \text{ mults/subs}
\end{align*}
\end{fs}

\begin{fs}{2. Determinant Formulas, Cramer's Rule, and Volumes}
\begin{align*}
  \text{Pivot Formula:}\quad & \det A = (\det P)(d_1 d_2 \cdots d_n), \qquad d_k = \frac{\det A_k}{\det A_{k-1}}, \qquad \det(AB) = (\det A)(\det B), \quad \det A^\T = \det A\\
  \text{Big Formula:}\quad & \det A = \sum (\det P)\, a_{1\alpha}a_{2\beta}\cdots a_{n\omega}, \qquad \text{Cofactors: } C_{ij} = (-1)^{i+j}\det M_{ij}\\
  \text{Cramer's Rule:}\quad & x_j = \frac{\det B_j}{\det A}, \qquad A^{-1} = \frac{C^\T}{\det A}, \qquad \text{Box Volume: } |\det A| = |(\vect{u}\times\vect{v})\cdot\vect{w}|
\end{align*}
\end{fs}

\formulasheet{Module 3 Formula Sheet: Orthogonality, Projections, and Statistics}

\begin{fs}{1. Dot Products, Orthogonal Subspaces, Projections, and Gram-Schmidt}
\begin{align*}
  \text{Dot Product / Length:}\quad & \vect{v}\cdot\vect{w} = \vect{v}^\T\vect{w}, \quad \|\vect{v}\| = \sqrt{\vect{v}^\T\vect{v}}, \quad \cos\theta = \frac{\vect{v}^\T\vect{w}}{\|\vect{v}\|\|\vect{w}\|}, \quad |\vect{v}^\T\vect{w}| \le \|\vect{v}\|\|\vect{w}\|\\
  \text{FTLA Part 2:}\quad & N(A) = C(A^\T)^\perp \ (\text{in } \mathbf{R}^n), \qquad N(A^\T) = C(A)^\perp \ (\text{in } \mathbf{R}^m)\\
  \text{Projection Onto Line:}\quad & \uv{x} = \frac{\vect{a}^\T\vect{b}}{\vect{a}^\T\vect{a}}, \quad \vect{p} = \vect{a}\uv{x} = P\vect{b}, \quad P = \frac{\vect{a}\vect{a}^\T}{\vect{a}^\T\vect{a}}, \quad P^2 = P = P^\T\\
  \text{Projection Onto } C(A):\quad & A^\T A\uv{\vect{x}} = A^\T\vect{b}, \quad \uv{\vect{x}} = (A^\T A)^{-1}A^\T\vect{b}, \quad \vect{p} = P\vect{b} = A(A^\T A)^{-1}A^\T\vect{b}\\
  \text{Gram-Schmidt:}\quad & \vect{q}_1 = \frac{\vect{a}}{\|\vect{a}\|}, \quad \vect{q}_2 = \frac{\vect{b} - (\vect{q}_1^\T\vect{b})\vect{q}_1}{\|\vect{b} - (\vect{q}_1^\T\vect{b})\vect{q}_1\|}, \quad A = QR \implies R\uv{\vect{x}} = Q^\T\vect{b}
\end{align*}
\end{fs}

\begin{fs}{2. Positive Definiteness, Weighted Least Squares, and Statistics}
\begin{align*}
  \text{Positive Definite:}\quad & \text{all } \lambda_i > 0 \iff \text{all pivots } d_i > 0 \iff \text{all upper-left } \det A_k > 0 \iff \vect{x}^\T A\vect{x} > 0 \ (\vect{x}\neq\vect{0})\\
  \text{Cholesky:}\quad & A = R^\T R = LDL^\T \quad (R = \sqrt{D}L^\T)\\
  \text{Weighted Least Squares:}\quad & A^\T C A\uv{\vect{x}} = A^\T C\vect{b} \qquad (C = \Sigma^{-1} = W^\T W)\\
  \text{Distributions:}\quad & \text{Binomial: } m = N/2, \ \sigma^2 = N/4; \quad \text{Poisson: } m = \lambda, \ \sigma^2 = \lambda\\
  \text{Sample Average:}\quad & \sigma^2/N \text{ for average of } N \text{ independent measurements with variance } \sigma^2
\end{align*}
\end{fs}

\formulasheet{Module 4 Formula Sheet: Eigenvalues, Diagonalization, SVD, and Complex Matrices}

\begin{fs}{1. Eigenvalues, Diagonalization, and Symmetric Matrices}
\begin{align*}
  \text{Characteristic Equation:}\quad & \det(A - \lambda I) = 0, \qquad \sum_{i=1}^n \lambda_i = \operatorname{trace}(A), \qquad \prod_{i=1}^n \lambda_i = \det A\\
  \text{Diagonalization:}\quad & AS = S\Lambda \iff A = S\Lambda S^{-1}, \qquad A^k = S\Lambda^k S^{-1}, \qquad \vect{u}_k = A^k\vect{u}_0 = \sum_{i=1}^n c_i\lambda_i^k\vect{x}_i\\
  \text{Spectral Theorem:}\quad & A = A^\T \implies A = Q\Lambda Q^\T = \sum_{i=1}^n \lambda_i\vect{q}_i\vect{q}_i^\T \quad (\text{real } \lambda_i, \text{ orthonormal } \vect{q}_i)\\
  \text{Similar Matrices:}\quad & B = M^{-1}AM \implies \text{same }\lambda_i, \quad J = \operatorname{diag}(J_1, \dots, J_s)
\end{align*}
\end{fs}

\begin{fs}{2. Singular Value Decomposition (SVD), Markov, and Complex Matrices}
\begin{align*}
  \text{SVD Factorization:}\quad & A = U\Sigma V^\T = \sum_{i=1}^r \sigma_i\vect{u}_i\vect{v}_i^\T \quad (A^\T A\vect{v}_i = \sigma_i^2\vect{v}_i, \ AA^\T\vect{u}_i = \sigma_i^2\vect{u}_i, \ \vect{u}_i = A\vect{v}_i/\sigma_i)\\
  \text{Four Subspaces:}\quad & \vect{v}_1, \dots, \vect{v}_r \in C(A^\T); \quad \vect{v}_{r+1}, \dots, \vect{v}_n \in N(A); \quad \vect{u}_1, \dots, \vect{u}_r \in C(A); \quad \vect{u}_{r+1}, \dots, \vect{u}_m \in N(A^\T)\\
  \text{Markov Matrices:}\quad & a_{ij} \ge 0, \ \sum_{i} a_{ij} = 1 \implies \lambda_1 = 1, \ |\lambda_i| \le 1, \ \vect{u}_k \to \vect{u}_\infty = \vect{x}_1\\
  \text{Complex Matrices:}\quad & \|\vect{z}\|^2 = \vect{z}^\Herm\vect{z}, \quad \vect{u}^\Herm\vect{v} = \overline{\vect{u}}^\T\vect{v}\\
  \text{Hermitian / Unitary:}\quad & A^\Herm = A \ (\text{real } \lambda_i), \qquad U^\Herm U = I \ (|\lambda|=1)
\end{align*}
\end{fs}

\end{document}
